% Lesson 22 — Vocabulary: Words and expressions related to managing personal finances — Anglais 1ère A — Cours signé OnBuch+
% Programme officiel MINESEC. Compiler avec : tectonic ENA22-vocabulary-words-and-expressions-related-to-managing-pe.tex
\newcommand{\DOCMATIERE}{Anglais}
\newcommand{\DOCNIVEAU}{1ère A}
\newcommand{\DOCMODULE}{Chapter 2 : Economic Life and Occupations}
\newcommand{\DOCLECON}{ENA22}
\newcommand{\DOCTITRE}{Lesson 22 — Vocabulary: Words and expressions related to managing personal finances}
% OnBuch+ — préambule « cours signés » ENRICHI (compilé avec tectonic / XeLaTeX)
% Système visuel « Pop » d'OnBuch+ : fond crème (jamais blanc), bordures encre
% épaisses, ombres « dures » décalées, titres Archivo Black en capitales.
% Version MATHÉMATIQUES : graphes (pgfplots/tikz), boîtes pédagogiques
% (définition, propriété, méthode, attention, le savais-tu, exercice résolu,
% exercice, TP, bilan), page de garde et signature OnBuch+ ancrée en bas.
%
% Macros attendues dans main.tex AVANT \input{preamble} :
%   \DOCMATIERE  \DOCNIVEAU  \DOCMODULE  \DOCLECON (numéro)  \DOCTITRE
\documentclass[11pt]{article}

\usepackage{fontspec}
\usepackage[english,french]{babel} % français = langue principale ; l'anglais via \eng{..} / textebox
\usepackage[a4paper,margin=2cm,top=3.4cm,bottom=2.8cm,headheight=1.4cm,footskip=1.3cm]{geometry}
\usepackage{amsmath,amssymb,amsfonts}
\usepackage{mathtools} % \xmapsto, \coloneqq... souvent utilisés par les modèles
\usepackage{cancel} % simplifications barrées (bilans de forces)
% \abs{z} (module, valeur absolue) et \norm{u} (norme), utilisés par les modèles
\providecommand{\abs}{}\let\abs\relax\DeclarePairedDelimiter{\abs}{\lvert}{\rvert}
\providecommand{\norm}{}\let\norm\relax\DeclarePairedDelimiter{\norm}{\lVert}{\rVert}
\usepackage[table]{xcolor} % [table] : fournit aussi \rowcolors (alternance de lignes)
\usepackage{pagecolor}
\usepackage{tikz}
\usetikzlibrary{arrows.meta,positioning,calc,shapes.geometric,shapes.misc,shapes.arrows,decorations.pathmorphing,decorations.pathreplacing,decorations.markings,patterns,shadows,mindmap,trees,fit,backgrounds,matrix,intersections,angles,quotes}
\usepackage{pgfplots}
\usepackage[version=4]{mhchem} % équations nucléaires \ce{^{235}_{92}U}
\usepackage[european,siunitx]{circuitikz} % schémas électriques
\pgfplotsset{compat=1.18}
\usepgfplotslibrary{fillbetween,groupplots} % aires entre courbes (calcul intégral)
\usepackage{enumitem}
\usepackage{fancyhdr}
% [most] charge toutes les bibliothèques tcolorbox (listings, minted, raster,
% documentation...) et gonfle inutilement la mémoire TeX sur un document long
% (~300 boîtes) : on ne charge que ce qui est réellement utilisé.
\usepackage[skins,breakable]{tcolorbox}
\usepackage{graphicx}
\usepackage{colortbl}
\usepackage{booktabs}
\usepackage{tabularx}
\usepackage{diagbox} % case d'en-tête diagonale des tableaux à double entrée
% Colonnes extensibles centrée / gauche / droite (C, L, R), souvent utilisées par les modèles
\newcolumntype{C}{>{\centering\arraybackslash}X}
\newcolumntype{L}{>{\raggedright\arraybackslash}X}
\newcolumntype{R}{>{\raggedleft\arraybackslash}X}
\usepackage{array}
\usepackage{multicol}
\usepackage{multirow}
\usepackage{siunitx}
% parse-numbers=false : accepte aussi \SI{2,5 \times 10^{-3}}{...}
\sisetup{locale=FR,output-decimal-marker={,},group-minimum-digits=4,parse-numbers=false}
\DeclareSIUnit\ohm{\ensuremath{\symup{\Omega}}} % U+2126 absent de la police
\DeclareSIUnit\division{div} % réglages d'oscilloscope (ms/div, V/div)
\DeclareSIUnit\tr{tr} % tours (vitesse de rotation, ex. \SI{1500}{\tr\per\minute})
\DeclareSIUnit\molar{M} % concentration molaire (ex. \SI{3}{\molar}, \SI{1}{\micro\molar})
\DeclareSIUnit\an{an} % année (le modèle confond parfois avec \year, un compteur TeX) — ex. \SI{5}{\kilo\meter\per\an}
\DeclareSIUnit\FCFA{FCFA} % franc CFA (ex. \SI{500}{\FCFA\per\kilo\gram})
\DeclareSIUnit\degreeBrix{\textdegree Brix} % teneur en sucre d'un sirop/jus (réfractométrie)
\DeclareSIUnit\month{mois} % le modèle utilise parfois \month, un compteur TeX, comme unité de durée
\usepackage{qrcode}
% \degree : commande fréquemment utilisée par les modèles pour un angle ou une
% température en dehors de \SI{}{} (non fournie par les paquets chargés).
\providecommand{\degree}{^{\circ}}
% \qed (fin de démonstration), utilisé par les modèles sans amsthm
\providecommand{\qed}{\hfill$\square$}
% \barre : double barre des valeurs interdites dans un tableau de variations (tkz-tab)
\providecommand{\barre}{\ensuremath{\|}}
% environnement proof (amsthm non chargé) utilisé par les modèles
\ifdefined\proof\else\newenvironment{proof}[1][Démonstration]{\par\noindent\textit{#1.}\ }{\qed\par}\fi
\usepackage{lastpage}
\usepackage{hyperref}
\hypersetup{hidelinks}

% ============================================================
% Palette « Pop » OnBuch+ — mêmes hex que la classe P (plus_ui.dart)
% ============================================================
\definecolor{poporange}{HTML}{F59321}
\definecolor{popdark}{HTML}{E07A0C}
\definecolor{popcream}{HTML}{FFF6E8}
\definecolor{popink}{HTML}{1C1714}
\definecolor{poppurple}{HTML}{7A5AE0}
\definecolor{popgreen}{HTML}{2E9E6A}
\definecolor{poppink}{HTML}{E0457B}
\definecolor{popblue}{HTML}{2F7DE1}
\definecolor{popgold}{HTML}{C98A12}
\definecolor{popmuted}{HTML}{8A7F76}
\definecolor{popline}{HTML}{EADFD2}
\definecolor{popgray}{HTML}{8A7F76}
\colorlet{popgrey}{popgray}
% Alias tolérants (le modèle nomme parfois les couleurs différemment)
\colorlet{popwhite}{white}
\colorlet{popblack}{popink}
\colorlet{popred}{poppink}
\colorlet{popyellow}{popgold}
% Teintes claires (fonds de tableaux/encadrés) — le modèle nomme parfois
% les couleurs avec un suffixe L (ex. popblueL) sans les avoir définies.
\colorlet{poporangeL}{poporange!20}
\colorlet{popdarkL}{popdark!20}
\colorlet{poppurpleL}{poppurple!20}
\colorlet{popgreenL}{popgreen!20}
\colorlet{poppinkL}{poppink!20}
\colorlet{popblueL}{popblue!20}
\colorlet{popgoldL}{popgold!20}
\colorlet{popredL}{popred!20}
\colorlet{popgrayL}{popgray!20}
% Teintes claires pour fonds de tableaux / zones
\colorlet{popblueL}{popblue!10!white}
\colorlet{poporangeL}{poporange!14!white}
\colorlet{popgreenL}{popgreen!12!white}
\colorlet{poppurpleL}{poppurple!12!white}
\colorlet{poppinkL}{poppink!12!white}
\colorlet{popgoldL}{popgold!14!white}

\pagecolor{popcream}

% ============================================================
% Typographie
% ============================================================
\defaultfontfeatures{Ligatures=TeX}
\setmainfont{PlusJakartaSans-Medium.ttf}[Path=../../../fonts/,BoldFont=PlusJakartaSans-Bold.ttf]
% Maths en sans-serif (Fira Math) avec chiffres et lettres droites de Plus
% Jakarta, pour que les formules se fondent dans le texte.
\usepackage[math-style=ISO,bold-style=ISO]{unicode-math}
\setmathfont{FiraMath-Regular.otf}[Path=../../../fonts/]
\setmathfont{PlusJakartaSans-Medium.ttf}[Path=../../../fonts/,range={up/num,up/{Latin,latin},\mathup}]
\usepackage{icomma} % virgule décimale sans espace en mode math
% Caractères mathématiques Unicode absents des polices de texte (Plus Jakarta,
% Archivo Black) : rendus par leur équivalent mathématique, en texte comme en
% formule — sinon ils s'affichent en carré vide (ex. « Arithmétique dans ℤ »).
\usepackage{newunicodechar}
\newunicodechar{ℝ}{\ensuremath{\mathbb{R}}}
\newunicodechar{ℤ}{\ensuremath{\mathbb{Z}}}
\newunicodechar{ℂ}{\ensuremath{\mathbb{C}}}
\newunicodechar{ℕ}{\ensuremath{\mathbb{N}}}
\newunicodechar{ℚ}{\ensuremath{\mathbb{Q}}}
\newunicodechar{𝔻}{\ensuremath{\mathbb{D}}}
\newunicodechar{α}{\ensuremath{\alpha}}
\newunicodechar{β}{\ensuremath{\beta}}
\newunicodechar{γ}{\ensuremath{\gamma}}
\newunicodechar{δ}{\ensuremath{\delta}}
\newunicodechar{ε}{\ensuremath{\varepsilon}}
\newunicodechar{θ}{\ensuremath{\theta}}
\newunicodechar{λ}{\ensuremath{\lambda}}
\newunicodechar{μ}{\ensuremath{\mu}}
\newunicodechar{σ}{\ensuremath{\sigma}}
\newunicodechar{τ}{\ensuremath{\tau}}
\newunicodechar{φ}{\ensuremath{\varphi}}
\newunicodechar{ψ}{\ensuremath{\psi}}
\newunicodechar{ω}{\ensuremath{\omega}}
\newunicodechar{Σ}{\ensuremath{\Sigma}}
% majuscules produites par \MakeUppercase dans les titres (α-aminés -> Α)
\newunicodechar{Α}{\ensuremath{\alpha}}
\newunicodechar{Β}{\ensuremath{\beta}}
\newunicodechar{Γ}{\ensuremath{\Gamma}}
\newunicodechar{Δ}{\ensuremath{\Delta}}
\newunicodechar{Ε}{\ensuremath{\varepsilon}}
\newunicodechar{Θ}{\ensuremath{\Theta}}
\newunicodechar{Λ}{\ensuremath{\Lambda}}
\newunicodechar{Μ}{\ensuremath{\mu}}
\newunicodechar{Τ}{\ensuremath{\tau}}
\newunicodechar{Φ}{\ensuremath{\Phi}}
\newunicodechar{Ψ}{\ensuremath{\Psi}}
\newunicodechar{Ω}{\ensuremath{\Omega}}
\newunicodechar{↦}{\ensuremath{\mapsto}}
\newunicodechar{∈}{\ensuremath{\in}}
\newunicodechar{∉}{\ensuremath{\notin}}
\newunicodechar{∧}{\ensuremath{\wedge}}
\newunicodechar{⇔}{\ensuremath{\Leftrightarrow}}
\newunicodechar{⟺}{\ensuremath{\Longleftrightarrow}}
\newunicodechar{⇒}{\ensuremath{\Rightarrow}}
\newunicodechar{⟹}{\ensuremath{\Longrightarrow}}
\newunicodechar{∘}{\ensuremath{\circ}}
\newunicodechar{∪}{\ensuremath{\cup}}
\newunicodechar{∩}{\ensuremath{\cap}}
\newunicodechar{≡}{\ensuremath{\equiv}}
\newunicodechar{⊂}{\ensuremath{\subset}}
\newunicodechar{⊥}{\ensuremath{\perp}}
\newunicodechar{∃}{\ensuremath{\exists}}
\newunicodechar{∀}{\ensuremath{\forall}}
\newunicodechar{∛}{\ensuremath{\sqrt[3]{\,}}}
\newunicodechar{∜}{\ensuremath{\sqrt[4]{\,}}}
\newunicodechar{⃗}{\ensuremath{{}^{\to}}}
\newunicodechar{ⁿ}{\ifmmode^{n}\else\textsuperscript{\ensuremath{n}}\fi}
\newunicodechar{ˣ}{\ifmmode^{x}\else\textsuperscript{\ensuremath{x}}\fi}
\newunicodechar{ᵇ}{\ifmmode^{b}\else\textsuperscript{\ensuremath{b}}\fi}
\newunicodechar{ʳ}{\ifmmode^{r}\else\textsuperscript{\ensuremath{r}}\fi}
\newunicodechar{ᵘ}{\ifmmode^{u}\else\textsuperscript{\ensuremath{u}}\fi}
\newunicodechar{ʸ}{\ifmmode^{y}\else\textsuperscript{\ensuremath{y}}\fi}
\newunicodechar{ᵃ}{\ifmmode^{a}\else\textsuperscript{\ensuremath{a}}\fi}
\newunicodechar{ᵉ}{\ifmmode^{e}\else\textsuperscript{\ensuremath{e}}\fi}
\newunicodechar{ᵏ}{\ifmmode^{k}\else\textsuperscript{\ensuremath{k}}\fi}
\newunicodechar{ᵖ}{\ifmmode^{p}\else\textsuperscript{\ensuremath{p}}\fi}
\newunicodechar{ᴺ}{\ifmmode^{N}\else\textsuperscript{\ensuremath{N}}\fi}
\newunicodechar{ᶜ}{\ifmmode^{c}\else\textsuperscript{\ensuremath{c}}\fi}
\newunicodechar{ʰ}{\ifmmode^{h}\else\textsuperscript{\ensuremath{h}}\fi}
\newunicodechar{ᵠ}{\ifmmode^{q}\else\textsuperscript{\ensuremath{q}}\fi}
\newunicodechar{ₙ}{\ifmmode_{n}\else\textsubscript{\ensuremath{n}}\fi}
\newunicodechar{ₐ}{\ifmmode_{a}\else\textsubscript{\ensuremath{a}}\fi}
\newunicodechar{ᵢ}{\ifmmode_{i}\else\textsubscript{\ensuremath{i}}\fi}
\newunicodechar{ᵣ}{\ifmmode_{r}\else\textsubscript{\ensuremath{r}}\fi}
\newunicodechar{ₖ}{\ifmmode_{k}\else\textsubscript{\ensuremath{k}}\fi}
\newunicodechar{ᵦ}{\ifmmode_{\beta}\else\textsubscript{\ensuremath{\beta}}\fi}
\newunicodechar{ₓ}{\ifmmode_{x}\else\textsubscript{\ensuremath{x}}\fi}
\newfontfamily\popdisplay{ArchivoBlack-Regular.ttf}[Path=../../../fonts/]
\newfontfamily\poplogo{SpaceGrotesk-Bold.ttf}[Path=../../../fonts/]
\newfontfamily\popsemi{PlusJakartaSans-SemiBold.ttf}[Path=../../../fonts/]
\newfontfamily\popxbold{PlusJakartaSans-ExtraBold.ttf}[Path=../../../fonts/]
\newfontfamily\popxboldls{PlusJakartaSans-ExtraBold.ttf}[Path=../../../fonts/,LetterSpace=3.0]

\setlist[enumerate]{leftmargin=1.7em,itemsep=5pt,topsep=4pt}
\setlist[itemize]{leftmargin=1.7em,itemsep=4pt,topsep=4pt,label={\color{poporange}\textbullet}}
\setlength{\parindent}{0pt}
\setlength{\parskip}{5pt}
\renewcommand{\arraystretch}{1.25}

% pgfplots : style Pop par défaut
\pgfplotsset{
  every axis/.append style={axis line style={popink,line width=1pt},tick style={popink},
    grid style={popline,line width=0.5pt},label style={font=\small},tick label style={font=\footnotesize}},
  popcurve/.style={poporange,line width=1.8pt,no markers,smooth},
}
% Même style disponible dans un \draw TikZ simple (hors pgfplots)
\tikzset{popcurve/.style={poporange,line width=1.8pt}}
\tikzset{popgrid/.style={popline,very thin}}
\colorlet{popcurve}{poporange} % le nom du style est parfois utilisé comme couleur

% ============================================================
% Pill / squiggle
% ============================================================
\newcommand{\poppill}[3]{%
  \tikz[baseline=-0.55ex]\node[draw=popink,line width=1.2pt,rounded corners=2.6pt,%
    fill=#2,inner xsep=6.5pt,inner ysep=3pt]{%
    \popxboldls\fontsize{7}{7}\selectfont\color{#3}\MakeUppercase{#1}};%
}
\newcommand{\popsquiggle}[1]{%
  \tikz[baseline=-0.4ex]{
    \draw[#1,line width=2.6pt,line cap=round]
      plot [smooth] coordinates {(0,0.09) (0.34,0.01) (0.68,0.21) (1.02,0.01) (1.36,0.10)};
  }%
}
% Mot-clé surligné (marqueur orange sous le texte)
\newcommand{\cle}[1]{{\popxbold\color{popdark}#1}}

% ============================================================
% En-tête / pied
% ============================================================
\pagestyle{fancy}
\fancyhf{}
\renewcommand{\headrulewidth}{0pt}
\renewcommand{\footrulewidth}{0pt}
\lhead{\raisebox{-0.15ex}{\poplogo\fontsize{13}{13}\selectfont\color{popink}OnBuch\color{poporange}+}}
\rhead{\poppill{\DOCMATIERE\ \textbullet\ \DOCNIVEAU\ \textbullet\ Leçon \DOCLECON}{white}{popink}}
\lfoot{\popxbold\fontsize{7.6}{7.6}\selectfont\color{popmuted}\MakeUppercase{Cours signé OnBuch+}\ \textbullet\ {\popsemi\DOCTITRE}}
\rfoot{\tikz[baseline=-0.6ex]\node[rounded corners=8.5pt,fill=popink,minimum height=17pt,minimum width=17pt,inner xsep=5pt,inner sep=0]{\popxbold\fontsize{8}{8}\selectfont\color{popcream}\thepage\,/\,\pageref*{LastPage}};}

% ============================================================
% Titres
% ============================================================
\newcommand{\chaptitre}[2]{%
  \begin{tcolorbox}[enhanced,colback=poporange,colframe=popink,boxrule=2.6pt,arc=11pt,
      drop shadow=popink,left=15pt,right=15pt,top=12pt,bottom=11pt,before skip=3pt,after skip=18pt]
    \poppill{#2}{popink}{popcream}\par\vspace{9pt}
    {\raggedright\hyphenpenalty=10000\exhyphenpenalty=10000\popdisplay\fontsize{22}{24}\selectfont\color{popcream}\MakeUppercase{#1}\par}
  \end{tcolorbox}
}
% Section numérotée : pastille orange avec le numéro + titre Archivo
\newcounter{popsec}
\newcommand{\coursec}[1]{%
  \stepcounter{popsec}\setcounter{popsub}{0}%
  \par\vspace{16pt}\noindent
  \begin{minipage}{\linewidth}
  \tikz[baseline=-0.7ex]\node[draw=popink,line width=1.6pt,fill=poporange,rounded corners=4pt,minimum size=22pt,inner sep=0]{\popdisplay\fontsize{12}{12}\selectfont\color{popcream}\thepopsec};%
  \hspace{8pt}{\popdisplay\fontsize{15}{17}\selectfont\color{popink}\MakeUppercase{#1}}\par
  \vspace{4pt}\noindent{\color{poporange}\rule{36pt}{3.6pt}}
  \end{minipage}\par\nopagebreak\vspace{8pt}
}
\newcounter{popsub}
\newcommand{\courssub}[1]{%
  \stepcounter{popsub}%
  \par\vspace{9pt}\noindent{\popxbold\fontsize{11.5}{13}\selectfont\color{popink}{\color{poporange}\thepopsec.\thepopsub}\hspace{6pt}#1}\par\nopagebreak\vspace{4pt}
}

% ============================================================
% Boîtes pédagogiques — même recette : fond blanc, bordure épaisse colorée,
% ombre dure encre, badge plein. Titre optionnel [#1].
% ============================================================
\tcbset{popbase/.style={breakable,enhanced,colback=white,boxrule=2.3pt,arc=9pt,
  shadow={3pt}{-3pt}{0pt}{popink},left=14pt,right=14pt,top=11pt,bottom=11pt,before skip=11pt,after skip=15pt,
  fontupper=\popsemi\fontsize{10.6}{14.5}\selectfont\color{popink},
  fontlower=\popsemi\fontsize{10.6}{14.5}\selectfont\color{popink}}}
% Coche dessinée (le glyphe ✓ n'existe pas dans les polices du document)
\newcommand{\popcheck}[1]{\tikz[baseline=-0.5ex]\draw[#1,line width=1.6pt,line cap=round,line join=round](-0.13,0.0)--(-0.04,-0.09)--(0.13,0.1);}
\providecommand{\coche}{\popcheck{popgreen}}
\providecommand{\resultat}[1]{\par\smallskip\noindent\fcolorbox{popgreen}{popgreenL}{\parbox{\dimexpr\linewidth-2\fboxsep-2\fboxrule\relax}{\textbf{Résultat.} #1}}\par\smallskip}
\newcommand{\popboxhead}[3]{\poppill{#1}{#2}{white}\ifx\relax#3\relax\else\hspace{6pt}{\popxbold\fontsize{10.6}{12}\selectfont\color{popink}#3}\fi\par\vspace{7pt}}

\newtcolorbox{aretenir}[1][]{popbase,colframe=popgreen,before upper={\popboxhead{À retenir}{popgreen}{#1}}}
% Texte en anglais (document de lecture, dialogue, extrait) : cadre
% d'étude. Un titre optionnel (source, auteur) s'affiche dans l'en-tête.
\newtcolorbox{textebox}[1][]{popbase,colframe=popblue,before upper={\popboxhead{Text}{popblue}{#1}\begin{otherlanguage}{english}},after upper={\end{otherlanguage}}}
% Marques ✓ / ✗ (forme correcte / fautive) souvent utilisées par les modèles
\providecommand{\cross}{\textcolor{poppink}{\ensuremath{\times}}}
\providecommand{\xmark}{\cross}\providecommand{\crossmark}{\cross}\providecommand{\cmark}{\ensuremath{\checkmark}}
% Ligne de réponse / justification à compléter (inventée par les modèles)
\providecommand{\justification}[1]{\par\noindent\textit{Justification~:} \dotfill\par}
% \textipa (package tipa non chargé) : la police ne couvre pas l'API -> simple passage
\providecommand{\textipa}[1]{#1}
% Soulignement (ulem non chargé) : \uline, \ul
\providecommand{\uline}[1]{\underline{#1}}\providecommand{\ul}[1]{\underline{#1}}
% \glossaire{\item ...} inventée par les modèles : liste de glossaire
\providecommand{\glossaire}[1]{\begin{itemize}#1\end{itemize}}
% \alert (beamer) inventée par les modèles : texte mis en évidence
\providecommand{\alert}[1]{\textbf{\textcolor{poppink}{#1}}}
% variantes ASCII d'ordinaux écrites en macro
\providecommand{\iere}{\textsuperscript{re}}\providecommand{\ieme}{\textsuperscript{e}}\providecommand{\ere}{\textsuperscript{re}}\providecommand{\eme}{\textsuperscript{e}}\providecommand{\er}{\textsuperscript{er}}\providecommand{\ie}{\textsuperscript{e}}
% Ordinaux français écrits en macro par les modèles : 1\ière, 3\ième
\newcommand{\ière}{\textsuperscript{re}}\newcommand{\ième}{\textsuperscript{e}}
% \exemple (inventée par les modèles) : étiquette « Exemple : »
\providecommand{\exemple}{\textbf{Exemple~:}\ }
% \square utilisable aussi hors mode math (cases à cocher)
\AtBeginDocument{\let\squareMath\square\renewcommand{\square}{\ensuremath{\squareMath}}}
% Mot ou phrase en anglais à l'intérieur d'un paragraphe français
\newcommand{\eng}[1]{\ifmmode\text{\textit{#1}}\else\begingroup\selectlanguage{english}\itshape #1\endgroup\fi}
\let\esp\eng % alias toléré
% flèches d'intonation inventées par les modèles
\providecommand{\textsearrow}{}\renewcommand{\textsearrow}{\ensuremath{\searrow}}\providecommand{\textnearrow}{}\renewcommand{\textnearrow}{\ensuremath{\nearrow}}
% \fr{..} : mot français (inventé par les modèles)
\providecommand{\fr}[1]{\textit{#1}}
\newtcolorbox{exemplebox}[1][]{popbase,colframe=poporange,before upper={\popboxhead{Exemple}{poporange}{#1}}}
\newtcolorbox{definition}[1][]{popbase,colframe=popblue,before upper={\popboxhead{Définition}{popblue}{#1}}}
\newtcolorbox{propriete}[1][]{popbase,colframe=poppurple,before upper={\popboxhead{Propriété}{poppurple}{#1}}}
\newtcolorbox{methode}[1][]{popbase,colframe=popgold,before upper={\popboxhead{Méthode}{popgold}{#1}}}
\newtcolorbox{attention}[1][]{popbase,colframe=poppink,before upper={\popboxhead{Attention}{poppink}{#1}}}
\newtcolorbox{experience}[1][]{popbase,colframe=popblue,colback=popblueL,before upper={\popboxhead{Expérience}{popblue}{#1}}}
% « Le savais-tu ? » : carte violette pleine (contexte, culture, Cameroun)
\newtcolorbox{savaistu}[1][]{popbase,colback=poppurple,colframe=popink,
  fontupper=\popsemi\fontsize{10.4}{14.2}\selectfont\color{white},
  before upper={\poppill{Le savais-tu\,?}{white}{poppurple}\ifx\relax#1\relax\else\hspace{6pt}{\popxbold\color{white}#1}\fi\par\vspace{7pt}}}
% Exercice résolu : énoncé en haut, \tcblower puis la solution sur fond crème
\newcounter{popexo}\newcounter{popexr}
\newtcolorbox{exoresolu}[1][]{popbase,colframe=poporange,colbacklower=popcream,
  segmentation style={popink,line width=1.2pt,dashed},
  before upper={\stepcounter{popexr}\popboxhead{Exercice résolu \thepopexr}{poporange}{#1}},
  before lower={\poppill{Solution}{popink}{popcream}\par\vspace{6pt}}}
% Exercice (non résolu en place) — la correction est dans la partie Corrigés
\newtcolorbox{exercice}[1][]{popbase,colframe=popink,
  before upper={\stepcounter{popexo}\popboxhead{Exercice \thepopexo}{popink}{#1}}}
% Corrigé
\newtcolorbox{corrige}[1][]{popbase,colframe=popgreen,colback=popgreenL,
  before upper={\popboxhead{Corrigé}{popgreen}{#1}}}
% Objectifs / prérequis / bilan
\newtcolorbox{objectifs}{popbase,colframe=popink,colback=poporangeL,
  before upper={\popboxhead{Objectifs de la leçon}{popink}{}}}
\newtcolorbox{prerequis}{popbase,colframe=popink,colback=popblueL,
  before upper={\popboxhead{Prérequis}{popblue}{}}}
\newtcolorbox{bilan}{popbase,colframe=popink,colback=white,boxrule=2.8pt,
  before upper={\popboxhead{Fiche bilan}{poporange}{L'essentiel à connaître pour l'examen}}}
% Figure encadrée avec légende
\newtcolorbox{popfigure}[1][]{enhanced,breakable=false,colback=white,colframe=popline,boxrule=1.4pt,arc=8pt,
  left=10pt,right=10pt,top=10pt,bottom=8pt,before skip=10pt,after skip=12pt,halign=center,
  lower separated=false,halign lower=center,
  fontlower=\popsemi\fontsize{9}{11.5}\selectfont\color{popmuted},
  #1}
\newcommand{\legende}[1]{\par\vspace{6pt}{\popsemi\fontsize{9}{11.5}\selectfont\color{popmuted}#1}}

% ============================================================
% Signature OnBuch+
% ============================================================
\newcommand{\poplogoinline}[1]{{\poplogo\fontsize{#1}{#1}\selectfont\color{popink}OnBuch\color{poporange}+}}
\newcommand{\signatureonbuch}{%
  \begin{tcolorbox}[enhanced,colback=popink,colframe=popink,boxrule=2.2pt,arc=10pt,
      drop shadow=poporange,left=14pt,right=14pt,top=10pt,bottom=10pt,before skip=0pt,after skip=0pt]
    \tikz[baseline=-0.7ex]\node[circle,fill=popgreen,draw=popcream,line width=1.3pt,minimum size=22pt,inner sep=0]{\popcheck{white}};%
    \hspace{9pt}\begin{minipage}[c]{0.62\linewidth}
      {\popxbold\fontsize{12}{13}\selectfont\color{popcream}Signé\ \textbullet\ {\poplogo OnBuch\color{poporange}+}}\par\vspace{2pt}
      {\raggedright\popsemi\fontsize{8}{10.5}\selectfont\color{popline}\MakeUppercase{Équipe pédagogique OnBuch+}\\\MakeUppercase{Conforme au programme officiel MINESEC}\par}
    \end{minipage}\hfill
    {\poplogo\fontsize{22}{22}\selectfont\color{popcream}OnBuch\color{poporange}+}
  \end{tcolorbox}%
}

% ============================================================
% Page de garde
% #1 titre · #2 matière · #3 niveau · #4 module · #5 n° leçon · #6 durée ·
% #7 description / objectifs (texte ou itemize)
% ============================================================
\newcommand{\pagegarde}[7]{%
  \thispagestyle{empty}
  \newgeometry{a4paper,margin=1.7cm,top=1.6cm,bottom=1.4cm}%
  \begin{tikzpicture}[remember picture,overlay]
    \fill[poporange!18] ([xshift=-3.2cm,yshift=-3.6cm]current page.north east) circle (4.6cm);
    \fill[poppurple!14] ([xshift=2.4cm,yshift=7.6cm]current page.south west) circle (3.8cm);
  \end{tikzpicture}%
  {\poplogo\fontsize{30}{30}\selectfont\color{popink}OnBuch\color{poporange}+}\hfill
  \raisebox{0.6ex}{\poppill{Cours signé \textbullet\ Premium}{popink}{popcream}}\par
  \vspace{1pt}\popsquiggle{poppurple}\par
  \vspace{8pt}
  \begin{tcolorbox}[enhanced,colback=poporange,colframe=popink,boxrule=2.8pt,arc=13pt,
      drop shadow=popink,left=18pt,right=18pt,top=12pt,bottom=13pt,after skip=10pt]
    \poppill{#2 \textbullet\ #3}{popink}{popcream}\hspace{5pt}\poppill{#4}{white}{popink}\par\vspace{8pt}
    {\popdisplay\fontsize{14}{14}\selectfont\color{popink}LEÇON #5}\par\vspace{4pt}
    {\raggedright\hyphenpenalty=10000\exhyphenpenalty=10000\popdisplay\fontsize{25}{27}\selectfont\color{popcream}\MakeUppercase{#1}\par}
    \vspace{8pt}
    {\popxbold\fontsize{9.5}{9.5}\selectfont\color{popink}\MakeUppercase{Durée conseillée : #6}}
  \end{tcolorbox}
  \begin{tcolorbox}[enhanced,colback=white,colframe=popink,boxrule=2.2pt,arc=10pt,
      drop shadow=popink,left=16pt,right=16pt,top=10pt,bottom=10pt,after skip=8pt]
    \poppill{Dans cette leçon}{poppurple}{white}\par\vspace{5pt}
    \popsemi\fontsize{9.3}{12.6}\selectfont\color{popink}#7
  \end{tcolorbox}
  \noindent\begin{minipage}[t]{0.487\linewidth}
    \begin{tcolorbox}[enhanced,colback=popgreen,colframe=popink,boxrule=2.2pt,arc=10pt,
        drop shadow=popink,left=11pt,right=11pt,top=10pt,bottom=10pt,height=3.1cm,valign=center]
      \tcbox[colback=white,colframe=popink,boxrule=1.3pt,arc=5pt,boxsep=0pt,nobeforeafter,
        left=3pt,right=3pt,top=3pt,bottom=3pt,valign=center]{\qrcode[height=1.9cm]{https://play.google.com/store/apps/details?id=com.luvvix.onbuch&hl=fr}}%
      \hspace{8pt}\begin{minipage}[c]{0.5\linewidth}
        {\popdisplay\fontsize{11}{12.5}\selectfont\color{white}\MakeUppercase{Télécharge OnBuch}}\par\vspace{4pt}
        {\raggedright\popsemi\fontsize{8.4}{11}\selectfont\color{white}Gratuit \textbullet\ Google Play\\Tuteur IA, annales, concours, résultats\par}
      \end{minipage}
    \end{tcolorbox}
  \end{minipage}\hfill
  \begin{minipage}[t]{0.487\linewidth}
    \begin{tcolorbox}[enhanced,colback=poppurple,colframe=popink,boxrule=2.2pt,arc=10pt,
        drop shadow=popink,left=11pt,right=11pt,top=10pt,bottom=10pt,height=3.1cm,valign=center]
      \tikz[baseline=-0.7ex]\node[circle,fill=white,draw=popink,line width=1.3pt,minimum size=1.6cm,inner sep=0]{\popdisplay\fontsize{24}{24}\selectfont\color{poppurple}+};%
      \hspace{8pt}\begin{minipage}[c]{0.6\linewidth}
        {\popdisplay\fontsize{11}{12.5}\selectfont\color{white}\MakeUppercase{Passe à OnBuch+}}\par\vspace{4pt}
        {\raggedright\popsemi\fontsize{8.4}{11}\selectfont\color{white}Tous les cours signés, Tuteur IA illimité, fiches PDF — onglet OnBuch+ de l'app\par}
      \end{minipage}
    \end{tcolorbox}
  \end{minipage}
  \vspace{8pt}
  \popcontenu
  \vfill
  \signatureonbuch
  \restoregeometry
  \newpage
}

% Rangée « Ce cours contient » (page de garde)
\newcommand{\poptile}[3]{% #1 couleur #2 gros texte #3 légende
  \begin{tcolorbox}[enhanced,colback=#1,colframe=popink,boxrule=2pt,arc=9pt,shadow={2.5pt}{-2.5pt}{0pt}{popink},
      left=6pt,right=6pt,top=9pt,bottom=9pt,halign=center,valign=center,height=1.75cm,width=\linewidth,nobeforeafter]
    {\popdisplay\fontsize{12.5}{13}\selectfont\color{white}\MakeUppercase{#2}}\par\vspace{3pt}
    {\centering\popsemi\fontsize{7.8}{9.5}\selectfont\color{white}#3\par}
  \end{tcolorbox}}
\newcommand{\popcontenu}{%
  \par\noindent{\popxboldls\fontsize{7.5}{7.5}\selectfont\color{popmuted}CE COURS SIGNÉ CONTIENT}\par\vspace{7pt}
  \noindent\begin{minipage}[t]{0.235\linewidth}\poptile{popblue}{Cours}{complet, illustré, pas à pas}\end{minipage}\hfill
  \begin{minipage}[t]{0.235\linewidth}\poptile{poporange}{Méthodes}{et exercices résolus}\end{minipage}\hfill
  \begin{minipage}[t]{0.235\linewidth}\poptile{poppink}{Exercices}{type examen + corrigés}\end{minipage}\hfill
  \begin{minipage}[t]{0.235\linewidth}\poptile{popgreen}{Fiche bilan}{l'essentiel à retenir}\end{minipage}\par}

% Sommaire
\newtcolorbox{sommaire}{enhanced,colback=white,colframe=popink,boxrule=2.2pt,arc=10pt,
  drop shadow=popink,left=18pt,right=18pt,top=13pt,bottom=13pt,before skip=6pt,after skip=20pt,
  before upper={\poppill{Sommaire}{poppurple}{white}\par\vspace{9pt}\popsemi\fontsize{10.6}{14.5}\selectfont\color{popink}}}

% Signature de fin de cours — poussée tout en bas de la dernière page
\newcommand{\findecours}{%
  \par\vfill\nopagebreak
  \begin{center}\popsquiggle{poporange}\end{center}\vspace{4pt}
  \signatureonbuch
}
\AtBeginDocument{\def\llbracket{\mathopen{[\mkern-3.5mu[}}\def\rrbracket{\mathclose{]\mkern-3.5mu]}}} % intervalles d'entiers (glyphe ⟦ absent de la police)
% Glyphes absents de Fira Math : redéfinis à partir de glyphes disponibles
\AtBeginDocument{%
  \def\setminus{\mathbin{\backslash}}\let\smallsetminus\setminus
  \def\vdots{\mathord{\vbox{\baselineskip3.2pt\lineskiplimit0pt\kern4pt\hbox{.}\hbox{.}\hbox{.}}}}%
  \def\ddots{\mathinner{\mkern1mu\raise6.5pt\hbox{.}\mkern2mu\raise3.6pt\hbox{.}\mkern2mu\raise0.7pt\hbox{.}\mkern1mu}}%
  \def\triangle{\mathord{\scalebox{0.9}{$\symup{\Delta}$}}}%
  \def\nabla{\mathord{\raisebox{\depth}{\scalebox{1}[-1]{$\symup{\Delta}$}}}}%
  \def\checkmark{\popcheck{popdark}}%
  \def\star{\mathbin{\ast}}\def\bigstar{\mathord{\ast}}%
  \def\diamond{\mathbin{\scalebox{0.6}{$\Diamond$}}}%
  \def\bigcup{\mathop{\mathchoice{\raisebox{-0.3em}{\scalebox{1.7}{$\cup$}}}{\scalebox{1.25}{$\cup$}}{\cup}{\cup}}\displaylimits}%
  \def\bigcap{\mathop{\mathchoice{\raisebox{-0.3em}{\scalebox{1.7}{$\cap$}}}{\scalebox{1.25}{$\cap$}}{\cap}{\cap}}\displaylimits}%
  \def\lesseqqgtr{\mathrel{\lessgtr}}\def\bigoplus{\mathop{\oplus}\displaylimits}%
}
\providecommand{\ssi}{si et seulement si} % abréviation fréquente
\AtBeginDocument{\providecommand{\wideparen}{\overparen}} % arc de cercle

\begin{document}

\pagegarde{Lesson 22 — Vocabulary: Words and expressions related to managing personal finances}{Anglais}{1ère A}{Chapter 2 : Economic Life and Occupations}{ENA22}{2 h}{Cette leçon de vocabulaire présente le lexique essentiel de la gestion des finances personnelles : gagner, dépenser, épargner, emprunter et budgétiser. À travers des champs lexicaux, des collocations, des expressions idiomatiques et des mises en garde contre les faux amis, l'élève apprend à parler avec précision de l'argent au quotidien.
\vspace{4pt}
\begin{multicols}{2}\small
\begin{itemize}
\item À la fin de cette leçon, je sais identifier et définir le vocabulaire de base des finances personnelles (income, expenses, budget, savings, debt...)
\item À la fin de cette leçon, je sais employer les verbes clés liés à l'argent (earn, spend, save, borrow, lend, owe, afford, waste)
\item À la fin de cette leçon, je sais utiliser les collocations courantes du domaine financier (make money, save up for, pay off a debt, live within one's means)
\item À la fin de cette leçon, je sais comprendre et employer des expressions idiomatiques sur l'argent (to be broke, to cost a fortune, to make ends meet)
\item À la fin de cette leçon, je sais éviter les faux amis et les calques du français (economy/économie, to economize, actually, to assist)
\item À la fin de cette leçon, je sais former des mots de la même famille (save/savings/saver, spend/spending, lend/lender)
\end{itemize}\end{multicols}}

\begin{sommaire}
\begin{multicols}{2}
\begin{enumerate}
\item Core vocabulary: money in, money out
\item Key verbs: earn, spend, save, borrow, lend
\item Collocations and idiomatic expressions
\item False friends and traps for francophones
\item Word families, pronunciation and spelling
\item Using the vocabulary in context
\item Activité d'intégration
\item Méthodes et exercices résolus
\item Exercices
\item Corrigés des exercices
\item Fiche bilan
\end{enumerate}
\end{multicols}
\end{sommaire}

\begin{objectifs}
\begin{itemize}
\item À la fin de cette leçon, je sais identifier et définir le vocabulaire de base des finances personnelles (income, expenses, budget, savings, debt...)
\item À la fin de cette leçon, je sais employer les verbes clés liés à l'argent (earn, spend, save, borrow, lend, owe, afford, waste)
\item À la fin de cette leçon, je sais utiliser les collocations courantes du domaine financier (make money, save up for, pay off a debt, live within one's means)
\item À la fin de cette leçon, je sais comprendre et employer des expressions idiomatiques sur l'argent (to be broke, to cost a fortune, to make ends meet)
\item À la fin de cette leçon, je sais éviter les faux amis et les calques du français (economy/économie, to economize, actually, to assist)
\item À la fin de cette leçon, je sais former des mots de la même famille (save/savings/saver, spend/spending, lend/lender)
\item À la fin de cette leçon, je sais prononcer et orthographier correctement les mots financiers difficiles (debt, receipt, mortgage, budget)
\item À la fin de cette leçon, je sais réemployer ce lexique dans des phrases, des dialogues et un court texte sur la gestion de mon argent
\end{itemize}
\end{objectifs}

\begin{prerequis}
\begin{itemize}
\item Vocabulaire de base de l'argent vu au collège et en Seconde (money, price, buy, sell, pay)
\item Les noms dénombrables et indénombrables (money est indénombrable)
\item Le présent simple et le present continuous pour décrire des habitudes financières
\item Les verbes suivis d'un complément d'objet direct ou indirect (lend somebody something)
\item Les chiffres et les montants en anglais (pounds, dollars, francs CFA)
\end{itemize}
\end{prerequis}

\newpage

\coursec{Core vocabulary: money in, money out}

\courssub{Warm-up: Aminatou's problem}

Imaginez la scène suivante :
\begin{textebox}[Aminatou's dilemma]
Aminatou, a student in Buea, receives 10,000 FCFA as pocket money at the beginning of the month, but by the second week she has nothing left. Her friend in London faces the same problem with her £20 weekly allowance.
\end{textebox}

Comment Aminatou et son amie peuvent-elles \eng{budget} (établir un budget), \eng{save} (épargner) et \eng{spend wisely} (dépenser judicieusement) ? Pour en parler en anglais, il faut d'abord maîtriser le vocabulaire de la gestion de l'argent. C'est l'objectif de cette leçon : apprendre les mots et expressions anglais qui permettent de parler de l'argent qui entre (\eng{money in}) et de l'argent qui sort (\eng{money out}).

\begin{methode}[Problématique de la leçon]
\begin{enumerate}
\item Observer un texte authentique où apparaît le vocabulaire des finances personnelles.
\item Classer les mots en trois champs lexicaux : l'argent qui entre, l'argent qui sort, la gestion de l'argent.
\item Distinguer les mots proches (\eng{salary/wages}, \eng{cash/money}, \eng{price/cost/value}).
\item Réemployer ce lexique dans des phrases et des dialogues.
\end{enumerate}
\end{methode}

\courssub{Reading: Mr Ndi's monthly budget}

Lisons d'abord ce court texte. Soulignez mentalement tous les mots liés à l'argent.

\begin{textebox}[Mr Ndi's money in and money out]
Every month, Mr Ndi receives his salary of 150,000 francs CFA. He pays his rent and his electricity bills, puts some money into his savings account, and gives pocket money to his children. What is left is for food and transport.

Mr Ndi is a careful man. At the beginning of each month, he sits at his table in Yaoundé and writes down his budget. First, he lists his income: his salary and a small profit from the tomatoes his wife sells at the market. Then he lists his expenses. The rent is 40,000 francs. The electricity and water bills come to about 15,000 francs. School fees for his three children are his biggest expenditure. He also pays taxes to the council.

Mr Ndi never borrows money if he can avoid it. “A loan means debt,” he always says, “and debt means interest.” Last year, his neighbour had to pay a heavy fine because he parked his car in the wrong place in town. Mr Ndi prefers to keep some cash at home for emergencies, but most of his money stays in his bank account.

At the end of the month, there is sometimes a little money left. Mr Ndi puts it into his savings. “One day,” he smiles, “I will buy a small piece of land in the village.” His children have learned the lesson too: they save part of their pocket money in a wooden box, and they count their coins every Saturday evening.
\end{textebox}

\textbf{Glossaire}\par
\begin{itemize}
\item \cle{salary} (n.) : salaire (mensuel)
\item \cle{rent} (n.) : loyer
\item \cle{bill} (n.) : facture
\item \cle{savings account} (n.) : compte d'épargne
\item \cle{pocket money} (n.) : argent de poche
\item \cle{budget} (n.) : budget, plan de dépenses
\item \cle{income} (n.) : revenu(s)
\item \cle{profit} (n.) : bénéfice
\item \cle{expense} (n.) : dépense
\item \cle{school fees} (n. pl.) : frais de scolarité
\item \cle{expenditure} (n.) : dépenses (terme soutenu)
\item \cle{tax} (n.) : impôt, taxe
\item \cle{loan} (n.) : prêt, emprunt
\item \cle{debt} (n.) : dette
\item \cle{interest} (n.) : intérêts (d'un prêt)
\item \cle{fine} (n.) : amende
\item \cle{cash} (n.) : espèces, liquide
\item \cle{coin} (n.) : pièce de monnaie
\end{itemize}

\textbf{Comprehension questions}\par
\begin{enumerate}
\item How much is Mr Ndi's monthly salary?
\item Name three of Mr Ndi's expenses.
\item Why does Mr Ndi dislike loans?
\item What does he do with the money left at the end of the month?
\item Where do his children keep their savings?
\end{enumerate}

\courssub{Lexical field 1: money in — l'argent qui entre}

Observons les phrases suivantes tirées ou inspirées du texte :

\begin{exemplebox}[Observe and compare]
\eng{Mr Ndi receives his salary of 150,000 francs CFA.} « M. Ndi reçoit son salaire de 150 000 francs CFA. »

\eng{His wife makes a small profit from the tomatoes she sells.} « Sa femme fait un petit bénéfice sur les tomates qu'elle vend. »

\eng{He gives pocket money to his children.} « Il donne de l'argent de poche à ses enfants. »

\eng{First, he lists his income.} « D'abord, il dresse la liste de ses revenus. »
\end{exemplebox}

\begin{definition}[Money in — les mots du revenu]
\begin{itemize}
\item \cle{income} (indénombrable) : l'ensemble des revenus d'une personne ou d'un ménage. \eng{What is your monthly income?} « Quel est votre revenu mensuel ? »
\item \cle{salary} (dénombrable) : salaire fixe, versé en général \textbf{chaque mois}, pour un travail de bureau ou une profession qualifiée. \eng{Teachers receive their salary at the end of the month.}
\item \cle{wages} (pluriel) : salaire payé \textbf{à l'heure, à la journée ou à la semaine}, souvent pour un travail manuel. \eng{The workers collected their wages on Friday.} « Les ouvriers ont touché leur salaire vendredi. »
\item \cle{earnings} (pluriel) : ce que l'on gagne globalement par son travail. \eng{Her earnings increased last year.}
\item \cle{pocket money} (indénombrable, GB) : argent de poche donné aux enfants. Équivalent américain : \cle{allowance}. \eng{Aminatou gets 10,000 francs pocket money every month.}
\item \cle{profit} : bénéfice, gain réalisé dans une vente ou une affaire. \eng{They sold the cocoa at a profit.} « Ils ont vendu le cacao avec bénéfice. »
\end{itemize}
\end{definition}

\begin{propriete}[Salary ou wages ?]
La distinction est sociale et rythmique :
\begin{itemize}
\item \textbf{salary} = somme fixe, versée \emph{mensuellement}, typique des emplois de bureau et des professions intellectuelles : \eng{a nurse's salary}, \eng{an annual salary of two million francs}.
\item \textbf{wages} = paye calculée \emph{à l'heure ou à la semaine}, typique du travail manuel ou temporaire : \eng{daily wages}, \eng{the minimum wage} (le salaire minimum).
\end{itemize}
On dit \eng{to earn a salary} et \eng{to earn wages}, mais jamais \eng{to win a salary} : attention, \eng{win} signifie « gagner » à un jeu ou à un concours (\eng{win the lottery}), pas « gagner sa vie ».
\end{propriete}

\begin{center}
\begin{tabularx}{\linewidth}{|X|X|X|}\hline
\rowcolor{popblueL} Français & English & Remarque \\ \hline
salaire (mensuel, bureau) & a salary & \eng{He earns a good salary.} \\ \hline
salaire (horaire/hebdo, manuel) & wages (pluriel) & \eng{weekly wages} \\ \hline
revenu(s) & income (indénombrable) & pas de pluriel \\ \hline
argent de poche & pocket money (GB) / allowance (US) & indénombrable \\ \hline
bénéfice & profit & \eng{to make a profit} \\ \hline
gains, revenus du travail & earnings & toujours au pluriel \\ \hline
\end{tabularx}
\end{center}

\courssub{Lexical field 2: money out — l'argent qui sort}

\begin{definition}[Money out — les mots de la dépense]
\begin{itemize}
\item \cle{expenses} (pluriel) : les dépenses courantes. \eng{Food and transport are my main expenses.} « La nourriture et le transport sont mes principales dépenses. »
\item \cle{expenditure} (indénombrable, soutenu) : l'ensemble des dépenses, terme officiel. \eng{Government expenditure on education.}
\item \cle{bill} : facture (électricité, eau, téléphone). \eng{He pays his electricity bills.} « Il paie ses factures d'électricité. »
\item \cle{fees} (pluriel) : frais de service ou de scolarité. \eng{school fees} « frais de scolarité », \eng{lawyer's fees} « honoraires d'avocat ».
\item \cle{rent} : loyer. \eng{The rent is 40,000 francs a month.}
\item \cle{taxes} : impôts et taxes (payés à l'État). \eng{Traders pay taxes to the council.}
\item \cle{fine} : amende (punition financière). \eng{He had to pay a heavy fine for dangerous parking.} « Il a dû payer une lourde amende pour un stationnement dangereux. »
\end{itemize}
\end{definition}

\begin{attention}[Faux amis et pièges]
\begin{itemize}
\item \eng{a fine} = une \textbf{amende}, pas « fin » ni « bien ». \eng{a parking fine} = une amende pour stationnement interdit.
\item \eng{fees} = frais de service/scolarité ; ne pas confondre avec \eng{free} (gratuit) : \eng{The museum is free, but there are fees for the guided tour.}
\item \eng{bill} = facture ; au restaurant britannique, \eng{the bill} = l'addition (américain : \eng{the check}).
\item \eng{expenses} prend toujours un \textbf{s} dans le sens de « dépenses courantes » : \eng{travelling expenses} « frais de déplacement ».
\end{itemize}
\end{attention}

\courssub{Lexical field 3: managing money — gérer son argent}

\begin{definition}[Budget]
\cle{budget} (n.) : \eng{a plan of how to spend the money you have over a period of time} — un plan de dépenses sur une période donnée. C'est aussi un verbe : \eng{Aminatou must budget her pocket money carefully.} « Aminatou doit gérer son argent de poche avec soin. »
\end{definition}

\begin{definition}[Managing money — les mots de la gestion]
\begin{itemize}
\item \cle{savings} (pluriel) : économies, épargne. \eng{He puts money into his savings account.} « Il place de l'argent sur son compte d'épargne. »
\item \cle{account} : compte (en banque). \eng{a bank account} « un compte bancaire », \eng{to open an account} « ouvrir un compte ».
\item \cle{loan} : prêt, emprunt. \eng{to take out a loan} « contracter un prêt ».
\item \cle{debt} : dette. \eng{to be in debt} « être endetté ». Attention : le \textbf{b} ne se prononce pas — « dèt ».
\item \cle{interest} (indénombrable) : intérêts d'un prêt ou d'un placement. \eng{The bank charges 10\% interest.}
\item \cle{change} (indénombrable) : la monnaie (rendue ou petite monnaie). \eng{Keep the change!} « Gardez la monnaie ! »
\item \cle{cash} (indénombrable) : espèces, liquide. \eng{to pay in cash} « payer en espèces ».
\end{itemize}
\end{definition}

\begin{attention}[Money est INDÉNOMBRABLE]
En anglais, \cle{money} est \textbf{indénombrable} : on dit \eng{some money}, \eng{a lot of money}, \eng{much money} — jamais \eng{a money} ni \eng{moneys}. Pour compter, on utilise des noms dénombrables : \eng{coins} (pièces) et \eng{banknotes} (GB) / \eng{bills} (US) (billets).

\eng{How much money do you have? — I have a few coins and two banknotes.} « Combien d'argent as-tu ? — J'ai quelques pièces et deux billets. »

De même : \eng{cash}, \eng{income}, \eng{pocket money}, \cle{change} et \cle{interest} sont indénombrables.
\end{attention}

\begin{propriete}[Cash ou money ?]
\begin{itemize}
\item \textbf{money} = terme général, indénombrable : \eng{I need money to pay the rent.}
\item \textbf{cash} = uniquement les \emph{espèces} (billets et pièces), par opposition au chèque ou à la carte : \eng{Can I pay in cash?} « Puis-je payer en liquide ? », \eng{a cash machine} « un distributeur de billets ».
\end{itemize}
On ne dit donc pas \eng{I have no cashs} : \eng{cash} reste indénombrable.
\end{propriete}

\courssub{Price, cost, value, worth: le prix et la valeur}

\begin{definition}[Quatre mots pour parler du prix]
\begin{itemize}
\item \cle{price} : le prix affiché d'un article précis. \eng{What is the price of this phone?} « Quel est le prix de ce téléphone ? »
\item \cle{cost} : le coût (ce que quelque chose coûte réellement) ; aussi un verbe. \eng{the cost of living} « le coût de la vie » ; \eng{This bag costs 5,000 francs.}
\item \cle{value} : la valeur (ce que quelque chose vaut). \eng{good value for money} « un bon rapport qualité-prix ». \eng{This market offers good value for money.}
\item \cle{worth} (adjectif, suivi d'un nom ou d'un gérondif) : qui vaut. \eng{The land is worth two million francs.} « Le terrain vaut deux millions de francs. » \eng{Is it worth buying?} « Cela vaut-il la peine d'être acheté ? »
\end{itemize}
\end{definition}

\begin{exemplebox}[Exemples traduits]
\eng{She earns a good salary.} « Elle gagne un bon salaire. »

\eng{I can't afford this phone.} « Je ne peux pas me permettre d'acheter ce téléphone. »

\eng{The cost of living is rising in Douala.} « Le coût de la vie augmente à Douala. »

\eng{These shoes are good value for money.} « Ces chaussures ont un bon rapport qualité-prix. »

\eng{He pays his rent in cash every month.} « Il paie son loyer en espèces chaque mois. »
\end{exemplebox}

\courssub{Mind map: the four worlds of personal finances}

\begin{popfigure}
\begin{center}
\begin{tikzpicture}[mindmap, grow cyclic, every node/.style={concept, align=center, font=\small},
root concept/.append style={concept color=popblue, font=\bfseries},
level 1/.append style={level distance=3.4cm, sibling angle=90},
level 2/.append style={level distance=2.2cm, sibling angle=45, font=\footnotesize}]
\node[root concept]{PERSONAL FINANCES}
  child[concept color=popgreen!70]{ node{EARNING}
      child{ node{salary} }
      child{ node{wages} }
      child{ node{income} }
      child{ node{profit} } }
  child[concept color=poporange!80]{ node{SPENDING}
      child{ node{bills} }
      child{ node{rent} }
      child{ node{fees} }
      child{ node{taxes} } }
  child[concept color=popgold!80]{ node{SAVING}
      child{ node{savings} }
      child{ node{account} }
      child{ node{budget} } }
  child[concept color=poppurple!70]{ node{BORROWING}
      child{ node{loan} }
      child{ node{debt} }
      child{ node{interest} } };
\end{tikzpicture}
\end{center}
\legende{Carte mentale du vocabulaire des finances personnelles : quatre branches, quatre familles de mots.}
\end{popfigure}

\begin{aretenir}[À retenir]
\begin{itemize}
\item \eng{salary} (mensuel, bureau) $\neq$ \eng{wages} (à l'heure/semaine, manuel).
\item \eng{money}, \eng{cash}, \eng{income}, \eng{pocket money} sont \textbf{indénombrables}.
\item \eng{a fine} = une amende ; \eng{fees} = frais ; \eng{a bill} = une facture.
\item \eng{price} (prix d'un article), \eng{cost} (coût), \eng{value} (valeur), \eng{worth} (qui vaut).
\item \eng{debt} se prononce « dèt » : le \textbf{b} est muet.
\end{itemize}
\end{aretenir}

\begin{exoresolu}[Vocabulary in context]
Énoncé — Complétez les phrases avec le mot qui convient : \emph{salary, wages, fine, cash, budget, loan, change, income}.

\begin{enumerate}
\item The builders receive their \ldots\ every Friday.
\item Aminatou makes a weekly \ldots\ so that her pocket money lasts the whole month.
\item “Sorry, we don't accept cards. You must pay in \ldots .”
\item The driver paid a \ldots\ of 25,000 francs for driving without a licence.
\item My father's total \ldots\ comes from his job and his farm.
\item Here is your \ldots : 500 francs. Thank you, madam!
\end{enumerate}
\tcblower
Solution détaillée :
\begin{enumerate}
\item \textbf{wages} — payés chaque semaine (\eng{every Friday}) pour un travail manuel : c'est la définition de \eng{wages}.
\item \textbf{budget} — un plan de dépenses sur une période donnée ; on dit \eng{to make a budget}.
\item \textbf{cash} — par opposition à la carte bancaire (\eng{cards}), il faut payer en espèces.
\item \textbf{fine} — une amende, punition financière pour une infraction.
\item \textbf{income} — l'ensemble des revenus (indénombrable).
\item \textbf{change} — la monnaie rendue par le commerçant.
\end{enumerate}
\end{exoresolu}

\begin{exercice}[Your turn]
Traduisez en anglais :
\begin{enumerate}
\item Les enseignants reçoivent leur salaire à la fin du mois.
\item Je ne peux pas me permettre d'acheter ces chaussures.
\item Il paie son loyer et ses factures d'électricité en espèces.
\item Combien d'argent mets-tu sur ton compte d'épargne chaque mois ?
\end{enumerate}
\end{exercice}

\begin{corrige}[Exercice « Your turn »]
\begin{enumerate}
\item \eng{Teachers receive their salary at the end of the month.}
\item \eng{I can't afford these shoes.} (ou \eng{I can't afford to buy these shoes.})
\item \eng{He pays his rent and his electricity bills in cash.}
\item \eng{How much money do you put into your savings account every month?}
\end{enumerate}
\end{corrige}

\begin{savaistu}[Njangi and tontines]
Au Cameroun anglophone, l'épargne rotative entre amis ou voisins s'appelle \eng{njangi} ; on parle aussi de \eng{tontine} dans tout le pays. En anglais standard, on dit \eng{a rotating savings club} ou \eng{a savings circle} : chaque membre cotise chaque mois, et à tour de rôle, l'un d'eux reçoit la totalité de la cagnotte. \eng{My mother belongs to a savings circle with ten other women.} « Ma mère fait partie d'une tontine avec dix autres femmes. » C'est une excellente façon de pratiquer le vocabulaire de cette leçon : \eng{to save, to contribute, to receive, to share} !
\end{savaistu}

\coursec{Key verbs: earn, spend, save, borrow, lend}

Dans la section précédente, nous avons appris à nommer les flux d'argent : \eng{money in} (salaire, revenus, pourboires) et \eng{money out} (dépenses, factures, dettes). Mais nommer ne suffit pas : pour raconter la vie économique d'une personne, il faut des \cle{verbes}. Comment dire « gagner », « dépenser », « épargner », « emprunter », « prêter » en anglais correct ? C'est l'objet de cette section. Ces verbes sont parmi les plus fréquents de la langue — et parmi les plus maltraités par les francophones, car leurs constructions ne se calquent pas sur le français.

\courssub{Observation : un texte pour découvrir les verbes en action}

Lisons d'abord ce court texte. Soulignez mentalement tous les verbes qui parlent d'argent.

\begin{textebox}[Mama Ndi's monthly budget]
Mama Ndi sells tomatoes and onions at the Mokolo market in Yaoundé. She \textbf{earns} about 90,000 francs a month — not a fortune, but honest money. Every evening, she counts what she has made and writes it in a small notebook.

She \textbf{spends} most of her money \textbf{on} food, school fees and transport. She refuses to \textbf{waste} money \textbf{on} things she does not need. “A phone with a broken screen still makes calls,” she laughs.

Last year, she decided to \textbf{save up for} a deep freezer, so that she could also sell frozen fish. Every week, she put 5,000 francs into her savings box. When she had almost enough, her freezer supplier raised the price. So she \textbf{borrowed} 30,000 francs \textbf{from} her cousin Bertille, who agreed to \textbf{lend} her the money for three months, without interest. Mama Ndi now \textbf{owes} Bertille 30,000 francs, but she is not worried: with the new freezer, her income has grown, and she will \textbf{pay} her cousin \textbf{back} before the end of the term.

Her son Etienne, a student, once asked her: “Mum, can you \textbf{lend} me 10,000 francs? I want to buy new trainers.” She answered: “I \textbf{can't afford} to give money away, my son. But if you help me at the stall on Saturdays, you will \textbf{earn} it yourself.” Etienne worked for four Saturdays, earned his 10,000 francs, and understood a lesson no schoolbook could teach him: money you earn yourself is the money you respect the most.
\end{textebox}

\textbf{Glossaire}

\begin{center}
\begin{tabularx}{\linewidth}{|l|l|X|}
\hline
\rowcolor{popblueL} Mot anglais & Nature & Traduction française \\
\hline
school fees & nom pluriel & les frais de scolarité \\
\hline
a deep freezer & nom & un congélateur \\
\hline
a stall & nom & un étal (au marché) \\
\hline
interest & nom & les intérêts (d'un prêt) \\
\hline
trainers & nom pluriel & des baskets (britannique) \\
\hline
to raise the price & expression & augmenter le prix \\
\hline
\end{tabularx}
\end{center}

\textbf{Comprehension questions}

\begin{enumerate}
\item How much does Mama Ndi earn every month?
\item What does she spend most of her money on?
\item What was she saving up for, and why?
\item Who did she borrow money from? How much does she owe?
\item How did Etienne finally get his 10,000 francs?
\end{enumerate}

\courssub{Earn vs win : deux façons de « gagner »}

Le français utilise un seul verbe, \emph{gagner}, pour deux réalités très différentes. L'anglais, lui, en impose deux.

\begin{propriete}[EARN ou WIN ?]
\begin{itemize}
\item \cle{earn} = gagner de l'argent \emph{par son travail, son mérite, ses efforts}.
\eng{She earns 90,000 francs a month.} « Elle gagne 90 000 francs par mois. »
\item \cle{win} = gagner \emph{au jeu, à la loterie, à un concours, à un match}.
\eng{He won 50,000 francs in the lottery.} « Il a gagné 50 000 francs à la loterie. »
\end{itemize}
Formes : \eng{earn – earned – earned} (régulier) ; \eng{win – won – won} (irrégulier ; « won » se prononce « ouann »).
\end{propriete}

\begin{attention}[Le piège du verbe « gagner »]
Ne dites jamais \eng{He won a good salary.} Un salaire ne se « gagne » pas à la chance : il se mérite. Dites \eng{He earns a good salary.} De même, on dit \eng{earn a living} (« gagner sa vie »), jamais \eng{win a living} : \eng{She earns a living by selling vegetables.} « Elle gagne sa vie en vendant des légumes. »
\end{attention}

\courssub{Spend : dépenser de l'argent, passer du temps}

\begin{propriete}[Les constructions de SPEND]
\cle{spend} (spend – spent – spent) a deux constructions à connaître par cœur :
\begin{itemize}
\item \textbf{spend + argent/temps + ON + nom} : \eng{He spends a lot of money on clothes.} « Il dépense beaucoup d'argent en vêtements. »
\item \textbf{spend + temps + (IN) + verbe en -ing} : \eng{He spent the evening counting his money.} « Il a passé la soirée à compter son argent. »
\end{itemize}
Remarquez : le français dit « dépenser \emph{en} vêtements », l'anglais dit \eng{spend ON}. La préposition \eng{in} devant le verbe en -ing est facultative et souvent omise : \eng{She spent two hours (in) bargaining at the market.}
\end{propriete}

\begin{attention}[Jamais « spend money for »]
L'influence du français « dépenser pour » pousse à écrire \eng{I spent 2,000 francs for data.} C'est incorrect. Dites : \eng{I spent 2,000 francs on data.} « J'ai dépensé 2 000 francs en forfait internet. »
\end{attention}

\courssub{Save et save up for : épargner en vue d'un objectif}

\begin{propriete}[SAVE et SAVE UP FOR]
\begin{itemize}
\item \cle{save money} = mettre de l'argent de côté, épargner : \eng{She saves 5,000 francs every week.} « Elle épargne 5 000 francs chaque semaine. »
\item \cle{save up FOR something} = économiser \emph{en vue d'acheter quelque chose de précis} : \eng{She is saving up for a laptop.} « Elle économise pour s'acheter un ordinateur portable. »
\end{itemize}
Attention : \eng{save} signifie aussi « sauver » (\eng{save a life}) et « économiser du temps » (\eng{save time}). Le contexte décide.
\end{propriete}

\courssub{Borrow, lend, owe : le triangle du prêt}

C'est le point le plus délicat de la leçon. Le français possède le verbe \emph{prêter} qui s'utilise dans les deux sens familièrement (« il m'a prêté » / « j'ai prêté »), mais l'anglais distingue strictement les deux directions du mouvement.

\begin{propriete}[BORROW ... FROM / LEND ... TO]
\begin{itemize}
\item \cle{borrow something FROM somebody} = emprunter quelque chose \emph{à} quelqu'un (on \textbf{reçoit}).
\eng{I borrowed 5,000 francs from my brother.} « J'ai emprunté 5 000 francs à mon frère. »
\item \cle{lend something TO somebody} = prêter quelque chose \emph{à} quelqu'un (on \textbf{donne}).
\eng{My brother lent me 5,000 francs.} « Mon frère m'a prêté 5 000 francs. »
\end{itemize}
Les deux phrases décrivent \textbf{la même action vue de deux côtés opposés} : celui qui reçoit (\eng{borrow}) et celui qui donne (\eng{lend}). Formes : \eng{borrow – borrowed – borrowed} (régulier) ; \eng{lend – lent – lent} (irrégulier).
\end{propriete}

\begin{popfigure}
\begin{center}
\begin{tikzpicture}[node distance=4.2cm]
\node[draw, circle, fill=popblueL, minimum size=1.2cm, font=\bfseries] (A) {A};
\node[draw, circle, fill=poporangeL, minimum size=1.2cm, font=\bfseries, right=of A] (B) {B};
\draw[-{Stealth[length=3mm]}, very thick, popblue] (A) to[bend left=25] node[above, align=center, font=\small]{A \textbf{lends} money \textbf{to} B} (B);
\draw[-{Stealth[length=3mm]}, very thick, poporange] (B) to[bend left=25] node[below, align=center, font=\small]{B \textbf{borrows} money \textbf{from} A\\ B \textbf{owes} money \textbf{to} A} (A);
\end{tikzpicture}
\legende{Le même prêt vu des deux côtés : A prête (lend), B emprunte (borrow) et B doit de l'argent à A (owe).}
\end{center}
\end{popfigure}

\begin{propriete}[OWE : devoir de l'argent]
\cle{owe somebody money} / \cle{owe money TO somebody} = devoir de l'argent à quelqu'un (être le débiteur).
\eng{B owes A 30,000 francs.} = \eng{B owes 30,000 francs to A.} « B doit 30 000 francs à A. »
On peut aussi dire \eng{owe somebody a favour} : « devoir un service à quelqu'un ».
\end{propriete}

\begin{attention}[L'erreur n°1 des francophones]
Ne dites \textbf{jamais} \eng{I borrowed him money} pour dire « je lui ai prêté de l'argent ». \eng{borrow} signifie \emph{recevoir}, pas \emph{donner}. Pour « je lui ai prêté », il faut : \eng{I lent him money.} ou \eng{I lent money to him.} Testez-vous : qui reçoit ? Si c'est le sujet, c'est \eng{borrow} ; si c'est l'autre personne, c'est \eng{lend}.
\end{attention}

\courssub{Afford : avoir les moyens}

\begin{propriete}[CAN/CAN'T + AFFORD]
\cle{afford} = avoir les moyens (financiers ou temporels) de faire quelque chose. Il s'emploie \textbf{presque toujours avec can, could ou be able to}, très souvent à la forme négative :
\begin{itemize}
\item \eng{I can't afford it.} « Je n'ai pas les moyens (de l'acheter). »
\item \eng{We can't afford a car this year.} « Nous ne pouvons pas nous offrir une voiture cette année. »
\item \eng{Can you afford to miss three days of work?} « Peux-tu te permettre de manquer trois jours de travail ? »
\end{itemize}
Construction : \eng{afford + nom} ou \eng{afford TO + infinitif}. Jamais de \eng{to} après \eng{can} : \eng{I can't afford to buy it}, mais \eng{I can't afford it}.
\end{propriete}

\begin{savaistu}[D'où vient « afford » ?]
\eng{Afford} descend du vieil anglais \emph{geforthian}, qui signifiait « mener à bien, accomplir, fournir ». Au fil des siècles, le sens s'est spécialisé : « avoir les ressources suffisantes pour ». C'est pourquoi le verbe a gardé un goût de « capacité » et refuse de vivre sans son compagnon \eng{can} ou \eng{could} : dire \eng{I afford it} tout seul sonne très étrange à une oreille anglophone. Retenez le duo inséparable : \textbf{can + afford}.
\end{savaistu}

\courssub{Waste vs invest ; pay, pay for, pay back}

\begin{propriete}[WASTE ... ON / INVEST ... IN]
\begin{itemize}
\item \cle{waste money ON something} = gaspiller de l'argent \emph{pour} quelque chose : \eng{Don't waste your money on cheap shoes that break in a week.} « Ne gaspille pas ton argent en chaussures bon marché qui cassent en une semaine. »
\item \cle{invest money IN something} = investir de l'argent \emph{dans} quelque chose : \eng{She invested her savings in her shop.} « Elle a investi ses économies dans sa boutique. »
\end{itemize}
Même argent, deux destins : \eng{waste} (l'argent disparaît sans retour) et \eng{invest} (l'argent doit rapporter). Notez les prépositions différentes : \eng{ON} pour le gaspillage, \eng{IN} pour l'investissement.
\end{propriete}

\begin{propriete}[PAY, PAY FOR, PAY BACK]
\begin{itemize}
\item \cle{pay + une somme / une facture / une personne} : \eng{pay a bill} « payer une facture » ; \eng{pay the mechanic} « payer le mécanicien » ; \eng{pay 2,000 francs} « payer 2 000 francs ».
\item \cle{pay FOR + l'objet acheté} : \eng{I paid 15,000 francs for this bag.} « J'ai payé ce sac 15 000 francs. »
\item \cle{pay somebody BACK} = rembourser quelqu'un : \eng{She will pay her cousin back before December.} « Elle remboursera sa cousine avant décembre. »
\end{itemize}
Formes : \eng{pay – paid – paid} (attention à l'orthographe : \eng{paid}, pas \eng{payed}).
\end{propriete}

\begin{attention}[Pay ou pay for ?]
On paie \emph{une personne}, \emph{une somme} ou \emph{une facture} sans préposition ; on paie \emph{pour un objet} avec \eng{for}. Comparez : \eng{I paid the taxi driver.} « J'ai payé le chauffeur de taxi. » / \eng{I paid for the taxi ride.} « J'ai payé la course de taxi. » Et ne confondez pas \eng{pay} (payer) avec \eng{pay attention} (« faire attention ») !
\end{attention}

\courssub{Tableau récapitulatif des verbes}

\begin{center}
\begin{tabularx}{\linewidth}{|l|X|X|}
\hline
\rowcolor{popblueL} Verbe anglais & Construction et exemple & Traduction française \\
\hline
earn & earn money / earn a living — \eng{She earns a good salary.} & gagner (par le travail) \\
\hline
win & win money (game, lottery) — \eng{He won 10,000 francs.} & gagner (au jeu) \\
\hline
spend & spend money ON sth / spend time DOING sth & dépenser ; passer (du temps) \\
\hline
save (up) & save money / save up FOR sth & épargner ; économiser pour \\
\hline
borrow & borrow sth FROM sb & emprunter (à) — on reçoit \\
\hline
lend & lend sth TO sb / lend sb sth & prêter (à) — on donne \\
\hline
owe & owe sb money / owe money to sb & devoir de l'argent (à) \\
\hline
afford & can/can't afford sth / to do sth & avoir les moyens de \\
\hline
waste & waste money ON sth & gaspiller \\
\hline
invest & invest money IN sth & investir \\
\hline
pay (back) & pay a bill / pay FOR sth / pay sb BACK & payer ; rembourser \\
\hline
\end{tabularx}
\end{center}

\courssub{Exercice résolu et entraînement}

\begin{exoresolu}[Choose the correct verb]
Complétez avec \eng{earn, win, spend, save, borrow, lend, owe, afford, waste, pay back} à la forme correcte.

1. My uncle \ldots\ a lot of money in the casino last night — what a surprise!

2. « Can you \ldots\ me your bicycle for an hour? » — « Yes, but bring it back before six. »

3. Amina's father works two jobs; he \ldots\ about 150,000 francs a month.

4. I can't \ldots\ a new phone this term; my old one must survive.

5. Stop \ldots\ your pocket money on sweets!

6. Paul borrowed 20,000 francs from me in March and still hasn't \ldots\ me \ldots.
\tcblower
\textbf{Solutions détaillées}

1. \eng{won} — il s'agit d'un jeu de hasard (casino), donc \eng{win}, au prétérit (\eng{last night}) : \eng{won}.

2. \eng{lend} — le locuteur \emph{donne} le vélo temporairement ; après \eng{can}, base verbale : \eng{lend}. (Si la phrase était « Can I \ldots\ your bicycle? », ce serait \eng{borrow}, car « je » reçoit.)

3. \eng{earns} — revenu du travail, présent simple, 3e personne : \eng{earns}.

4. \eng{afford} — construction obligatoire \eng{can't afford + nom}.

5. \eng{wasting} — après \eng{stop}, verbe en -ing ; « gaspiller » : \eng{wasting}.

6. \eng{paid ... back} — « rembourser » : \eng{pay somebody back} ; au present perfect : \eng{hasn't paid me back}.
\end{exoresolu}

\begin{exercice}[Your turn]
Traduisez en anglais :

1. Mon frère m'a prêté 10 000 francs hier.

2. Elle économise pour s'acheter une machine à coudre.

3. Il dépense tout son argent en jeux vidéo.

4. Nous n'avons pas les moyens d'acheter une voiture cette année.

5. Je dois 5 000 francs à ma voisine ; je la rembourserai samedi.
\end{exercice}

\begin{corrige}[Exercice « Your turn »]
1. \eng{My brother lent me 10,000 francs yesterday.} (ou \eng{lent 10,000 francs to me}) — jamais \eng{borrowed} ici : c'est le frère qui donne.

2. \eng{She is saving up for a sewing machine.}

3. \eng{He spends all his money on video games.} — préposition \eng{on}, pas \eng{for}.

4. \eng{We can't afford a car this year.} (ou \eng{We can't afford to buy a car this year.})

5. \eng{I owe my neighbour 5,000 francs; I will pay her back on Saturday.}
\end{corrige}

\begin{aretenir}[L'essentiel de la section]
\begin{itemize}
\item \eng{earn} = par le travail ; \eng{win} = au jeu.
\item \eng{spend money ON sth} ; \eng{spend time (IN) DOING sth}.
\item \eng{save up FOR sth} = économiser en vue de.
\item \eng{borrow FROM} (recevoir) / \eng{lend TO} (donner) / \eng{owe TO} (devoir).
\item \eng{afford} ne vit qu'avec \eng{can/could} : \eng{I can't afford it.}
\item \eng{waste ON} (gaspiller) / \eng{invest IN} (investir) ; \eng{pay FOR} un objet, \eng{pay sb BACK} (rembourser).
\end{itemize}
\end{aretenir}

\coursec{Collocations and idiomatic expressions}

Dans la section précédente, nous avons étudié les verbes clés \eng{earn}, \eng{spend}, \eng{save}, \eng{borrow} et \eng{lend}. Mais un verbe seul ne suffit pas pour parler naturellement : en anglais comme en français, les mots ont des \cle{partenaires habituels}. On dit \eng{make money} et non \eng{do money}, \eng{pay off a debt} et non \eng{finish a debt}. Ces combinaisons naturelles s'appellent des \cle{collocations}. Nous verrons ensuite les \cle{expressions idiomatiques}, ces images figées dont le sens ne se devine pas mot à mot.

\courssub{What is a collocation?}

\begin{definition}[Collocation]
Une \cle{collocation} est une combinaison de mots qui « vont ensemble » naturellement dans une langue : verbe + nom (\eng{earn money}), adjectif + nom (\eng{a tight budget}), verbe + préposition. On ne peut pas la deviner à partir du français : il faut l'apprendre comme un bloc.
\end{definition}

Observez d'abord ce court texte, puis relevez les collocations qu'il contient.

\begin{textebox}[Tabi's tight month — a reading text]
After paying his children's school fees, Tabi is completely broke. He has to tighten his belt and stick to a strict budget until the end of the month. As his mother always says: money doesn't grow on trees!

Tabi is a taxi driver in Douala. Last year, he decided to draw up a family budget with his wife, Marthe. They listed all their expenses: food, transport, school fees, rent and mobile data. At first, it was difficult to stick to the budget, because life in the city is expensive and the children always need something new. Once, Tabi ran up a debt at the provision store near the market. It took him three months to pay it off, and he promised himself he would never be in debt again.

Now the family tries to put money aside every month. Marthe opened a savings account at a microfinance bank, and they only dip into their savings for real emergencies, like hospital bills. “We are not rich,” Tabi says, “but we live within our means, and we manage to make ends meet. One day, our little nest egg will pay for a plot of land in the village.”

\emph{Glossaire :} fees (n. pl.) : frais (scolaires) ; broke (adj.) : fauché ; to tighten one's belt : se serrer la ceinture ; to stick to (v.) : s'en tenir à ; to run up a debt : accumuler une dette ; to pay off (v.) : rembourser ; to put aside : mettre de côté ; to dip into : puiser dans ; nest egg (n.) : pécule ; plot of land (n.) : terrain.
\end{textebox}

\textbf{Comprehension questions}
\begin{enumerate}
\item Why is Tabi broke at the beginning of the text?
\item What did Tabi and Marthe do to control their expenses?
\item How long did it take Tabi to pay off his debt at the provision store?
\item When do they dip into their savings?
\end{enumerate}

\courssub{Collocations with money, budget, debt and savings}

\begin{propriete}[Les quatre familles de collocations]
Les collocations se mémorisent par \cle{nom-clé}. Autour de chaque nom (\eng{money}, \eng{budget}, \eng{debt}, \eng{savings}), un petit nombre de verbes et d'adjectifs reviennent toujours.
\end{propriete}

\begin{center}
\begin{tabularx}{\linewidth}{|l|X|X|}
\hline
\rowcolor{popblueL} Collocation & Traduction & Exemple \\
\hline
\eng{make money} & gagner / faire de l'argent & \eng{He makes money by repairing phones.} \\
\hline
\eng{earn money} & gagner de l'argent (par son travail) & \eng{She earns money as a nurse.} \\
\hline
\eng{save money} & économiser de l'argent & \eng{We save money every month.} \\
\hline
\eng{spend money (on sth)} & dépenser de l'argent (pour qqch) & \eng{Don't spend money on useless things.} \\
\hline
\eng{waste money} & gaspiller de l'argent & \eng{He wasted his money on gambling.} \\
\hline
\eng{owe money (to sb)} & devoir de l'argent (à qqn) & \eng{I owe 5,000 francs to my brother.} \\
\hline
\eng{change money} & changer de l'argent (devises) & \eng{You can change money at the bank.} \\
\hline
\end{tabularx}
\end{center}

\begin{center}
\begin{tabularx}{\linewidth}{|l|X|X|}
\hline
\rowcolor{popgreenL} Collocation & Traduction & Exemple \\
\hline
\eng{draw up a budget} & établir un budget & \eng{They drew up a budget for the term.} \\
\hline
\eng{stick to a budget} & respecter un budget & \eng{It is hard to stick to a tight budget.} \\
\hline
\eng{cut the budget} & réduire le budget & \eng{The school cut its budget this year.} \\
\hline
\eng{a tight budget} & un budget serré & \eng{We live on a tight budget.} \\
\hline
\eng{run up a debt} & accumuler une dette & \eng{He ran up a debt at the shop.} \\
\hline
\eng{pay off a debt} & rembourser une dette & \eng{She paid off her debt in June.} \\
\hline
\eng{be in debt} & être endetté & \eng{Many traders are in debt.} \\
\hline
\eng{get out of debt} & sortir de l'endettement & \eng{He worked hard to get out of debt.} \\
\hline
\eng{open a savings account} & ouvrir un compte d'épargne & \eng{Marthe opened a savings account.} \\
\hline
\eng{put money aside} & mettre de l'argent de côté & \eng{Put some money aside for emergencies.} \\
\hline
\eng{dip into one's savings} & puiser dans ses économies & \eng{We dipped into our savings for the fees.} \\
\hline
\end{tabularx}
\end{center}

\begin{attention}[Pièges pour les francophones]
\begin{itemize}
\item On dit \eng{make money} et non \eng{do money} : \eng{He makes good money.} « Il gagne bien sa vie. »
\item On dit \eng{spend money ON something} (préposition \eng{on}) : \eng{She spends money on clothes.} — jamais \eng{spend money for}.
\item \eng{to be in debt} = être endetté. Ne dites pas \eng{I have a debt with him} pour « je lui dois de l'argent » ; dites \eng{I owe him money}.
\item \eng{to change money} signifie « changer de la monnaie » (euros contre francs CFA), pas « changer d'argent » au sens de modifier.
\item \eng{borrow} (emprunter) $\neq$ \eng{lend} (prêter) : \eng{I borrow from the bank}, \eng{The bank lends to me}. Le pronom suit le verbe : \eng{lend me 1,000 francs} = « prête-moi », \eng{borrow 1,000 francs from me} = « emprunte-moi ».
\end{itemize}
\end{attention}

\courssub{Idiomatic expressions}

Les \cle{expressions idiomatiques} sont des images figées : leur sens global ne se déduit pas des mots. Bonne nouvelle : beaucoup ont un équivalent français presque identique !

\begin{definition}[To make ends meet]
\eng{To make ends meet} : \emph{to have just enough money to pay for the things you need} — joindre les deux bouts. \eng{With two children at school, it is not easy to make ends meet.} « Avec deux enfants à l'école, il n'est pas facile de joindre les deux bouts. »
\end{definition}

\begin{exemplebox}[Idioms in sentences — observez et traduisez]
\begin{itemize}
\item \eng{The new phone costs a fortune.} « Le nouveau téléphone coûte une fortune. »
\item \eng{Many families struggle to make ends meet.} « Beaucoup de familles peinent à joindre les deux bouts. »
\item \eng{I am broke — can you lend me 2,000 francs?} « Je suis fauché — peux-tu me prêter 2 000 francs ? »
\item \eng{They live within their means and never borrow.} « Ils vivent selon leurs moyens et n'empruntent jamais. »
\item \eng{We had to tighten our belts after father lost his job.} « Nous avons dû nous serrer la ceinture après que père a perdu son emploi. »
\item \eng{Money doesn't grow on trees, my son!} « L'argent ne pousse pas sur les arbres, mon fils ! »
\item \eng{Our shop is in the red this year.} « Notre boutique est à découvert cette année. »
\item \eng{After the cocoa season, the cooperative is in the black.} « Après la saison du cacao, la coopérative est créditrice. »
\item \eng{She keeps a little nest egg for her old age.} « Elle garde un petit pécule pour ses vieux jours. »
\end{itemize}
\end{exemplebox}

\begin{popfigure}
\begin{center}
\begin{tikzpicture}
\draw[fill=popblueL, draw=popblue, thick] (0,0) rectangle (7.2,1);
\node[font=\bfseries] at (3.6,0.5) {IDIOM};
\draw[fill=poporangeL, draw=poporange, thick] (7.4,0) rectangle (15,1);
\node[font=\bfseries] at (11.2,0.5) {MEANING};
\draw[fill=white, draw=popline] (0,-1) rectangle (7.2,0);
\node[align=center] at (3.6,-0.5) {\eng{to be broke}};
\draw[fill=white, draw=popline] (7.4,-1) rectangle (15,0);
\node[align=center] at (11.2,-0.5) {être fauché, sans argent};
\draw[fill=popcream, draw=popline] (0,-2) rectangle (7.2,-1);
\node[align=center] at (3.6,-1.5) {\eng{to cost a fortune}};
\draw[fill=popcream, draw=popline] (7.4,-2) rectangle (15,-1);
\node[align=center] at (11.2,-1.5) {coûter une fortune, être très cher};
\draw[fill=white, draw=popline] (0,-3) rectangle (7.2,-2);
\node[align=center] at (3.6,-2.5) {\eng{to make ends meet}};
\draw[fill=white, draw=popline] (7.4,-3) rectangle (15,-2);
\node[align=center] at (11.2,-2.5) {joindre les deux bouts};
\draw[fill=popcream, draw=popline] (0,-4) rectangle (7.2,-3);
\node[align=center] at (3.6,-3.5) {\eng{to tighten one's belt}};
\draw[fill=popcream, draw=popline] (7.4,-4) rectangle (15,-3);
\node[align=center] at (11.2,-3.5) {se serrer la ceinture, se restreindre};
\end{tikzpicture}
\end{center}
\legende{Quatre idiomes essentiels et leur équivalent français.}
\end{popfigure}

\begin{aretenir}[Les huit idiomes à connaître par cœur]
\begin{itemize}
\item \eng{to be broke} = être fauché ; \eng{to cost a fortune} = coûter une fortune ; \eng{to make ends meet} = joindre les deux bouts.
\item \eng{to live within one's means} = vivre selon ses moyens ; \eng{to tighten one's belt} = se serrer la ceinture ; \eng{money doesn't grow on trees} = l'argent ne pousse pas sur les arbres.
\item \eng{to be in the red} = être à découvert ; \eng{to be in the black} = être créditeur ; \eng{a nest egg} = un pécule.
\end{itemize}
Attention à la grammaire : le possessif s'adapte — \eng{I live within \textbf{my} means}, \eng{she tightens \textbf{her} belt}.
\end{aretenir}

\begin{savaistu}[Why “in the red”?]
L'expression \eng{to be in the red} vient de la couleur rouge utilisée autrefois par les comptables pour noter les dettes et les pertes dans les registres ; les sommes positives étaient écrites à l'encre noire — d'où \eng{to be in the black}, « être créditeur ». Quant à \eng{nest egg} (« œuf de nid »), les fermiers anglais laissaient un œuf dans le nid pour encourager la poule à pondre : l'image est devenue celle d'une petite réserve d'argent qui « fait grossir » l'épargne.
\end{savaistu}

\courssub{Word families, pronunciation and spelling}

\begin{propriete}[Familles de mots : dérivation nom / verbe / adjectif]
Connaître les dérivés permet de varier le style et de comprendre des textes plus complexes. Le suffixe \eng{-er} / \eng{-or} forme souvent l'agent, \eng{-ing} le nom d'action, \eng{-ings} (pluriel) les économies.
\end{propriete}

\begin{center}
\begin{tabular}{|l|l|l|l|}
\hline
\rowcolor{poppurpleL} \textbf{Verbe} & \textbf{Nom (action/agent)} & \textbf{Adjectif / Participe} & \textbf{Traduction principale} \\
\hline
\eng{earn} & \eng{earnings (n. pl.), earner} & \eng{earned} & gagner (salaire) \\
\hline
\eng{spend} & \eng{spending, spender} & \eng{spent} & dépenser \\
\hline
\eng{save} & \eng{savings (n. pl.), saver} & \eng{saved} & économiser, sauver \\
\hline
\eng{borrow} & \eng{borrower, loan (n.)} & \eng{borrowed} & emprunter \\
\hline
\eng{lend} & \eng{lender, loan (n.)} & \eng{lent} & prêter \\
\hline
\eng{owe} & \eng{debt (n.), debtor} & \eng{owing} & devoir (argent) \\
\hline
\eng{budget (v.)} & \eng{budget (n.)} & \eng{budgeted} & budgétiser \\
\hline
\eng{invest} & \eng{investment, investor} & \eng{invested} & investir \\
\hline
\end{tabular}
\end{center}

\begin{definition}[Prononciation — repères pour francophones]
L'anglais ne se prononce pas comme il s'écrit. Voici les pièges fréquents de ce thème (notation approximative entre guillemets français) :
\end{definition}

\begin{center}
\begin{tabular}{|l|l|l|}
\hline
\rowcolor{poppinkL} \textbf{Mot} & \textbf{Prononciation (approx.)} & \textbf{Observation} \\
\hline
\eng{budget} & « beu-djite » /bʌdʒɪt/ & \eng{dge} = [dʒ] comme dans \eng{bridge} ; \eng{u} = [ʌ] (son « e » ouvert) \\
\hline
\eng{debt} & « dette » /dɛt/ & \textbf{b muet} — ne prononcez jamais le \eng{b} ! \\
\hline
\eng{savings} & « sé-vi-ngz » /ˈseɪvɪŋz/ & \eng{ai} = [eɪ] (é long) ; \eng{ng} = [ŋ] (nasal vélaire) ; \eng{s} final = [z] \\
\hline
\eng{owe} & « o » /əʊ/ & se prononce exactement comme \eng{oh} — \eng{ow} = [əʊ] \\
\hline
\eng{borrow} & « bo-ro » /ˈbɒrəʊ/ & double \eng{r} ; \eng{ow} final = [əʊ] \\
\hline
\eng{lent} & « lente » /lɛnt/ & voyelle courte [ɛ] — irrégulier (pas \eng{lend-ed}) \\
\hline
\eng{fortune} & « for-tchioune » /ˈfɔːtʃuːn/ & \eng{tu} = [tʃuː] (son « tchou ») \\
\hline
\end{tabular}
\end{center}

\begin{attention}[Orthographe : pièges classiques]
\begin{itemize}
\item \eng{budget} : un seul \eng{t} à la fin (pas \eng{budgett}).
\item \eng{debt} : \eng{b} muet étymologique (latin \emph{debitum}) — écrivez-le mais ne le prononcez pas.
\item \eng{savings} : toujours au pluriel pour « économies » ; le singulier \eng{saving} = « action d'économiser » ou « réduction ».
\item \eng{owe} / \eng{own} : \eng{owe} (devoir de l'argent) se prononce « o » ; \eng{own} (posséder) se prononce « one » /əʊn/ avec [n] final.
\item \eng{lend} / \eng{borrow} : \eng{lent} (passé de \eng{lend}) s'écrit avec \eng{e} ; \eng{borrowed} (passé de \eng{borrow}) est régulier.
\end{itemize}
\end{attention}

\begin{methode}[Comment mémoriser le vocabulaire financier]
\begin{enumerate}
\item Apprenez chaque collocation et idiome dans une \cle{phrase complète liée à votre vie}, jamais dans une liste isolée : \eng{I am broke after buying my school books.}
\item Associez l'expression anglaise à son équivalent français quand il existe (\eng{tighten one's belt} = se serrer la ceinture) : l'image commune facilite la mémoire.
\item Créez un \cle{tableau de familles de mots} (verbe / nom / adjectif) pour chaque notion clé : \eng{save — savings — saved}.
\item Entraînez-vous à la prononciation à voix haute : \eng{budget} (« beu-djite »), \eng{debt} (« dette »), \eng{owe} (« o »).
\item Réutilisez le vocabulaire à l'oral dans la semaine (dialogue avec un camarade, présentation courte).
\item Relisez vos phrases et tableaux après quelques jours : la répétition espacée fixe le vocabulaire durablement.
\end{enumerate}
\end{methode}

\begin{exoresolu}[Complétez avec la bonne collocation ou le bon idiome]
Énoncé — Complete the sentences with: \eng{pays off}, \eng{in the red}, \eng{make ends meet}, \eng{stick to}, \eng{nest egg}, \eng{owes}.
\begin{enumerate}
\item My uncle \dots\ 10,000 francs to the shopkeeper.
\item With her small salary, she can hardly \dots .
\item If you \dots\ your budget, you will not be broke.
\item The business is \dots\ : it loses money every month.
\item He finally \dots\ his debt after six months.
\item Grandmother keeps a \dots\ for emergencies.
\end{enumerate}
\tcblower
Solution détaillée :
\begin{enumerate}
\item \eng{owes} — « devoir de l'argent à quelqu'un » : \eng{My uncle owes 10,000 francs to the shopkeeper.}
\item \eng{make ends meet} — joindre les deux bouts : avec un petit salaire, elle peine à joindre les deux bouts.
\item \eng{stick to} — respecter : \eng{stick to a budget} est la collocation exacte.
\item \eng{in the red} — à découvert : l'entreprise perd de l'argent chaque mois.
\item \eng{pays off} — rembourser totalement : \eng{pay off a debt}.
\item \eng{nest egg} — un pécule, une réserve d'argent pour les urgences.
\end{enumerate}
\end{exoresolu}

\begin{attention}[Ne traduisez pas les idiomes mot à mot]
\begin{itemize}
\item \eng{to be broke} ne veut pas dire « être cassé » : c'est « être fauché ». Ne dites pas \eng{I have no money in my pocket} si vous voulez utiliser l'idiome : dites simplement \eng{I am broke}.
\item \eng{money doesn't grow on trees} se dit toujours ainsi, à la forme négative et au présent — ne changez pas la structure.
\item \eng{a fortune} dans \eng{it costs a fortune} est exagéré volontairement : ne le remplacez pas par \eng{it costs money}, qui n'a pas le même sens.
\item \eng{to be in the black} ne signifie pas « être dans le noir » mais « être créditeur / avoir des bénéfices ».
\end{itemize}
\end{attention}

\begin{exercice}[À vous de jouer]
\begin{enumerate}
\item Traduisez en anglais : « Après l'achat de ses fournitures scolaires, je suis fauché. » / « Ils ont ouvert un compte d'épargne et mettent de l'argent de côté chaque mois. » / « Ce sac coûte une fortune ! »
\item Trouvez la collocation : \eng{to \dots\ up a debt} ; \eng{to \dots\ into one's savings} ; \eng{to \dots\ money on clothes}.
\item Complétez le tableau des familles de mots :
\begin{center}
\begin{tabular}{|l|l|l|}
\hline
\rowcolor{popblueL} Verbe & Nom (action/résultat) & Adjectif / Participe passé \\
\hline
\eng{invest} & \dots & \dots \\
\hline
\dots & \eng{loan} & \dots \\
\hline
\eng{borrow} & \dots & \eng{borrowed} \\
\hline
\end{tabular}
\end{center}
\item Rédigez trois phrases personnelles avec \eng{tighten my belt}, \eng{live within my means} et \eng{make ends meet}.
\item Prononcez à voix haute : \eng{budget}, \eng{debt}, \eng{savings}, \eng{owe}, \eng{fortune}. Notez la prononciation approximative en lettres françaises pour chaque mot.
\end{enumerate}
\end{exercice}

\begin{corrige}[Exercice « À vous de jouer »]
\begin{enumerate}
\item \eng{After buying my school supplies, I am broke.} / \eng{They opened a savings account and put money aside every month.} / \eng{This bag costs a fortune!}
\item \eng{run up a debt} ; \eng{dip into one's savings} ; \eng{spend money on clothes}.
\item \eng{invest} — \eng{investment / investor} — \eng{invested} ; \eng{lend} — \eng{loan} — \eng{lent} ; \eng{borrow} — \eng{borrower / loan} — \eng{borrowed}.
\item Exemples : \eng{I had to tighten my belt after the holidays.} / \eng{I try to live within my means.} / \eng{With my pocket money, I can just make ends meet.}
\item \eng{budget} : « beu-djite » ; \eng{debt} : « dette » ; \eng{savings} : « sé-vi-ngz » ; \eng{owe} : « o » ; \eng{fortune} : « for-tchioune ».
\end{enumerate}
\end{corrige}

\coursec{False friends and traps for francophones}

Dans la section précédente, nous avons appris à combiner les mots de la finance en collocations naturelles (\eng{to save money}, \eng{to run into debt}, \eng{to make ends meet}). Mais il reste un danger bien plus sournois : les mots qui \emph{ressemblent} au français et qui, justement pour cela, nous trahissent. Ce sont les \cle{faux amis} (\eng{false friends}). Un élève de Première qui écrit \eng{I have no economies} croit dire « je n'ai pas d'économies » ; en réalité, sa phrase ne veut rien dire pour un anglophone. Cette section vous donne les armes pour repérer et neutraliser ces pièges, un par un.

\courssub{Observons d'abord : un dialogue piégé}

\begin{textebox}[At the bank in Buea — a dialogue full of traps]
\textbf{Clerk:} “Good morning, Madam. How can I help you?”

\textbf{Mrs Efosi:} “Good morning. I would like to open a savings account. You see, I have no savings at all, and I want to change that.”

\textbf{Clerk:} “A wise decision. Are you currently employed?”

\textbf{Mrs Efosi:} “Yes, I teach at a secondary school. Actually, I also run a small shop in the evening, so I earn a little extra money.”

\textbf{Clerk:} “Excellent. Last month I attended a workshop on budgeting, and I can share some useful tips with you. Our bank also assists young entrepreneurs with small loans.”

\textbf{Mrs Efosi:} “That is interesting. My brother asked the bank for a loan last year, but he is still paying back his debt.”

\textbf{Clerk:} “I see. Will you be depositing cash today? Notes or coins?”

\textbf{Mrs Efosi:} “Notes, mostly. Oh, and if I pay in foreign currency, will you give me my change in CFA francs?”

\textbf{Clerk:} “Of course, Madam. Here is your change and your receipt.”
\end{textebox}

\begin{description}
\item[\eng{savings account} (n.)] : compte d'épargne
\item[\eng{currently} (adv.)] : actuellement
\item[\eng{actually} (adv.)] : en fait
\item[\eng{to attend} (v.)] : assister à
\item[\eng{to assist} (v.)] : aider
\item[\eng{loan} (n.)] : prêt
\item[\eng{debt} (n.)] : dette
\item[\eng{change} (n.)] : monnaie rendue
\item[\eng{currency} (n.)] : devise
\end{description}

\textbf{Questions de compréhension :}
\begin{enumerate}
\item Why does Mrs Efosi want to open a savings account?
\item What is the difference between what the clerk \emph{attended} and what the bank \emph{assists}?
\item Who asked for a loan, and what is he still paying back?
\end{enumerate}

\courssub{Le grand tableau des faux amis de la finance}

\begin{center}
\begin{tabularx}{\linewidth}{|X|X|X|}
\hline
\rowcolor{popblueL} \textbf{Faux ami} & \textbf{Sens anglais réel} & \textbf{Traduction du français visé} \\
\hline
economy & l'économie (d'un pays) & « économies » (argent épargné) = \textbf{savings} \\
\hline
actually & en fait, en réalité & « actuellement » = \textbf{currently} \\
\hline
to assist & aider & « assister à » = \textbf{to attend} \\
\hline
to demand & exiger & « demander » = \textbf{to ask / to ask for} \\
\hline
currency & devise, monnaie (système) & « monnaie » (rendue) = \textbf{change} \\
\hline
a coin & une pièce de monnaie & « coin » (angle) = \textbf{a corner} \\
\hline
\end{tabularx}
\end{center}

\begin{popfigure}
\begin{tikzpicture}[
  box/.style={draw=popink, fill=popblueL, rounded corners=2pt, align=center, font=\small, text width=3.1cm, minimum height=1.1cm},
  real/.style={draw=popgreen, fill=popgreenL, rounded corners=2pt, align=center, font=\small, text width=3.4cm, minimum height=1.1cm},
  fr/.style={draw=poporange, fill=poporangeL, rounded corners=2pt, align=center, font=\small, text width=3.6cm, minimum height=1.1cm},
  arr/.style={-{Stealth[length=2.5mm]}, thick, popdark}
]
\node[font=\bfseries\small, popink] at (0,0) {FAUX AMI};
\node[font=\bfseries\small, popgreen!60!black] at (4.4,0) {SENS ANGLAIS RÉEL};
\node[font=\bfseries\small, poporange!70!black] at (9.2,0) {FRANÇAIS VISÉ};

\node[box] (a1) at (0,-1.3) {economy};
\node[real] (a2) at (4.4,-1.3) {l'économie du pays};
\node[fr] (a3) at (9.2,-1.3) {« économies » = savings};
\draw[arr] (a1) -- (a2); \draw[arr, poporange] (a3) -- (a1);

\node[box] (b1) at (0,-2.8) {actually};
\node[real] (b2) at (4.4,-2.8) {en fait};
\node[fr] (b3) at (9.2,-2.8) {« actuellement » = currently};
\draw[arr] (b1) -- (b2); \draw[arr, poporange] (b3) -- (b1);

\node[box] (c1) at (0,-4.3) {to assist};
\node[real] (c2) at (4.4,-4.3) {aider};
\node[fr] (c3) at (9.2,-4.3) {« assister à » = to attend};
\draw[arr] (c1) -- (c2); \draw[arr, poporange] (c3) -- (c1);

\node[box] (d1) at (0,-5.8) {to demand};
\node[real] (d2) at (4.4,-5.8) {exiger};
\node[fr] (d3) at (9.2,-5.8) {« demander » = to ask for};
\draw[arr] (d1) -- (d2); \draw[arr, poporange] (d3) -- (d1);
\end{tikzpicture}
\legende{Quatre faux amis à connaître : le mot anglais ne dit jamais ce que le français suggère.}
\end{popfigure}

\courssub{Étude détaillée des pièges, un par un}

\textbf{1. \eng{economy} et \eng{savings} : le piège numéro un.}

\begin{propriete}[Economy ou savings ?]
\eng{Economy} (nom, prononcé « i-ko-no-mi ») désigne \cle{l'économie d'un pays}, c'est-à-dire son système économique. L'argent qu'une personne met de côté se dit \cle{savings} (toujours au pluriel). Le verbe \eng{to economise} (orthographe britannique) existe mais il est rare et soutenu ; on préfère \eng{to save money} ou \eng{to cut down on expenses} (réduire ses dépenses).
\end{propriete}

\begin{exemplebox}[Comparez]
\eng{The Cameroonian economy depends heavily on cocoa and oil.} « L'économie camerounaise dépend fortement du cacao et du pétrole. »

\eng{She has no savings, so she cannot pay the school fees.} « Elle n'a pas d'économies, donc elle ne peut pas payer les frais de scolarité. »

\eng{We are saving money to buy a motorbike.} « Nous économisons de l'argent pour acheter une moto. »
\end{exemplebox}

\begin{attention}[« I have no economies » : un calque du français]
Ne dites jamais \eng{I have no economies} : c'est un calque mot à mot du français « je n'ai pas d'économies ». Dites \textbf{I have no savings}. Réservez \eng{economy} au système économique d'un pays : \eng{The Cameroonian economy is growing} (« L'économie camerounaise est en croissance »). De même, « faire des économies » se dit \eng{to save money}, et non « to make economies ».
\end{attention}

\textbf{2. \eng{actually} et \eng{currently} : le temps contre la réalité.}

\begin{propriete}[Actually ne veut pas dire « actuellement »]
\eng{Actually} signifie \cle{en fait}, \cle{en réalité} : il sert à corriger ou préciser une idée. Pour dire « actuellement », utilisez \cle{currently}, \cle{at the moment} ou \cle{now}.
\end{propriete}

\begin{exemplebox}[Comparez]
\eng{I currently have 3,000 francs left in my account.} « J'ai actuellement 3 000 francs sur mon compte. »

\eng{You think he is rich? Actually, he is deeply in debt.} « Tu crois qu'il est riche ? En fait, il est lourdement endetté. »

\eng{She is currently working in Douala.} « Elle travaille actuellement à Douala. »
\end{exemplebox}

\textbf{3. \eng{to assist} et \eng{to attend} : aider ou être présent ?}

\begin{propriete}[Assist = aider ; attend = assister à]
\eng{To assist} quelqu'un signifie \cle{l'aider}. Pour dire « assister à » une réunion, un cours, une cérémonie, utilisez \cle{to attend} (suivi directement du nom, sans préposition : \eng{to attend a meeting}).
\end{propriete}

\begin{exemplebox}[Comparez]
\eng{The bank assists young entrepreneurs with small loans.} « La banque aide les jeunes entrepreneurs avec de petits prêts. »

\eng{I attended a workshop on budgeting last week.} « J'ai assisté à un atelier sur la budgétisation la semaine dernière. »
\end{exemplebox}

\begin{attention}[« I assist the meeting » est faux]
Ne dites jamais \eng{I assist the meeting} pour « j'assiste à la réunion » : un anglophone comprendra que vous \emph{aidez} la réunion, ce qui n'a pas de sens. Dites \textbf{I attend the meeting}. Et attention à la préposition : on dit \eng{to attend a meeting} (sans « at »), mais \eng{to assist somebody} (sans « to »).
\end{attention}

\textbf{4. \eng{to demand} et \eng{to ask (for)} : exiger ou demander ?}

\begin{propriete}[Demand est beaucoup plus fort que « demander »]
\eng{To demand} signifie \cle{exiger}, avec une idée d'autorité ou de colère. Le « demander » poli et ordinaire du français se traduit par \cle{to ask} (quelqu'un) ou \cle{to ask for} (quelque chose).
\end{propriete}

\begin{exemplebox}[Comparez]
\eng{He asked the bank for a loan.} « Il a demandé un prêt à la banque. »

\eng{The workers demanded higher wages.} « Les ouvriers ont exigé des salaires plus élevés. »

\eng{She asked me for some change.} « Elle m'a demandé de la monnaie. »
\end{exemplebox}

\textbf{5. \eng{currency} et \eng{change} : la devise et la monnaie rendue.}

\begin{propriete}[Deux « monnaies » différentes]
\eng{Currency} désigne la \cle{devise} d'un pays (le franc CFA, le dollar, la livre). La monnaie que le commerçant vous rend se dit \cle{change} (invariable).
\end{propriete}

\begin{exemplebox}[Comparez]
\eng{The currency of Cameroon is the CFA franc.} « La monnaie (devise) du Cameroun est le franc CFA. »

\eng{The shopkeeper gave me my change.} « Le commerçant m'a rendu la monnaie. »

\eng{Here is your change: five hundred francs.} « Voici votre monnaie : cinq cents francs. »
\end{exemplebox}

\textbf{6. \eng{coin} et \eng{corner} : le piège orthographique.}

\begin{attention}[Une seule lettre change tout]
\eng{A coin} (une pièce de monnaie) et \eng{a corner} (un coin, un angle) se ressemblent dangereusement à l'écrit. Comparez : \eng{I found a coin on the floor.} « J'ai trouvé une pièce par terre. » / \eng{The bank is at the corner of the street.} « La banque est au coin de la rue. » Relisez toujours ce mot dans vos rédactions.
\end{attention}

\textbf{7. \eng{debt} : le « b » muet.}

\begin{savaistu}[Pourquoi « debt » s'écrit-il avec un b ?]
Le mot \eng{debt} vient du latin \emph{debitum}. Au Moyen Âge, les érudits anglais ont rajouté le « b » pour rappeler cette origine latine, mais on ne l'a jamais prononcé ! On dit donc simplement « dette ». Même phénomène dans \eng{doubt} (« doute ») et \eng{subtle} (« subtil ») : le « b » est muet. À l'oral, \eng{He is in debt} se prononce « hi iz in dette ».
\end{savaistu}

\courssub{Méthode : comment éviter les calques dans vos rédactions}

\begin{methode}[La règle du « mot trop beau pour être vrai »]
\begin{enumerate}
\item En rédigeant, repérez tout mot anglais qui \emph{ressemble fortement} au français (economy, assist, demand, actually, currency...).
\item Méfiez-vous : plus la ressemblance est grande, plus le risque de faux ami est élevé.
\item Vérifiez le mot dans un \textbf{dictionnaire anglais-anglais} (ou bilingue fiable) : lisez la définition anglaise et l'exemple, pas seulement la traduction.
\item Si le mot ne convient pas, cherchez l'équivalent correct (savings, attend, ask for, currently, change).
\item Relisez votre phrase à voix haute : un calque « sonne » souvent bizarre, même avant de comprendre pourquoi.
\end{enumerate}
\end{methode}

\courssub{Exercice résolu et entraînement}

\begin{exoresolu}[Corrige les erreurs]
Énoncé — Chaque phrase contient un faux ami ou un calque du français. Corrige-la.
\begin{enumerate}
\item I have no economies, so I cannot buy the books.
\item I assist the English club every Friday.
\item He demanded a loan from the microfinance bank. (sens : il a simplement demandé)
\item Actually, I live in Bamenda. (sens : actuellement)
\item The trader gave me my currency.
\end{enumerate}
\tcblower
Solution détaillée —
\begin{enumerate}
\item \eng{I have no \textbf{savings}, so I cannot buy the books.} — « économies » (argent épargné) = \eng{savings} ; \eng{economy} désigne l'économie d'un pays.
\item \eng{I \textbf{attend} the English club every Friday.} — « assister à » = \eng{to attend} ; \eng{to assist} signifie « aider ».
\item \eng{He \textbf{asked for} a loan from the microfinance bank.} — « demander » = \eng{to ask for} ; \eng{to demand} = exiger, ce qui serait agressif ici.
\item \eng{\textbf{Currently}, I live in Bamenda.} — « actuellement » = \eng{currently} ; \eng{actually} = en fait.
\item \eng{The trader gave me my \textbf{change}.} — la monnaie rendue = \eng{change} ; \eng{currency} = la devise d'un pays.
\end{enumerate}
\end{exoresolu}

\begin{exercice}[À toi de jouer]
Traduis en anglais en évitant les faux amis :
\begin{enumerate}
\item J'assiste à un cours sur la gestion de l'argent tous les mardis.
\item En fait, elle a beaucoup d'économies à la banque.
\item Il a demandé un prêt à son oncle.
\item Actuellement, l'économie du pays se porte bien.
\item Gardez la monnaie !
\end{enumerate}
\end{exercice}

\begin{corrige}[Exercice « À toi de jouer »]
\begin{enumerate}
\item \eng{I attend a class on money management every Tuesday.}
\item \eng{Actually, she has a lot of savings in the bank.}
\item \eng{He asked his uncle for a loan.}
\item \eng{Currently, the country's economy is doing well.}
\item \eng{Keep the change!}
\end{enumerate}
\end{corrige}

\begin{aretenir}[L'essentiel de la section]
\begin{itemize}
\item « économies » (argent) = \eng{savings} ; \eng{economy} = l'économie d'un pays.
\item \eng{actually} = en fait ; « actuellement » = \eng{currently}.
\item \eng{to assist} = aider ; « assister à » = \eng{to attend}.
\item \eng{to demand} = exiger ; « demander » = \eng{to ask / to ask for}.
\item \eng{currency} = devise ; la monnaie rendue = \eng{change}.
\item \eng{coin} (pièce) et \eng{corner} (coin) : attention à l'orthographe.
\item \eng{debt} : le « b » est muet, on prononce « dette ».
\end{itemize}
\end{aretenir}

\coursec{Word families, pronunciation and spelling}

Après avoir identifié les \textbf{faux amis} et les pièges de traduction dans la section précédente, nous devons maintenant consolider notre maîtrise lexicale en comprenant comment les mots se construisent (morphologie), comment ils se prononcent (phonologie) et comment ils s'écrivent (orthographe). Une bonne connaissance des \textbf{familles de mots} permet de multiplier son vocabulaire à partir d'une racine unique, tandis que la vigilance sur les \textbf{lettres muettes} et les \textbf{doublements consonantiques} évite les erreurs classiques des francophones.

\courssub{Word families: derivation and conversion}

En anglais, la dérivation (ajout de préfixes/suffixes) et la conversion (changement de catégorie sans changement de forme) sont les deux moteurs principaux de l'enrichissement lexical. Observons les cinq familles clés de cette leçon.

\begin{center}
\begin{tabularx}{\linewidth}{|X|X|X|X|}
\hline
\rowcolor{popblueL} \textbf{Base verbale / Nominale} & \textbf{Noms (Nouns)} & \textbf{Adjectifs / Participles} & \textbf{Agent / Personne} \\
\hline
\textbf{SAVE} & savings (pl.), saving (acte) & saving (adj.: a saving grace) & a saver \\
\hline
\textbf{SPEND} & spending (n. indénombrable) & spending (adj.: spending power) & a spender / a big spender \\
\hline
\textbf{LEND} (lent, lent) & a loan (nom distinct), lending & lending (adj.: lending rate) & a lender \\
\hline
\textbf{BORROW} & borrowing (n.) & borrowed (adj.: borrowed time) & a borrower \\
\hline
\textbf{PAY} (paid, paid) & payment, pay (salaire) & payable (adj.), paid (adj.) & a payee (bénéficiaire) / a payer (payeur) \\
\hline
\textbf{ECONOMY} (n.) & economics (science), economy (système) & economic / economical & an economist \\
\hline
\end{tabularx}
\end{center}

\medskip
\noindent \textbf{Observation :} Remarquez la conversion pour \eng{save}, \eng{spend}, \eng{pay} (verbe $\to$ nom sans suffixe ou avec \textit{-ing}) et la dérivation suffixale pour \eng{economy} (\textit{-ic, -ical, -ist, -ics}). Pour \eng{lend/borrow}, le nom d'action principal est souvent un mot distinct (\eng{loan}) ou le gérondif (\eng{lending/borrowing}).

\begin{propriete}[Economic vs Economical : distinction cruciale]
Ces deux adjectifs dérivent de \eng{economy} mais leurs emplois sont radicalement différents~:
\begin{itemize}
    \item \textbf{Economic} (pron. /ˌiːkəˈnɒmɪk/ ou /ˌekəˈnɒmɪk/) : se rapporte à \textbf{l'économie en tant que système, science ou secteur} (macro-économie, finances publiques).
    \eng{The government announced new economic reforms.} « Le gouvernement a annoncé de nouvelles réformes économiques. »
    \item \textbf{Economical} (pron. /ˌiːkəˈnɒmɪkl/ ou /ˌekəˈnɒmɪkl/) : signifie \textbf{« économe, qui évite le gaspillage, rentable »} (micro-niveau : personne, machine, gestion).
    \eng{This car is very economical; it does 20 km per litre.} « Cette voiture est très économe ; elle fait 20 km au litre. »
    \eng{She is an economical housewife.} « C'est une ménagère économe. »
\end{itemize}
\emph{Astuce mnémotechnique :} \textbf{Economical} finit par \textbf{-al} comme \textbf{personn\textbf{al}} ou \textbf{loc\textbf{al}} $\to$ cela concerne un comportement \textbf{personnel} ou \textbf{local}.
\end{propriete}

\begin{exemplebox}[Familles \eng{SAVE} et \eng{SPEND} en contexte]
\begin{itemize}
    \item \eng{She is a careful \textbf{saver}; she puts 20\% of her salary into her \textbf{savings account} every month.}
    « Elle est une épargnante prudente ; elle met 20\% de son salaire sur son compte d'épargne chaque mois. »
    \item \eng{His \textbf{spending} habits have changed since he lost his job; he is no longer a \textbf{big spender}.}
    « Ses habitudes de dépenses ont changé depuis qu'il a perdu son emploi ; il n'est plus un grand dépensier. »
    \item \eng{The bank is the biggest \textbf{lender} in the region; it checks the \textbf{borrower}'s credit history carefully.}
    « La banque est le plus grand prêteur de la région ; elle vérifie soigneusement l'historique de crédit de l'emprunteur. »
    \item \eng{The \textbf{payment} is \textbf{payable} by the 30th of the month. The \textbf{payee} is the landlord.}
    « Le paiement est exigible pour le 30 du mois. Le bénéficiaire est le propriétaire. »
\end{itemize}
\end{exemplebox}

\begin{attention}[Savings : toujours un \textbf{'s'} final pour l'argent mis de côté]
En français, « l'épargne » est un singulier. En anglais, \textbf{\eng{savings}} (l'argent épargné) est un \textbf{nom pluriel uniquement} (pluralia tantum). Il s'accorde avec un verbe au pluriel.
\begin{itemize}
    \item \eng{My \textbf{savings are} in a safe place.} (Correct)
    \item \eng{My \textbf{savings is} in a safe place.} (Incorrect $\leftarrow$ calque du français « mon épargne est »)
\end{itemize}
\emph{Note :} \eng{A saving} (singulier) existe mais signifie « une économie réalisée sur un achat » (ex. \eng{We made a saving of 5000 francs}).
\end{attention}

\begin{savaistu}[Economics : une science au singulier]
Le mot \textbf{\eng{economics}} (la science économique, l'économie politique) se termine par \textbf{-s} mais prend \textbf{un verbe au singulier}, comme \eng{mathematics}, \eng{physics}, \eng{news}.
\eng{Economics \textbf{is} a difficult subject for many students.}
« L'économie est une matière difficile pour beaucoup d'étudiants. »
De même : \eng{The economics of the project \textbf{are} complex} (rare, sens « aspects financiers »), mais standard : \eng{Economics \textbf{is}...}
\end{savaistu}

\begin{popfigure}
\legende{Arbre lexical de la famille \textit{ECONOMY}}
\begin{tikzpicture}[
    mindmap,
    grow cyclic,
    every node/.style={concept, execute at begin node=\hskip0pt, align=center},
    root concept/.append style={concept color=popblue, minimum size=2.5cm, text=white, font=\large\bfseries},
    level 1 concept/.append style={level distance=4cm, sibling angle=90, minimum size=1.8cm, font=\bfseries},
    level 2 concept/.append style={level distance=3cm, sibling angle=45},
    concept color=popblue!70,
    text=white
]
\node[root concept, align=center]{economy\\ (noun)}
    child[concept color=popgreen] { node[align=center]{economic\\ (adj.)\\ \textit{économique (système)}}
        child { node[align=center]{economically\\ (adv.)} }
    }
    child[concept color=poporange] { node[align=center]{economical\\ (adj.)\\ \textit{économe (gaspillage)}}
        child { node[align=center]{economically\\ (adv.)} }
    }
    child[concept color=poppurple] { node[align=center]{economist\\ (noun)\\ \textit{économiste}} }
    child[concept color=poppink] { node[align=center]{economics\\ (noun, sing.)\\ \textit{science éco.}} };
\end{tikzpicture}
\end{popfigure}

\courssub{Prononciation : les pièges des lettres muettes}

L'orthographe anglaise conserve souvent des traces étymologiques (latin, français ancien) qui ne se prononcent plus. Trois mots de notre thème illustrent parfaitement ce phénomène : \eng{debt}, \eng{receipt}, \eng{mortgage}. Ajoutons \eng{budget} (accentuation) et \eng{afford} (voyelle réduite).

\begin{methode}[Le trio des lettres muettes : \textit{b, p, t}]
Pour ne plus jamais prononcer ces lettres parasites, suivez cette routine en 3 étapes~:
\begin{enumerate}
    \item \textbf{Visualisez} le mot en surlignant la lettre muette mentalement : \eng{de\textcolor{red}{b}t}, \eng{reci\textcolor{red}{p}t}, \eng{mor\textcolor{red}{t}gage}.
    \item \textbf{Associez} un mot-référence français qui a le même son final : \eng{debt} $\to$ « dette » (pas de b) ; \eng{receipt} $\to$ « recette » (pas de p) ; \eng{mortgage} $\to$ « mort-gage » (le t de mort ne s'entend pas devant g).
    \item \textbf{Répétez à voix haute} 5 fois la phrase : \eng{“I paid my debt for the mortgage and kept the receipt.”} en veillant à ne \textbf{jamais} articuler le \textit{b}, le \textit{p}, le \textit{t}.
\end{enumerate}
\end{methode}

\begin{center}
\begin{tabularx}{\linewidth}{|X|X|X|}
\hline
\rowcolor{popblueL} \textbf{Mot} & \textbf{Prononciation approximative (français)} & \textbf{Traduction / Note} \\
\hline
\textbf{debt} & « dette » & \textit{b} muet (etym. latin \textit{debitum}). \eng{National debt} (dette nationale). \\
\hline
\textbf{receipt} & « ri-sit » & \textit{p} muet. \eng{Keep the receipt} (Gardez le reçu). Différent de \eng{recipe} (« ré-ci-pe », recette cuisine). \\
\hline
\textbf{mortgage} & « mor-gage » & \textit{t} muet. Litt. « gage mort » (vieux français). \eng{To pay off the mortgage} (rembourser l'hypothèque). \\
\hline
\textbf{budget} & « beu-djè-te » & Accent sur 1ère syllabe \textbf{BUD}-get. Pas « bu-djé ». \\
\hline
\textbf{afford} & « a-ford » & \textit{a} = schwa /ə/. \eng{Can you afford it?} (Peux-tu te le permettre ?). \\
\hline
\end{tabularx}
\end{center}

\begin{exoresolu}[Entraînement à la prononciation et discrimination auditive]
\textbf{Énoncé :} Écoutez (ou lisez à voix haute) les paires de phrases suivantes. Identifiez le mot-cible en gras et cochez la bonne prononciation approximative parmi les propositions.

1. \eng{He has a huge \textbf{debt}.}
   a) « dé-bette » \quad b) « dette » \quad c) « dé-bé-té »

2. \eng{Please show me the \textbf{receipt}.}
   a) « re-seipte » \quad b) « ri-sit » \quad c) « re-ki-pe-te »

3. \eng{They signed the \textbf{mortgage} papers.}
   a) « mor-ti-gage » \quad b) « mor-gage » \quad c) « mor-ti-ge »

4. \eng{We must discuss the \textbf{budget}.}
   a) « bu-djé » \quad b) « beu-djè-te » \quad c) « beu-djé »

\tcblower

\textbf{Solution détaillée :}
\begin{enumerate}
    \item \textbf{b) « dette »}. Le \textit{b} est muet. Erreur fréquente : prononcer le \textit{b} par analogie avec \eng{debit} (débit) où il se prononce.
    \item \textbf{b) « ri-sit »}. Le \textit{p} est muet. Attention à ne pas confondre avec \eng{recipe} (recette de cuisine) où le \textit{p} se prononce /p/ $\to$ « ré-ci-pe ».
    \item \textbf{b) « mor-gage »}. Le \textit{t} est muet. L'étymologie \textit{mort-gage} (gage mort) aide à mémoriser l'orthographe mais pas la prononciation du \textit{t}.
    \item \textbf{b) « beu-djè-te »}. Accent tonique sur la première syllabe \textbf{BUD}. Le \textit{u} se prononce /ʌ/ (son « beu » ouvert), le \textit{dge} = /dʒ/ (son « dj »).
\end{enumerate}
\end{exoresolu}

\courssub{Orthographe : doublement des consonnes}

Les francophones ont tendance à simplifier les doubles consonnes (une seule consonne en français souvent) ou à en ajouter là où il n'y en a pas. Quatre mots de cette leçon sont des pièges classiques.

\begin{center}
\begin{tabularx}{\linewidth}{|X|X|X|}
\hline
\rowcolor{popblueL} \textbf{Mot} & \textbf{Règle / Explication} & \textbf{Piège francophone} \\
\hline
\textbf{saving} & Verbe \eng{save} (termine par \textit{e} muet) + \textit{-ing} $\to$ on garde le \textit{v} simple. \textbf{Un seul \textit{v}}. & Écrire *\eng{savving} (double \textit{v} par analogie avec \eng{running}, mais \eng{run} finit par consonne unique voyelle brève). \\
\hline
\textbf{spending} & Verbe \eng{spend} (termine par deux consonnes \textit{nd}) + \textit{-ing} $\to$ \textbf{pas de doublement} (règle du CVC non applicable). & Écrire *\eng{spendding} (double \textit{d} inexistant). \\
\hline
\textbf{borrow} & Racine historique \textit{borrow} $\to$ \textbf{deux \textit{r}} médians. & Écrire *\eng{borow} (un seul \textit{r}, calque sur la prononciation « bor-o »). \\
\hline
\textbf{afford} & Préfixe \textit{ad-} (devenu \textit{af-} par assimilation) + racine \textit{ford} $\to$ \textbf{deux \textit{f}}. & Écrire *\eng{aford} (un seul \textit{f}). Penser à \eng{affirm}, \eng{affect}, \eng{aggressive} (assimilation préfixe \textit{ad-}). \\
\hline
\end{tabularx}
\end{center}

\courssub{Mise en contexte : texte et activités}

Le texte suivant réemploie l'essentiel du vocabulaire, des familles de mots et des points de prononciation/orthographe vus ci-dessus.

\begin{textebox}[Managing a Student Budget in Yaoundé]
\textbf{Armand}, a second-year student at the University of Yaoundé I, sat on a cheap bench near the faculty canteen, staring at his phone. His \textbf{savings}, carefully built up during the long vacation selling mangoes at \textbf{Mokolo} market, were melting fast. “I cannot \textbf{afford} this new textbook,” he whispered. The \textbf{receipt} for the photocopies he had just bought showed 15,000 FCFA — a huge \textbf{spending} for a single week.

He recalled his father’s advice: “Be an \textbf{economical} student, not a \textbf{big spender}. Avoid \textbf{debt} at all costs.” But the \textbf{economic} situation in the country was tough; prices had risen. Armand needed a \textbf{loan} for his \textbf{tuition fees}. He went to the microfinance agency near \textbf{Carrefour Warda}. The \textbf{lender} examined his file. “You are a reliable \textbf{borrower},” the agent said, “but the \textbf{interest rate} is high. The \textbf{payment} is \textbf{payable} monthly.”

Armand signed the \textbf{mortgage} agreement for his father’s land title used as collateral. He felt the weight of the \textbf{obligation}. Later, buying a cheap but \textbf{economical} second-hand laptop from a friend, he realised: “Being a good \textbf{saver} is not just about keeping money, it is about making \textbf{economical} choices every day.” He put the \textbf{receipt} in his wallet, promising himself to track all his \textbf{spending} in a notebook.
\end{textebox}

\medskip
\noindent \textbf{Glossaire (mots difficiles) :}
\begin{itemize}
    \item \textbf{melting fast} (loc.) : fondre rapidement (disparaître).
    \item \textbf{textbook} (n.) : manuel scolaire.
    \item \textbf{tuition fees} (n. pl.) : frais de scolarité / inscription.
    \item \textbf{microfinance agency} (n.) : institution de microfinance.
    \item \textbf{collateral} (n.) : garantie, sûreté (bien mis en gage).
    \item \textbf{interest rate} (n.) : taux d'intérêt.
    \item \textbf{track} (v.) : suivre, tracer, noter (dépenses).
    \item \textbf{budget} (adj.) : bon marché, peu coûteux (ex. \eng{a budget airline}, \eng{a budget bench}).
\end{itemize}

\medskip
\noindent \textbf{Questions de compréhension et de réemploi lexical :}
\begin{enumerate}
    \item Trouvez dans le texte le nom au pluriel qui désigne l'argent mis de côté et soulignez le verbe qui l'accompagne (indice : accord).
    \item Identifiez les deux adjectifs en \textit{-ical} dérivés de \eng{economy}. Lequel qualifie l'étudiant ? Lequel qualifie la situation du pays ?
    \item Relevez les trois mots contenant une lettre muette (\textit{b, p, t}) et écrivez leur prononciation approximative.
    \item Quel mot de la famille \eng{PAY} est utilisé comme adjectif signifiant « exigible » ?
    \item Reformulez la phrase \eng{“Being a good saver is not just about keeping money...”} en utilisant le verbe \eng{to save} et le nom \eng{savings}.
\end{enumerate}

\begin{experience}[Budget Role-Play : \eng{The Loan Interview}]
\textbf{Objectif :} Réemployer le vocabulaire (lender, borrower, loan, interest rate, payable, collateral, afford, economical) à l'oral.
\textbf{Déroulement (par paires) :}
\begin{enumerate}
    \item \textbf{Rôle A (Lender / Agent de microfinance) :} Vous recevez un étudiant. Posez des questions sur ses \eng{savings}, son \eng{spending} mensuel, sa capacité à \eng{afford} les mensualités. Demandez une \eng{collateral}. Annoncez l'\eng{interest rate} et la date où le \eng{payment} est \eng{payable}.
    \item \textbf{Rôle B (Borrower / Étudiant) :} Expliquez pourquoi vous avez besoin du \eng{loan} (frais de scolarité, \eng{economical} laptop). Montrez que vous êtes un \eng{careful saver} et pas un \eng{big spender}. Négociez le taux.
    \item \textbf{Échange :} Jouez la scène 3 min. Inversez les rôles.
    \item \textbf{Restitution :} Un binôme joue devant la classe. La classe note le vocabulaire utilisé correctement.
\end{enumerate}
\end{experience}

\begin{exercice}[Consolidation : Familles, Prononciation, Orthographe]
\textbf{Exercice 1 — Dérivation :} Complétez le tableau en trouvant la forme manquante (nom, verbe, adjectif, personne). La première ligne est un exemple.
\begin{center}
\begin{tabular}{|l|l|l|l|l|}
\hline
\rowcolor{popblueL} \textbf{Verbe} & \textbf{Nom (chose/action)} & \textbf{Adjectif} & \textbf{Personne (Agent)} & \textbf{Traduction FR (Verbe)} \\
\hline
save & savings & saving & saver & épargner \\
\hline
spend & \dots & \dots & spender & \dots \\
\hline
lend & loan / lending & \dots & \dots & prêter \\
\hline
borrow & \dots & borrowed & \dots & emprunter \\
\hline
pay & payment & payable & payee / payer & payer \\
\hline
\dots & economy & economic / economical & economist & \dots \\
\hline
\end{tabular}
\end{center}

\textbf{Exercice 2 — Prononciation :} Écrivez la prononciation approximative (en lettres françaises) pour ces mots. Entourez la lettre muette dans le mot anglais.
\begin{enumerate}
    \item debt \hfill 2. receipt \hfill 3. mortgage \hfill 4. budget \hfill 5. afford
\end{enumerate}

\textbf{Exercice 3 — Orthographe :} Choisissez l'orthographe correcte.
\begin{enumerate}
    \item She is a good (saver / savver / saiver).
    \item His (spending / spendding) habits are dangerous.
    \item The bank is the main (lender / lender / lendor).
    \item We cannot (afford / aford / affordd) a new car.
    \item Did you keep the (receipt / receit / recipe)?
\end{enumerate}

\textbf{Exercice 4 — Choix lexical (Economic vs Economical) :} Complétez avec \eng{economic}, \eng{economical} ou \eng{economics}.
\begin{enumerate}
    \item The country faces a severe \dots crisis.
    \item Buying in bulk is often more \dots
    \item The \dots outlook for next year is uncertain.
    \item This washing machine is very \dots (low water consumption).
    \item She teaches \dots at university.
\end{enumerate}
\end{exercice}

\begin{corrige}[Exercice 1 à 4]
\textbf{Exercice 1 — Tableau complété :}
\begin{center}
\begin{tabular}{|l|l|l|l|l|}
\hline
\rowcolor{popblueL} \textbf{Verbe} & \textbf{Nom} & \textbf{Adjectif} & \textbf{Personne} & \textbf{Traduction} \\
\hline
spend & spending & spending & spender / big spender & dépenser \\
\hline
lend & loan / lending & lending & lender & prêter \\
\hline
borrow & borrowing & borrowed & borrower & emprunter \\
\hline
pay & payment & payable / paid & payee / payer & payer \\
\hline
--- & economy & economic / economical & economist & (gérer l'économie) \\
\hline
\end{tabular}
\end{center}
\emph{Note :} Pour \eng{economy}, il n'y a pas de verbe direct courant (\eng{to economise} existe mais signifie « faire des économies », pas « gérer l'économie »).

\textbf{Exercice 2 — Prononciation :}
1. \eng{debt} $\to$ « \textbf{dette} » (\textit{b} muet).
2. \eng{receipt} $\to$ « \textbf{ri-sit} » (\textit{p} muet).
3. \eng{mortgage} $\to$ « \textbf{mor-gage} » (\textit{t} muet).
4. \eng{budget} $\to$ « \textbf{beu-djè-te} » (accent sur \textbf{BUD}).
5. \eng{afford} $\to$ « \textbf{a-ford} » (\textit{a} = schwa, double \textit{ff} = /f/).

\textbf{Exercice 3 — Orthographe correcte :}
1. \textbf{saver} (un seul \textit{v}, comme \eng{save}).
2. \textbf{spending} (un seul \textit{n}, \eng{spend} finit par \textit{nd}).
3. \textbf{lender} (un seul \textit{n}, \eng{lend} $\to$ \eng{lender}).
4. \textbf{afford} (deux \textit{f}, assimilation \textit{ad-} + \textit{ford}).
5. \textbf{receipt} (\textit{p} muet, \textit{ei} = /i:/). \emph{Piège :} \eng{recipe} = recette cuisine.

\textbf{Exercice 4 — Economic vs Economical vs Economics :}
1. \textbf{economic} crisis (crise du système économique).
2. more \textbf{economical} (économe, évite gaspillage).
3. \textbf{economic} outlook (perspectives du système/secteur).
4. very \textbf{economical} (machine peu consommatrice).
5. \textbf{economics} (nom de la discipline/science).
\end{corrige}

\coursec{Using the vocabulary in context}

Dans la section précédente, nous avons travaillé sur les familles de mots (\eng{save} / \eng{savings} / \eng{saver}), la prononciation et l'orthographe (\eng{receive}, \eng{afford}, \eng{borrow} vs \eng{lend}). Connaître un mot isolément ne suffit pas : un vocabulaire n'est réellement acquis que lorsqu'on sait l'\cle{employer en contexte} — dans un texte, un dialogue, une rédaction. C'est exactement ce que l'examinateur attend de vous au BEPC comme au Probatoire : réutiliser le lexique de façon naturelle et correcte. Cette dernière section est donc un atelier pratique : lire, écouter, compléter, reformuler, rédiger et jouer.

\courssub{Reading: Aminatou's monthly budget}

Lisons d'abord un court texte authentique. Votre tâche : \textbf{soulignez mentalement tous les mots de la leçon} que vous reconnaissez.

\begin{textebox}[Aminatou's monthly budget]
Aminatou is a Form Five student at a secondary school in Bamenda. Every month, she receives 10,000 francs from her parents: that is her income. She also earns a little extra money by plaiting her friends' hair on Saturdays, which brings in about 2,000 francs more.

Aminatou is careful with money. At the beginning of each month, she draws up a simple budget. She spends about 5,000 francs on food and transport — the bendskin to school costs her 100 francs a day. She also buys phone data for 2,000 francs, because she needs the internet for her homework. She never spends money on things she cannot afford: when her friends invite her to eat out, she sometimes says no.

Every month, she puts 2,000 francs aside in a small wooden savings box at home. “I am saving up for a second-hand laptop,” she explains. The last 1,000 francs is her emergency fund, in case she falls ill or needs to print documents urgently.

Last week, her classmate Divine asked her: “Could you lend me 1,500 francs? I will pay you back on Monday.” Aminatou agreed, but she wrote the amount in her notebook. “I like helping my friends,” she says, “but I always stick to my budget. That way, I am never broke before the end of the month.”
\end{textebox}

\textbf{Glossaire} (mots difficiles) :

\begin{center}
\begin{tabularx}{\linewidth}{|l|l|X|}
\hline
\rowcolor{popblueL} Mot anglais & Nature & Traduction française \\
\hline
income & nom & revenu \\
\hline
to earn & verbe & gagner (de l'argent) \\
\hline
to draw up a budget & expression & établir un budget \\
\hline
to afford & verbe & avoir les moyens de \\
\hline
to put aside & verbe & mettre de côté \\
\hline
savings box & nom & tirelire, caisse d'épargne \\
\hline
to save up for & verbe & économiser en vue de \\
\hline
emergency fund & nom & fonds d'urgence \\
\hline
to pay back & verbe & rembourser \\
\hline
broke & adjectif (familier) & fauché, sans le sou \\
\hline
\end{tabularx}
\end{center}

\textbf{Comprehension questions} (répondez en phrases complètes) :

\begin{enumerate}
\item What is Aminatou's total monthly income?
\item How does she earn extra money?
\item How much does she spend on phone data?
\item What is she saving up for?
\item Why did she write Divine's loan in her notebook?
\item What advice does she give at the end of the text?
\end{enumerate}

\begin{corrige}[Questions de compréhension]
\begin{enumerate}
\item Her total income is about 12,000 francs a month (10,000 from her parents plus about 2,000 from plaiting hair).
\item She earns extra money by plaiting her friends' hair on Saturdays.
\item She spends 2,000 francs on phone data.
\item She is saving up for a second-hand laptop.
\item She wrote it down so as not to forget the amount and to make sure Divine pays her back.
\item She advises to always stick to a budget in order never to be broke before the end of the month.
\end{enumerate}
\end{corrige}

Pour visualiser la logique du texte, voici l'organigramme du budget d'Aminatou — un modèle que vous pourrez reproduire dans votre propre rédaction.

\begin{popfigure}
\begin{center}
\begin{tikzpicture}[
node distance=0.9cm,
box/.style={draw=popblue, thick, fill=popblueL, rounded corners, align=center, minimum width=4.2cm, minimum height=1cm, font=\small},
arrow/.style={-{Stealth[length=3mm]}, thick, popdark}
]
\node[box, align=center] (income){\textbf{INCOME}\\ pocket money + small jobs\\ (12,000 F)};
\node[box, below=of income, fill=poporangeL, draw=poporange, align=center] (expenses){\textbf{EXPENSES}\\ food, transport, data\\ (7,000 F)};
\node[box, below left=0.9cm and -1.2cm of expenses, fill=popgreenL, draw=popgreen, align=center] (savings){\textbf{SAVINGS}\\ about 10--20\%\\ (2,000 F)};
\node[box, below right=0.9cm and -1.2cm of expenses, fill=popgoldL, draw=popgold, align=center] (emergency){\textbf{EMERGENCY FUND}\\ for urgent needs\\ (1,000 F)};
\draw[arrow] (income) -- (expenses);
\draw[arrow] (expenses.south) -- (savings.north);
\draw[arrow] (expenses.south) -- (emergency.north);
\end{tikzpicture}
\end{center}
\legende{Organigramme : « My monthly budget » — du revenu (income) aux dépenses (expenses), à l'épargne (savings) et au fonds d'urgence (emergency fund).}
\end{popfigure}

\courssub{Listening and speaking: a model dialogue}

Voici maintenant un dialogue modèle entre deux amis qui comparent leurs façons de gérer leur argent de poche. Lisez-le à voix haute à deux, en soignant la prononciation de \eng{budget} (« beu-djète »), \eng{afford} (« a-ford ») et \eng{borrow} (« bo-ro »).

\begin{textebox}[Pocket money habits]
\textbf{Ngong:} Hi, Raissa! You look worried. What's the matter?

\textbf{Raissa:} It's only the middle of the month and I'm already broke. I spent all my pocket money on data and snacks.

\textbf{Ngong:} You should draw up a budget. I receive 8,000 francs a month, and I always put 1,000 aside before I spend anything.

\textbf{Raissa:} That's wise. The problem is, I can't resist buying things. Last week I even borrowed 2,000 francs from my cousin.

\textbf{Ngong:} Did you? When will you pay him back?

\textbf{Raissa:} On Friday, I hope. By the way... could you lend me 500 francs for lunch today? I'll pay you back tomorrow, I promise.

\textbf{Ngong:} Sure, but pay me back on time this time! And my advice is: stick to a budget, and never borrow money for things you don't really need.

\textbf{Raissa:} You're right. From next month, I'll save first and spend what is left — not the other way round!
\end{textebox}

\begin{attention}[La politesse dans les demandes d'argent]
En anglais, on ne demande jamais de l'argent avec l'impératif direct. \eng{Give me 500 francs!} est très impoli (c'est un ordre). Employez les formules atténuées :
\begin{itemize}
\item \eng{Could you lend me 500 francs, please?} — Peux-tu me prêter 500 francs, s'il te plaît ?
\item \eng{Can I borrow 500 francs?} — Je peux emprunter 500 francs ?
\end{itemize}
Et pour répondre : \eng{Sure, but pay me back on time.} (D'accord, mais rembourse-moi à temps.) ou \eng{Sorry, I'm broke myself.} (Désolé, je suis fauché moi aussi.) Notez bien : \eng{lend} = prêter (je donne), \eng{borrow} = emprunter (je reçois). Ne les inversez pas !
\end{attention}

\begin{exemplebox}[Deux phrases modèles à retenir par cœur]
\eng{Could you lend me 1,000 francs? I will pay you back on Monday.} = « Peux-tu me prêter 1 000 francs ? Je te rembourserai lundi. »

\eng{I can't afford to eat out every day.} = « Je ne peux pas me permettre de manger au restaurant tous les jours. »

Remarquez la construction : \eng{afford} est suivi de \textbf{to + infinitif} (\eng{afford to buy}, \eng{afford to pay}), et \eng{pay back} est un verbe à particule \textbf{séparable} : \eng{I will pay you back} (le pronom se place entre \eng{pay} et \eng{back}).
\end{exemplebox}

\courssub{Practice 1: fill in the blanks with the right verb}

\begin{exoresolu}[Choose the right verb: earn, spend, save, borrow, lend, afford]
\textbf{Énoncé.} Complétez chaque phrase avec le verbe qui convient, à la forme correcte.

\begin{enumerate}
\item My uncle \ldots\ldots\ldots 150,000 francs a month as a taxi driver in Douala.
\item I can't \ldots\ldots\ldots a new phone; I only have 5,000 francs.
\item She \ldots\ldots\ldots 3,000 francs from her sister yesterday, and she will pay her back on Saturday.
\item If you \ldots\ldots\ldots 500 francs every week, you will have 26,000 francs at the end of the year.
\item Don't \ldots\ldots\ldots all your pocket money on sweets!
\item Can you \ldots\ldots\ldots me your pen, please?
\end{enumerate}

\tcblower
\textbf{Solution détaillée.}
\begin{enumerate}
\item \textbf{earns} — gagner un salaire = \eng{to earn} ; 3e personne du singulier, d'où le \emph{-s}.
\item \textbf{afford} — « avoir les moyens de » ; après \eng{can't}, l'infinitif sans \eng{to}.
\item \textbf{borrowed} — elle a \emph{reçu} l'argent (emprunter) ; \eng{yesterday} impose le prétérit : \eng{borrowed}.
\item \textbf{save} — mettre de côté régulièrement = \eng{to save} ; dans la subordonnée avec \eng{if}, présent simple.
\item \textbf{spend} — structure \eng{spend money on something} ; après l'impératif négatif \eng{don't}, infinitif.
\item \textbf{lend} — ici c'est moi qui \emph{donne} temporairement : prêter = \eng{lend}. (Comparez : \eng{Can I borrow your pen?} — même situation, autre point de vue.)
\end{enumerate}
\end{exoresolu}

\begin{center}
\begin{tabularx}{\linewidth}{|X|X|X|}
\hline
\rowcolor{popblueL} Verbe & Construction & Exemple \\
\hline
earn & earn money / earn a salary & She earns 2,000 F a week. \\
\hline
spend & spend money \textbf{on} sth & He spends 500 F on data. \\
\hline
save & save money / save up \textbf{for} sth & I'm saving up for a bike. \\
\hline
borrow & borrow sth \textbf{from} sb & I borrowed money from Ali. \\
\hline
lend & lend sth \textbf{to} sb / lend sb sth & Lend me 1,000 F, please. \\
\hline
afford & (can) afford \textbf{to} + infinitive & We can't afford to eat out. \\
\hline
pay back & pay sb back / pay back sth & I'll pay you back on Monday. \\
\hline
\end{tabularx}
\end{center}

\courssub{Practice 2: reformulate — kill the calque!}

Les phrases suivantes sont des \textbf{calques du français}. Corrigez-les en anglais naturel.

\begin{exoresolu}[From French-style English to correct English]
\textbf{Énoncé.} Chaque phrase contient une erreur typique de francophone. Réécrivez-la correctement.

\begin{enumerate}
\item I don't have means to buy this bag. (« Je n'ai pas les moyens. »)
\item He wins a lot of money in his new job. (« Il gagne beaucoup. »)
\item She borrowed me 2,000 francs. (« Elle m'a prêté 2 000 francs. »)
\item I am sparing money for a computer. (« J'épargne pour un ordinateur. »)
\item He spent two hours for his homework... and all his money for games. (« dépenser pour »)
\end{enumerate}

\tcblower
\textbf{Solution détaillée.}
\begin{enumerate}
\item \eng{I can't afford to buy this bag.} — « avoir les moyens de » se dit \eng{afford}, presque toujours avec \eng{can/can't}.
\item \eng{He earns a lot of money in his new job.} — \eng{win} = gagner (un match, un prix, à la loterie) ; pour un salaire, c'est \eng{earn}.
\item \eng{She lent me 2,000 francs.} — « prêter » = \eng{lend} (prétérit \eng{lent}) ; \eng{borrow} = emprunter. On peut dire aussi : \eng{She lent 2,000 francs to me.}
\item \eng{I am saving money for a computer.} — \eng{to spare} signifie « se passer de / avoir de reste » ; « épargner » = \eng{save}.
\item \eng{He spent two hours on his homework... and all his money on games.} — la préposition après \eng{spend} est \eng{on} (ou \eng{-ing} : \eng{he spent all his money buying games}), jamais \eng{for}.
\end{enumerate}
\end{exoresolu}

\begin{attention}[Récapitulatif des pièges]
\begin{itemize}
\item \eng{earn} (salaire) $\neq$ \eng{win} (match, prix, loterie).
\item \eng{borrow from} (recevoir) $\neq$ \eng{lend to} (donner).
\item \eng{spend money \textbf{on} something}, jamais \eng{for}.
\item \eng{afford} s'emploie avec \eng{can} : \eng{I can't afford it}, pas \eng{I don't afford it}.
\item \eng{money} est \textbf{indénombrable} : \eng{money is}, \eng{a lot of money} — jamais \eng{monies} ni \eng{a money} au sens courant.
\end{itemize}
\end{attention}

\courssub{Guided writing: “How I manage my pocket money”}

\begin{methode}[Rédiger 5 à 6 lignes en quatre étapes]
\begin{enumerate}
\item \textbf{Dites d'où vient l'argent} (income) : \eng{I receive ... francs a month from my parents. I also earn ... by ...}
\item \textbf{Listez vos dépenses} (expenses) avec \eng{spend ... on ...} : \eng{I spend ... on food, transport and data.}
\item \textbf{Dites ce que vous épargnez} (savings) : \eng{I put ... aside every month. I am saving up for ...}
\item \textbf{Concluez par une habitude ou un conseil} : \eng{I always stick to my budget, so I am never broke.} ou \eng{My advice is: save first, spend later.}
\end{enumerate}
Consigne : utilisez \textbf{au moins six mots de la leçon} et soulignez-les.
\end{methode}

\begin{exemplebox}[Modèle de rédaction (6 lignes)]
\eng{Every month, I receive 6,000 francs of pocket money from my mother. I also earn a little money by selling groundnuts after school. I spend most of my income on transport and phone data, but I never buy things I cannot afford. Every week, I put 500 francs aside in my savings box, because I am saving up for a school bag. When my friends ask to borrow money, I only lend what I can. I always stick to my budget, so I am never broke at the end of the month.}

(Six mots-clefs de la leçon y figurent : \eng{pocket money}, \eng{earn}, \eng{spend}, \eng{afford}, \eng{put aside}, \eng{saving up}, \eng{borrow}, \eng{lend}, \eng{stick to a budget}, \eng{broke}.)
\end{exemplebox}

\begin{exercice}[À vous d'écrire]
Rédigez votre propre paragraphe \eng{How I manage my pocket money} (5 à 6 lignes) en suivant la méthode en quatre étapes. Soulignez au moins six mots de la leçon. Échangez ensuite votre copie avec votre voisin : il doit retrouver et entourer vos six mots de vocabulaire.
\end{exercice}

\courssub{Role play: negotiating a small loan}

\begin{experience}[Borrowing and paying back]
Travaillez à deux. \textbf{Élève A} : vous avez besoin de 2,000 francs pour acheter un manuel scolaire ; négociez un petit prêt avec votre ami et proposez une date de remboursement. \textbf{Élève B} : vous acceptez ou refusez, mais posez des conditions (montant, date, promesse).

Expressions à utiliser obligatoirement :
\begin{itemize}
\item \eng{Could you lend me ...?} / \eng{Can I borrow ...?}
\item \eng{I will pay you back on ...} / \eng{I promise to pay you back before Friday.}
\item \eng{Sure, but pay me back on time.} / \eng{Sorry, I'm broke myself.}
\item \eng{How much do you need?} / \eng{When exactly will you pay me back?}
\end{itemize}
Puis inversez les rôles. Le reste de la classe écoute et note chaque emploi correct de \eng{lend}, \eng{borrow} et \eng{pay back}.
\end{experience}

\begin{savaistu}[The 50/30/20 rule]
Les conseillers financiers anglo-saxons recommandent souvent la règle « 50/30/20 » pour gérer un revenu : \textbf{50\%} pour les besoins essentiels (\eng{needs} : nourriture, transport, logement), \textbf{30\%} pour les envies (\eng{wants} : loisirs, sorties, data), et \textbf{20\%} pour l'épargne (\eng{savings}). Aminatou, dans notre texte, épargne environ 25\% de ses revenus — elle est donc plus disciplinée que la règle ! Gardez cette règle en tête : c'est un excellent argument à mobiliser au Probatoire et au Baccalauréat dans un essai ou un débat sur l'argent (\eng{Young people should learn to manage money at school. Do you agree?}).
\end{savaistu}

\begin{aretenir}[L'essentiel de la section]
\begin{itemize}
\item Un vocabulaire n'est acquis que s'il est \cle{réemployé en contexte} : lecture, dialogue, rédaction, jeu de rôle.
\item Les constructions à maîtriser : \eng{spend on}, \eng{borrow from}, \eng{lend to}, \eng{afford to + infinitif}, \eng{pay sb back}, \eng{save up for}, \eng{stick to a budget}.
\item Politesse : \eng{Could you lend me...?} et non \eng{Give me...}.
\item Pour rédiger : income $\rightarrow$ expenses $\rightarrow$ savings $\rightarrow$ conseil final.
\end{itemize}
\end{aretenir}

\newpage

\coursec{Activité d'intégration}

\coursec{Lesson 22 — Vocabulary: Words and expressions related to managing personal finances}

\courssub{Objectifs de l'activité}
Cette activité de synthèse vise à mobiliser l'ensemble du vocabulaire et des structures étudiés dans la leçon 22 pour réaliser une tâche de communication complexe et authentique. À l'issue de cette séquence, l'élève doit être capable de :
\begin{itemize}
    \item \textbf{Mobiliser} le lexique des finances personnelles (\cle{income}, \cle{expenses}, \cle{savings}, \cle{loan}, \cle{budget}, \cle{afford}, \cle{put aside}, \cle{make ends meet}, etc.) dans une production écrite cohérente et une présentation orale fluide.
    \item \textbf{Structurer} un budget mensuel simple en distinguant clairement les postes de recettes (\eng{income sources}) et de dépenses (\eng{expense categories}), en justifiant les choix par des collocations appropriées (\eng{cover the rent}, \eng{pay off a debt}, \cle{save up for}).
    \item \textbf{Interagir} en situation de \eng{role-play} : présenter un plan financier, répondre à des questions précises (\eng{How much do you allocate to...?}, \eng{Why can't you afford...?}) et formuler des conseils en utilisant des modaux de conseil (\eng{should}, \eng{ought to}, \eng{had better}) et le vocabulaire cible.
    \item \textbf{Identifier et corriger} les erreurs fréquentes des francophones (faux amis \eng{benefit} vs \eng{profit}, \eng{expenses} vs \eng{charges}, confusion \eng{borrow}/\eng{lend}, \eng{earn}/\eng{win}) lors de la phase de relecture et d'évaluation par les pairs.
\end{itemize}

\courssub{Situation}
\begin{textebox}[Appel à projets : « My Personal Budget Challenge » — Youth Finance Club, Lycée de Nkolbisson, Yaoundé]
\textbf{My Personal Budget Challenge}

\textit{Are you ready to take control of your money? The Youth Finance Club invites all Form 5 students to participate in our annual budgeting competition.}

\textbf{The Task:}
Work in groups of three. Your team is a \textbf{Financial Planning Unit}. You must create a realistic \textbf{monthly budget} for one of the three profiles below. The monthly \textbf{net income} is fixed. You must decide how to allocate every franc.

\textbf{Profile A — Étienne, 17, Form 5 Student (Yaoundé)}
\begin{itemize}
    \item \textbf{Net Monthly Income:} 45,000 FCFA (pocket money 25,000 + weekend tutoring 20,000).
    \item \textbf{Context:} Lives with parents (no rent). Needs: transport (bendskin/bus), phone credit/data, school supplies, snacks, savings for university entrance fees (concours), leisure.
\end{itemize}

\textbf{Profile B — Maman Amina, 38, Food Vendor (Mokolo Market, Douala)}
\begin{itemize}
    \item \textbf{Net Monthly Income:} 180,000 FCFA (daily profit after food costs).
    \item \textbf{Context:} Single mother, two children in primary school. Rent: 30,000 FCFA. Market stall fee: 15,000 FCFA. Needs: food, school fees, transport, medical emergencies, \textbf{tontine} contribution, savings to expand business.
\end{itemize}

\textbf{Profile C — Ngwa, 24, Junior Accountant (Bamenda)}
\begin{itemize}
    \item \textbf{Net Monthly Income:} 250,000 FCFA (entry-level salary).
    \item \textbf{Context:} Lives alone in a one-room apartment. Rent: 50,000 FCFA. Electricity/Water: 15,000 FCFA. Transport: 20,000 FCFA. Has a student \textbf{loan} repayment: 30,000 FCFA/month. Wants to \textbf{save up for} a professional certification (ACCA) and an emergency fund.
\end{itemize}

\textbf{Deliverables:}
\begin{enumerate}
    \item A written \textbf{Budget Sheet} (6–8 lines, table format) using \textbf{at least 8 key words/expressions} from Lesson 22 (underlined in your copy).
    \item A 2-minute \textbf{Oral Pitch} to the class: “Here is [Name]'s budget. His/Her income is... Expenses include... He/She manages to... but struggles to...”
    \item A \textbf{Q\&A Session}: Another group acts as \textbf{Financial Advisors} and asks two challenging questions.
\end{enumerate}

\textbf{Assessment Criteria:} Accuracy of vocabulary (spelling, collocations), realism of figures, fluency and pronunciation, relevance of advisors' questions, teamwork.
\textit{Deadline: End of the double period. Best budget wins a “Smart Saver” certificate!}
\end{textebox}

\courssub{Consignes}
\begin{enumerate}
    \item \textbf{Formation des équipes et choix du profil (5 min).} Constituez des groupes de 3 élèves. Choisissez l'un des trois profils (Étienne, Amina ou Ngwa). Lisez attentivement le contexte et les contraintes financières. Discutez en anglais des priorités : \eng{“What are the fixed costs?”}, \eng{“How much can we allocate to savings?”}, \eng{“Is there room for unexpected expenses?”}.
    \item \textbf{Rédaction du Budget Sheet (15 min).} Sur une feuille A4 (ou au tableau), tracez un tableau à deux colonnes : \eng{Income} / \eng{Expenses}. Remplissez-le avec des montants réalistes (total expenses = income). Rédigez 6 à 8 phrases commentaires \textbf{sous le tableau} pour expliquer les choix. \textbf{Contrainte lexicale :} Vous devez intégrer \textbf{au minimum 8 unités lexicales} distinctes de la leçon 22 (verbes clés, collocations, idiomes, mots de la famille \eng{finance}). Soulignez-les. Exemples attendus : \eng{Etienne's income covers his transport but he cannot afford a new phone. He tries to put aside 5,000 FCFA monthly to save up for university fees. He must stick to his budget to avoid borrowing from friends.}
    \item \textbf{Préparation de l'Oral Pitch (10 min).} Répartissez les rôles : \textit{Speaker 1} présente les revenus et charges fixes ; \textit{Speaker 2} détaille les dépenses variables et l'épargne ; \textit{Speaker 3} conclut sur les difficultés (\eng{make ends meet}, \eng{tighten one's belt}) et les objectifs. Entraînez-vous à prononcer correctement les termes techniques (\eng{budget} [« beu-djit »], \eng{expenses} [« èk-spen-siz »], \eng{loan} [« loun »], \eng{afford} [« a-ford »]). Vérifiez l'accord sujet-verbe (\eng{Expenses are...}, \eng{Income is...}).
    \item \textbf{Présentation et Phase « Financial Advisors » (20 min).} Chaque groupe présente son budget (2 min max). Un autre groupe désigné joue le rôle de \textit{Financial Advisors}. Ils doivent poser \textbf{deux questions précises} en réemployant le vocabulaire. Exemples : \eng{“Why did you allocate only 5,000 FCFA to savings? Isn't it risky if an emergency occurs?”}, \eng{“She spends 20,000 on data. Could she cut down on that to pay off her loan faster?”}, \eng{“He wants to splash out on sneakers. Shouldn't he prioritize his emergency fund?”}. Le groupe présentateur répond sur le champ.
    \item \textbf{Vote et Débriefing (10 min).} La classe vote pour le budget le plus \eng{realistic}, le mieux \eng{argued} et le plus \eng{lexically rich}. L'enseignant note au tableau les meilleures formulations entendues et corrige les erreurs récurrentes (faux amis, prépositions, prononciation).
\end{enumerate}

\courssub{Ressources à mobiliser}
Voici un rappel synthétique des outils linguistiques de la leçon 22 à garder sous les yeux pendant l'activité.

\begin{center}
\begin{tabularx}{\linewidth}{|X|X|X|}
\hline
\rowcolor{popblueL}
\textbf{Verbes clés (Formes \& Emploi)} & \textbf{Collocations \& Expressions idiomatiques} & \textbf{Faux amis \& Pièges (Francophones)} \\
\hline
\eng{earn} (earned, earned) — gagner (salaire) & \eng{earn a living / a salary} & \eng{earn} $\neq$ \eng{win} (gagner un jeu/prix) \\
\eng{spend} (spent, spent) — dépenser & \eng{spend money on sth / doing sth} & \eng{spend time} (passer du temps) — pas *\eng{pass time} \\
\eng{save} (saved, saved) — économiser & \eng{save up for sth}, \eng{put money aside}, \eng{set aside} & \eng{save} $\neq$ \eng{spare} (épargner qqn / avoir en rab) \\
\eng{borrow} (borrowed) — emprunter & \eng{borrow sth from sb}, \eng{take out a loan} & \eng{borrow} (emprunter) vs \eng{lend} (prêter) — \textbf{Ne pas confondre !} \\
\eng{lend} (lent, lent) — prêter & \eng{lend sth to sb}, \eng{lend a hand} & \eng{lend} $\rightarrow$ \eng{loan} (nom) \\
\eng{owe} (owed, owed) — devoir de l'argent & \eng{owe sb money}, \eng{be in debt} & \eng{owe} $\neq$ \eng{own} (posséder) \\
\eng{afford} (can/could afford) — avoir les moyens & \eng{afford to do sth}, \eng{can't afford sth} & Toujours avec modal \eng{can/could/be able to}. Pas *\eng{I afford} \\
\eng{budget} (v.) — budgétiser & \eng{stick to a budget}, \eng{draw up a budget}, \eng{balance the budget} & \eng{budget} (nom) = budget ; \eng{forecast} = prévision \\
\hline
\end{tabularx}
\end{center}

\vspace{0.5cm}
\textbf{Mots de la même famille (Word Families) — Attention à la catégorie grammaticale :}
\begin{itemize}
    \item \eng{Finance} (n.) $\rightarrow$ \eng{Financial} (adj.) $\rightarrow$ \eng{Financially} (adv.) $\rightarrow$ \eng{Financier} (n. person).
    \item \eng{Economy} (n. système) $\rightarrow$ \eng{Economic} (adj. macro) / \eng{Economical} (adj. qui coûte peu) $\rightarrow$ \eng{Economist} (n.).
    \item \eng{Expense} (n.) $\rightarrow$ \eng{Expensive} (adj. cher) $\neq$ \eng{Expansive} (vaste/étendu — \textbf{faux ami fréquent}).
    \item \eng{Invest} (v.) $\rightarrow$ \eng{Investment} (n.) $\rightarrow$ \eng{Investor} (n.).
    \item \eng{Save} (v.) $\rightarrow$ \eng{Savings} (n. pl. économies) $\rightarrow$ \eng{Saver} (n.).
\end{itemize}

\textbf{Structures utiles pour les « Financial Advisors » (Questions \& Conseils) :}
\begin{itemize}
    \item \eng{Why did you decide to...? / What made you allocate... to...?}
    \item \eng{Don't you think you should / ought to / had better...?}
    \item \eng{If I were you, I would... (Second Conditional).}
    \item \eng{Have you considered... + V-ing?}
    \item \eng{It might be wise to... / It is crucial to...}
\end{itemize}

\courssub{Proposition de production et corrigé commenté}
\begin{exoresolu}[Profil A — Étienne (Étudiant à Yaoundé) — Budget Sheet \& Oral Script]
\textbf{Énoncé :} Rédigez le budget mensuel d'Étienne (Revenu net : 45 000 FCFA) et préparez le script de présentation orale (2 min), en intégrant au moins 8 mots de la leçon soulignés.

\tcblower

\textbf{Solution détaillée — Budget Sheet : Étienne, Form 5 Student, Yaoundé}

\begin{center}
\begin{tabular}{|l|r|}
\hline
\textbf{INCOME (Monthly Net)} & \textbf{Amount (FCFA)} \\
\hline
Pocket money (parents) & 25,000 \\
Weekend tutoring (Maths/Physics) & 20,000 \\
\hline
\textbf{TOTAL INCOME} & \textbf{45,000} \\
\hline
\end{tabular}
\hspace{1cm}
\begin{tabular}{|l|rc|}
\hline
\textbf{EXPENSES (Monthly)} & \textbf{Amount (FCFA)} \\
\hline
Transport (Bendskin/Bus to school) & 12,000 \\
Phone Credit \& Data Bundle & 8,000 \\
School Supplies \& Photocopies & 5,000 \\
Snacks / Street Food (Lunch break) & 10,000 \\
Leisure (Video games / Football match) & 3,000 \\
\hline
\textbf{TOTAL FIXED \& VARIABLE EXPENSES} & \textbf{38,000} \\
\hline
\textbf{SAVINGS (Put aside for Univ. Fees)} & \textbf{7,000} \\
\hline
\textbf{BALANCE} & \textbf{0} \\
\hline
\end{tabular}
\end{center}

\textbf{Commentaires écrits (sous le tableau) :}
\begin{enumerate}
    \item Étienne's total monthly \uline{income} amounts to 45,000 FCFA, which \uline{covers} his basic \uline{expenses} but leaves a tight margin.
    \item He \uline{cannot afford} to buy a new smartphone this month, so he \uline{puts aside} 7,000 FCFA to \uline{save up for} his university entrance exam fees (\eng{concours}).
    \item To \uline{stick to his budget}, he avoids \uline{splashing out} on expensive snacks and walks short distances instead of taking a \eng{bendskin}.
    \item He \uline{owes} nothing to anyone and has no \uline{loan} to \uline{repay}, which helps him \uline{make ends meet} without \uline{borrowing} from friends.
    \item His \uline{financial} goal is to build a small \uline{emergency fund} before the long holidays.
\end{enumerate}

\textbf{Script de présentation orale (Speaker 1, 2, 3) :}
\begin{quote}
\textit{Speaker 1:} “Good morning everyone. We are the Financial Planning Unit for \textbf{Étienne}. His monthly \textbf{income} \textbf{stands at} 45,000 francs. It \textbf{comes from} two sources: pocket money from his parents and earnings from weekend tutoring.”

\textit{Speaker 2:} “On the \textbf{expenses} side, fixed costs like \textbf{transport} and \textbf{phone data} take 20,000 francs. Variable costs — food, supplies, leisure — reach 18,000. Total \textbf{expenses} are 38,000. He \textbf{allocates} the remaining 7,000 to \textbf{savings}. He is \textbf{saving up for} his university entrance fees.”

\textit{Speaker 3:} “Étienne \textbf{manages to} \textbf{balance his budget}, but it is tight. He \textbf{can't afford} unexpected costs like a medical bill. He \textbf{sticks to} his plan by \textbf{cutting down on} leisure. He \textbf{doesn't need to borrow} money, which is good. Our advice: he \textbf{should} try to \textbf{increase his income} with more tutoring to build an \textbf{emergency fund} faster. Thank you.”
\end{quote}

\textbf{Questions des « Financial Advisors » (Groupe adverse) \& Réponses :}
\begin{itemize}
    \item \eng{Q1: “You allocate only 7,000 FCFA to savings. Don't you think he should increase this amount to cover potential emergencies?”} \\
    \eng{A1: “We agree, but his income is fixed. He could only save more if he cut transport costs by walking more.”}
    \item \eng{Q2: “He spends 8,000 on data. Is it possible to switch to a cheaper plan to free up cash?”} \\
    \eng{A2: “Yes, he could downgrade his bundle. That would free up 2,000 FCFA for his emergency fund. Good point.”}
\end{itemize}

\textbf{Commentaire des choix de langue (Pour l'enseignant) :}
\begin{itemize}
    \item \textbf{Verbes clés bien employés :} \eng{amounts to} (s'élève à), \textbf{comes from} (provient de), \textbf{covers} (couvre), \textbf{allocate} (allouer), \textbf{manage to} (réussir à), \textbf{stick to} (respecter/s'en tenir à), \textbf{cut down on} (réduire), \textbf{free up} (libérer), \textbf{downgrade} (réduire l'abonnement).
    \item \textbf{Collocations/Idiomes ciblés :} \eng{make ends meet} (boucler son budget), \eng{save up for} (économiser pour), \eng{put aside} (mettre de côté), \eng{splash out on} (se ruiner pour/claquer l'argent dans — utilisé à la forme négative), \eng{tight margin} (marge serrée), \eng{balance the budget} (équilibrer le budget), \eng{emergency fund} (fonds d'urgence).
    \item \textbf{Grammaire :} \eng{Income} (singulier) $\rightarrow$ \eng{amounts/stands/comes}. \eng{Expenses} (pluriel) $\rightarrow$ \eng{are/take/reach}. Modaux \eng{can't afford}, \eng{should increase}, \eng{could cut}, \eng{doesn't need to}. Second conditional \eng{If I were you / he could only save more if...}.
    \item \textbf{Prononciation surveillée :} \eng{Budget} [« beu-djit »], \eng{Expenses} [« èk-spen-siz »], \eng{Data} [« déi-ta »], \eng{Allocate} [« a-lo-keïte »].
    \item \textbf{Évités (Pièges) :} Pas de *\eng{he wins 45,000} (win), pas de *\eng{he spares money} (spare), pas de *\eng{he borrows money from the bank} pour dire « il prend un prêt » (préférer \eng{take out a loan}), pas de *\eng{expensives} (adjectif invariable), pas de *\eng{he affords} (sans modal).
\end{itemize}
\end{exoresolu}

\courssub{Bilan de l'activité}
Cette activité d'intégration a permis de valider les acquis de la leçon 22 à travers une tâche finale complexe (\eng{Task-Based Learning}) ancrée dans le réel camerounais (Mokolo, bendskin, concours, tontine, salaires locaux).

\begin{itemize}
    \item \textbf{Sur le plan lexical :} Les élèves ont dû \textbf{sélectionner}, \textbf{orthographier} et \textbf{combiner} correctement les unités lexicales (verbes, noms, collocations, idiomes). L'obligation d'en utiliser 8 minimum a forcé le recours au dictionnaire mental et à la fiche de vocabulaire. La correction par les pairs (Financial Advisors) a mis en lumière les confusions persistantes (\eng{borrow/lend}, \eng{win/earn}, \eng{expensive/expansive}).
    \item \textbf{Sur le plan grammatical :} La contrainte du budget a imposé l'accord sujet-verbe (\eng{Income is} vs \eng{Expenses are}), l'usage des modaux (\eng{can't afford}, \eng{should}, \eng{had better}) et la structure \eng{spend/put/save + préposition}. L'oral a révélé des hésitations sur \eng{afford to + V} vs \eng{afford + N}.
    \item \textbf{Sur le plan des compétences :}
    \begin{itemize}
        \item \textit{Writing :} Production d'un texte court, dense, formaté (tableau + commentaires), register semi-formel.
        \item \textit{Speaking :} Prise de parole en continu (pitch) et en interaction (Q\&A), gestion du stress, prononciation du vocabulaire technique.
        \item \textit{Listening/Reading :} Compréhension des profils, des questions adverses, prise de notes pour le vote.
    \end{itemize}
    \item \textbf{Sur le plan culturel/citoyen :} La discussion sur les profils (élève, commerçante informelle, jeune diplômé) a sensibilisé à la diversité des réalités économiques au Cameroun, à l'importance de l'\eng{emergency fund}, de la \eng{tontine} (système d'épargne communautaire) et de la planification (\eng{stick to a budget}) face à l'aléa.
\end{itemize}

\textbf{Pistes de remédiation / Différenciation :}
\begin{itemize}
    \item \textit{Élèves en difficulté :} Fournir un squelette de phrases à trous (\eng{cloze test}) pour le commentaire écrit ; limiter à 5 mots obligatoires ; autoriser des notes pour l'oral.
    \item \textit{Élèves avancés :} Demander d'introduire un imprévu (\eng{unexpected expense} : phone broken, medical issue) et de réajuster le budget en direct ; rédiger un \eng{formal email} de demande de bourse ou de prêt étudiant en réinvestissant le lexique.
    \item \textit{Approfondissement :} Étude de cas réel : « Comment un jeune de Douala a financé sa startup via \eng{crowdfunding} et \eng{microfinance} » (lecture/compréhension pour le devoir).
\end{itemize}

\cle{Bilan} : L'élève est désormais capable de \textbf{parler d'argent} en anglais avec précision, sans calquer le français, et de structurer un raisonnement financier simple — compétence transférable en éducation financière (vie quotidienne) et en anglais des affaires (prochain chapitre).

\newpage

\coursec{Méthodes et exercices résolus}

\courssub{Fiches méthodologiques et exercices résolus}

\begin{methode}[Identifier et employer le vocabulaire du thème]
Pour maîtriser le lexique des finances personnelles (\eng{managing personal finances}), suis cette démarche :
\begin{enumerate}
\item \textbf{Repère le sens général du texte} : parle-t-on d'argent qui entre (\eng{money in} : \cle{earn}, \cle{salary}, \cle{income}) ou d'argent qui sort (\eng{money out} : \cle{spend}, \cle{expenses}, \cle{bills}) ?
\item \textbf{Classe les mots en familles} : verbe (\eng{to save}), nom (\eng{savings}), adjectif (\eng{economical}). Cela t'aide à deviner le sens d'un mot inconnu.
\item \textbf{Apprends les mots avec leurs prépositions} : \eng{to spend money \cle{on} something}, \eng{to borrow money \cle{from} somebody}, \eng{to lend money \cle{to} somebody}, \eng{to pay \cle{for} something}.
\item \textbf{Vérifie le contexte} avant de traduire : \eng{to save} peut signifier « économiser » (argent) ou « sauver » (une vie).
\item \textbf{Constitue une fiche de vocabulaire} en deux colonnes (anglais / français) avec un exemple complet pour chaque mot.
\end{enumerate}
\end{methode}

\begin{propriete}[Faux amis fréquents : argent et finances]
\begin{center}
\begin{tabularx}{\linewidth}{|X|X|X|}
\hline
\rowcolor{popblueL} \textbf{Anglais} & \textbf{Sens correct} & \textbf{Piège (faux ami)} \\
\hline
\eng{to earn} & gagner (un salaire) & \eng{to win} = gagner un match, un prix \\
\hline
\eng{to save} & économiser (de l'argent) & \eng{to save} = sauver (une vie) — sens différent selon contexte \\
\hline
\eng{a debt} & une dette & \eng{debit} = un débit (bancaire) \\
\hline
\eng{a bill} & une facture (UK) / un billet de banque (US) & \eng{a bill} $\neq$ une bille \\
\hline
\eng{to lend} & prêter (donner) & \eng{to borrow} = emprunter (recevoir) \\
\hline
\eng{earnings} & gains, revenus & \eng{winnings} = gains (jeu, pari) \\
\hline
\end{tabularx}
\end{center}
\end{propriete}

\begin{propriete}[Collocations et expressions idiomatiques utiles]
\begin{center}
\begin{tabularx}{\linewidth}{|X|X|X|}
\hline
\rowcolor{popgreenL} \textbf{Expression} & \textbf{Traduction} & \textbf{Exemple} \\
\hline
\eng{to make ends meet} & joindre les deux bouts & \eng{With this salary, it's hard to make ends meet.} \\
\hline
\eng{to save up for} & mettre de côté pour & \eng{She is saving up for a new phone.} \\
\hline
\eng{to live from hand to mouth} & vivre au jour le jour & \eng{Many families live from hand to mouth.} \\
\hline
\eng{a rainy day} & un jour difficile (épargne de précaution) & \eng{Save something for a rainy day.} \\
\hline
\eng{to be broke} & être sans le sou (informel) & \eng{I'm completely broke until payday.} \\
\hline
\eng{to tighten one's belt} & se serrer la ceinture & \eng{We'll have to tighten our belts this month.} \\
\hline
\end{tabularx}
\end{center}
\end{propriete}

\begin{exoresolu}[Exercice 1 — Vocabulary in context]
\textbf{Énoncé.} Read the text and answer the questions in English.

\begin{textebox}[Extrait : Managing a budget in Douala]
“Mireille works as a shop assistant in Douala. She earns 150,000 francs a month. Every month, she pays her rent and her electricity bill, and she spends about 40,000 francs on food and transport. She saves 20,000 francs in a tontine because she wants to buy a sewing machine. Last week, she borrowed 10,000 francs from her brother to pay for her son's school fees.”
\end{textebox}

a) How much does Mireille earn per month?\\
b) What does she spend money on?\\
c) Why does she save money?\\
d) Who did she borrow money from?
\tcblower
\textbf{Solution détaillée.}

\textbf{a)} La question porte sur le revenu (\eng{earn} = gagner de l'argent par son travail). Le texte dit : \eng{She earns 150,000 francs a month.}\\
Réponse rédigée : \eng{Mireille earns 150,000 francs a month.} « Mireille gagne 150 000 francs par mois. »\\
Attention à ne pas écrire \eng{wins} : \eng{to win} signifie « gagner » un match ou un prix, jamais un salaire.

\textbf{b)} On cherche les dépenses (\eng{to spend money \cle{on}} + nom). Le texte mentionne le loyer, la facture d'électricité, la nourriture et le transport.\\
Réponse rédigée : \eng{She spends money on her rent, her electricity bill, food and transport.} « Elle dépense de l'argent pour son loyer, sa facture d'électricité, la nourriture et le transport. »\\
Note la préposition : \eng{to spend \cle{on}}, jamais \eng{to spend for}.

\textbf{c)} On cherche le but de l'épargne (\eng{to save} = mettre de l'argent de côté). Le texte donne la raison avec \eng{because}.\\
Réponse rédigée : \eng{She saves money because she wants to buy a sewing machine.} « Elle économise de l'argent car elle veut acheter une machine à coudre. »

\textbf{d)} Structure : \eng{to borrow \cle{from} somebody} (emprunter à quelqu'un).\\
Réponse rédigée : \eng{She borrowed 10,000 francs from her brother.} « Elle a emprunté 10 000 francs à son frère. »\\
Ne pas confondre avec \eng{to lend \cle{to}} (prêter à) : c'est son frère qui a prêté, c'est Mireille qui a emprunté.

\textbf{Conclusion.} Les quatre réponses sont des phrases complètes au bon temps, avec les bonnes prépositions : c'est ce que l'examinateur attend.
\end{exoresolu}

\begin{methode}[Réemployer le lexique à l'oral et à l'écrit]
Pour réutiliser le vocabulaire financier dans tes propres productions :
\begin{enumerate}
\item \textbf{Prépare des structures toutes faites} : \eng{I earn... / I spend... on... / I am saving up to buy... / I never borrow money because...}.
\item \textbf{Utilise les verbes-clés au bon temps} : habitudes au présent simple (\eng{I save 5,000 francs every month}), actions passées au prétérit (\eng{I lent him money last week}).
\item \textbf{Varie le vocabulaire} : au lieu de répéter \eng{money}, utilise \cle{income}, \cle{salary}, \cle{savings}, \cle{cash}, \cle{pocket money}.
\item \textbf{À l'oral}, entraîne-toi avec des questions types : \eng{How do you manage your pocket money? Do you think it is important to save? Why?}
\item \textbf{À l'écrit}, organise ton paragraphe : idée générale, deux ou trois exemples avec le lexique du thème, conclusion personnelle.
\end{enumerate}
\end{methode}

\begin{exoresolu}[Exercice 2 — Short guided writing]
\textbf{Énoncé.} Write a short paragraph (5 to 7 lines) answering this question: \emph{“How do you manage your pocket money?”} Use at least four of these words: \eng{earn, spend, save, borrow, lend, pocket money, budget}.
\tcblower
\textbf{Solution détaillée.}

\textbf{Étape 1 — Plan.} 1) Combien je reçois ; 2) comment je dépense ; 3) ce que j'économise ; 4) mon avis.

\textbf{Étape 2 — Choix des temps.} Il s'agit d'habitudes : présent simple partout.

\textbf{Étape 3 — Rédaction.}\\
\eng{Every month, my parents give me 10,000 francs as pocket money. I try to manage it carefully. I spend most of it on school supplies, snacks and phone credit. I always save 2,000 francs because I want to buy a bicycle. I never borrow money from my friends, and I do not like lending money either, because it can destroy friendships. In my opinion, making a simple budget is the best way to avoid problems.}

\textbf{Étape 4 — Vérification.}
\begin{itemize}
\item \eng{pocket money} : argent de poche (attention à ne pas dire \eng{silver money} — calque du français « argent »).
\item \eng{spend \cle{on}} : préposition correcte.
\item \eng{lending} : gérondif après \eng{like}, correct.
\item 6 mots de la liste réemployés : \eng{spend, save, borrow, lending, pocket money, budget} — objectif dépassé.
\end{itemize}

\textbf{Conclusion.} Un paragraphe court, au présent simple, avec le lexique exigé et des prépositions correctes : c'est un modèle de production réussie.
\end{exoresolu}

\begin{methode}[Éviter les faux amis et les calques du français]
Le lexique de l'argent est plein de pièges pour les francophones. Réflexes à adopter :
\begin{enumerate}
\item \textbf{Mémorise les paires de faux amis} : \eng{to earn} = gagner (un salaire) $\neq$ \eng{to win} = gagner (un match) ; \eng{to save} = économiser $\neq$ « sauver » ; \eng{a debt} = une dette $\neq$ « un débit » ; \eng{a bill} = une facture (ou un billet en américain) $\neq$ « une bille ».
\item \textbf{Méfie-toi des calques} : « faire de l'argent » se dit \eng{to make money} (correct), mais « gagner sa vie » se dit \eng{to earn a living}, jamais \eng{to win his life}.
\item \textbf{Vérifie la préposition} dans un dictionnaire : \eng{to depend \cle{on}}, \eng{to pay \cle{for}}, \eng{to borrow \cle{from}}, \eng{to lend \cle{to}}.
\item \textbf{Attention à l'ordre des mots} : \eng{I lent him some money} (pas d'inversion, pas de \eng{to} avant \eng{him} avec cette structure : \eng{lend somebody something} ou \eng{lend something to somebody}).
\item \textbf{Relis-toi} en traquant spécifiquement ces pièges avant de rendre ta copie.
\end{enumerate}
\end{methode}

\begin{exoresolu}[Exercice 3 — Correct the mistakes]
\textbf{Énoncé.} Each sentence contains one mistake (false friend, wrong preposition or word order). Correct it.\\
a) \eng{I win 100,000 francs a month in my job.}\\
b) \eng{She spends a lot of money for clothes.}\\
c) \eng{He borrowed 5,000 francs to his neighbour.}\\
d) \eng{My father saves money for to buy a car.}\\
e) \eng{They must pay their debts before December.}
\tcblower
\textbf{Solution détaillée.}

\textbf{a)} Faux ami : on ne « gagne » pas un salaire avec \eng{win}. Correction : \eng{I earn 100,000 francs a month in my job.} « Je gagne 100 000 francs par mois dans mon travail. » (\eng{win} = gagner un match, un prix, une compétition.)

\textbf{b)} Mauvaise préposition : \eng{to spend money \cle{on} something}. Correction : \eng{She spends a lot of money on clothes.} « Elle dépense beaucoup d'argent en vêtements. »

\textbf{c)} Confusion de direction : on emprunte \emph{à} quelqu'un, donc \eng{borrow \cle{from}}. Correction : \eng{He borrowed 5,000 francs from his neighbour.} « Il a emprunté 5 000 francs à son voisin. » (S'il prête, ce serait \eng{He lent 5,000 francs to his neighbour.})

\textbf{d)} Calque du français « pour acheter » : après \eng{for} on met un nom, pas un infinitif. Correction : \eng{My father saves money to buy a car.} ou \eng{My father saves up to buy a car.} « Mon père économise pour acheter une voiture. »

\textbf{e)} Phrase correcte ! \eng{debts} = dettes (orthographe avec un \textbf{b} muet, prononcé « dèts »), \eng{pay} sans préposition devant le nom de chose (dette, facture, salaire). Aucune erreur : c'est le piège de l'exercice.

\textbf{Conclusion.} Toujours vérifier le trio verbe + préposition + direction (qui donne, qui reçoit) : c'est là que les francophones perdent le plus de points.
\end{exoresolu}

\begin{methode}[Exploiter les familles de mots, la prononciation et l'orthographe]
\begin{enumerate}
\item \textbf{Construis des familles} autour d'un mot-clé : \eng{to save} $\rightarrow$ \eng{savings} (économies), \eng{a saver} (épargnant), \eng{safe} (sûr) ; \eng{to lend} $\rightarrow$ \eng{a loan} (un prêt) ; \eng{to pay} $\rightarrow$ \eng{payment}, \eng{paid} (prétérit irrégulier).
\item \textbf{Mémorise les formes irrégulières} : \eng{lend – lent – lent} ; \eng{pay – paid – paid} ; \eng{spend – spent – spent} ; \eng{buy – bought – bought} ; \eng{sell – sold – sold}.
\item \textbf{Attention à l'orthographe piège} : \eng{money} (invariable, jamais \eng{monies} au niveau lycée), \eng{debt} (avec un \textbf{b} muet), \eng{business} (prononcé « biz-nèss »), \eng{receipt} (le \textbf{p} est muet, prononcé « ri-ssit »).
\item \textbf{Note la prononciation en lettres françaises} sur ta fiche : \eng{earn} « eurn », \eng{budget} « beu-djèt », \eng{borrow} « bo-ro ».
\item \textbf{Utilise la famille du mot pour deviner} le sens d'un mot inconnu dans un texte d'examen : si tu connais \eng{to employ}, tu peux comprendre \eng{employer} (employeur), \eng{employee} (employé), \eng{unemployment} (chômage).
\end{enumerate}
\end{methode}

\begin{exoresolu}[Exercice 4 — Word families and verb forms]
\textbf{Énoncé.}\\
a) Complete the table with the missing word from the same family:\\
\eng{to save} (verb) $\rightarrow$ \dots\ (noun) ; \eng{to pay} (verb) $\rightarrow$ \dots\ (noun) ; \eng{to employ} (verb) $\rightarrow$ \dots\ (person who works).\\
b) Put the verbs in brackets in the correct tense:\\
1. Last year, my uncle (lend) \dots\ me 50,000 francs.\\
2. I (already / spend) \dots\ all my pocket money this week.\\
3. Every month, she (pay) \dots\ her bills on time.
\tcblower
\textbf{Solution détaillée.}

\textbf{a)} Familles de mots :
\begin{itemize}
\item \eng{to save} $\rightarrow$ \eng{savings} (nom : les économies ; toujours au pluriel dans ce sens).
\item \eng{to pay} $\rightarrow$ \eng{payment} (nom : le paiement ; suffixe \eng{-ment} fréquent : \eng{pay $\rightarrow$ payment}, \eng{develop $\rightarrow$ development}).
\item \eng{to employ} $\rightarrow$ \eng{employee} (la personne qui travaille ; attention : \eng{employer} = l'employeur, celui qui embauche — les suffixes \eng{-er} et \eng{-ee} ont des sens opposés).
\end{itemize}

\textbf{b)} Conjugaison :

1. \eng{Last year} indique le prétérit simple. \eng{lend} est irrégulier : \eng{lend – lent – lent}. Réponse : \eng{Last year, my uncle lent me 50,000 francs.} « L'an dernier, mon oncle m'a prêté 50 000 francs. » (Jamais \eng{lended}.)

2. \eng{already} + \eng{this week} (période non terminée) imposent le present perfect : \eng{I have already spent all my pocket money this week.} « J'ai déjà dépensé tout mon argent de poche cette semaine. » (\eng{spend – spent – spent}, irrégulier.)

3. \eng{Every month} indique une habitude : présent simple, troisième personne : \eng{Every month, she pays her bills on time.} « Chaque mois, elle paie ses factures à temps. » (Attention au \textbf{s} : \eng{she pays}, et orthographe \eng{pays}, pas \eng{paies}.)

\textbf{Conclusion.} Trois réflexes gagnants : repérer l'indice de temps (\eng{last year}, \eng{already}, \eng{every month}), connaître les verbes irréguliers, vérifier le \textbf{s} de la troisième personne.
\end{exoresolu}

\begin{attention}[Les 5 erreurs qui coûtent des points au Bac]
\begin{enumerate}
\item \textbf{Faux ami \eng{win} / \eng{earn}} : un salaire se \emph{gagne} avec \eng{earn}. Écris \eng{He earns a good salary}, jamais \eng{He wins a good salary}. \eng{Win} est réservé aux matchs, prix et concours.
\item \textbf{Prépositions de l'argent} : \eng{spend \cle{on}}, \eng{pay \cle{for}}, \eng{borrow \cle{from}}, \eng{lend \cle{to}}, \eng{save \cle{for}} (un but) ou \eng{save \cle{to}} + infinitif. Une préposition fausse = un point perdu à chaque occurrence.
\item \textbf{Confusion \eng{borrow} / \eng{lend}} : \eng{borrow} = emprunter (je reçois), \eng{lend} = prêter (je donne). Demande-toi toujours : dans quel sens circule l'argent ?
\item \textbf{Orthographe et prononciation pièges} : \eng{debt} (avec un \textbf{b} muet), \eng{business} (« biz-nèss »), \eng{money} (invariable : \eng{money is}, jamais \eng{money are}), \eng{savings} (avec un \textbf{s} final).
\item \textbf{Calques du français} : « économiser pour acheter » = \eng{save to buy} ou \eng{save up to buy}, jamais \eng{save for to buy} ; « gagner sa vie » = \eng{earn a living} ; « l'argent est important » = \eng{money is important} (singulier, sans article : jamais \eng{the money is important} pour une généralité).
\end{enumerate}
\end{attention}

\newpage

\coursec{Exercices}

\courssub{Je vérifie mes connaissances}

\begin{exercice}[QCM : le bon mot]
Entoure la bonne réponse dans chaque phrase.
\begin{enumerate}
\item My father receives his \textbf{a)} salary \textbf{b)} wage \textbf{c)} tip at the end of every month.
\item Can you \textbf{a)} borrow \textbf{b)} lend \textbf{c)} win me 2,000 francs until Friday?
\item This small car is very \textbf{a)} economic \textbf{b)} economical \textbf{c)} economy: it uses little fuel.
\item The \textbf{a)} price \textbf{b)} cost \textbf{c)} value of living in Douala is higher than in Bafoussam.
\item After paying all his \textbf{a)} bills \textbf{b)} coins \textbf{c)} notes, he had nothing left.
\item She keeps her \textbf{a)} savings \textbf{b)} economies \textbf{c)} savingses in a bank account.
\item I can't \textbf{a)} afford \textbf{b)} offer \textbf{c)} pay a new phone this year.
\item The market woman gave me my \textbf{a)} change \textbf{b)} exchange \textbf{c)} rest after I paid with a 10,000-franc note.
\item He is deeply in \textbf{a)} debt \textbf{b)} debit \textbf{c)} doubt because of his loan.
\item A \textbf{a)} budget \textbf{b)} budgetting \textbf{c)} wallet helps you plan your expenses.
\end{enumerate}
\end{exercice}

\begin{exercice}[Vrai ou faux ? Justifie]
Dis si chaque affirmation est vraie (V) ou fausse (F) et justifie ta réponse en une phrase.
\begin{enumerate}
\item “To lend” and “to borrow” mean exactly the same thing.
\item “Actually” means “actuellement” in English.
\item A person who is “broke” has no money left.
\item “Savings” is always used in the plural.
\item “Wages” are usually paid monthly to office workers.
\end{enumerate}
\end{exercice}

\begin{exercice}[Vocabulaire : définitions]
Trouve le mot de la leçon qui correspond à chaque définition.
\begin{enumerate}
\item Money you receive regularly for your work: \ldots
\item Money you put aside for the future: \ldots
\item A plan of how you will spend your money: \ldots
\item Money you must give back to someone: \ldots
\item Papers you receive asking you to pay for water, electricity or telephone: \ldots
\item To have enough money to buy something: \ldots
\end{enumerate}
\end{exercice}

\begin{exercice}[Conjugaison et formes]
Complète avec la forme correcte du verbe entre parenthèses.
\begin{enumerate}
\item Yesterday, my uncle \ldots (lend) me his bicycle to go to the market.
\item She has never \ldots (borrow) money from anyone.
\item They \ldots (save) 50,000 francs since January.
\item He \ldots (spend) all his pocket money last week.
\item My mother usually \ldots (pay) the electricity bill on time.
\item We \ldots (earn) some money by selling groundnuts during the holidays.
\end{enumerate}
\end{exercice}

\courssub{Je m'entraîne}

\begin{exercice}[Traduction]
Traduis les phrases suivantes en anglais.
\begin{enumerate}
\item Mon frère gagne sa vie en réparant des motos à Bamenda.
\item Je ne peux pas me permettre d'acheter ces chaussures : elles coûtent une fortune.
\item Elle a prêté 10,000 francs à sa voisine la semaine dernière.
\item Nous devons économiser de l'argent pour payer les frais de scolarité.
\item Après avoir payé toutes ses factures, il était complètement fauché.
\end{enumerate}
\end{exercice}

\begin{exercice}[Transformation avec le verbe imposé]
Réécris chaque phrase en utilisant le verbe entre parenthèses, sans changer le sens.
\begin{enumerate}
\item My brother gave me 5,000 francs as a loan. (LEND)
\item She received 3,000 francs from me and must give it back. (BORROW)
\item This bag is too expensive for me. (AFFORD)
\item He puts 2,000 francs aside every week. (SAVE)
\item They gave the taxi driver 500 francs for the ride. (PAY)
\item I have no money left at all. (BROKE)
\end{enumerate}
\end{exercice}

\begin{exercice}[Texte à trous : type Bac]
Complète le texte avec les dix mots de la banque suivante. Chaque mot ne peut être utilisé qu'une seule fois.
\begin{center}
\textbf{salary — bills — savings — afford — borrow — budget — broke — debt — change — expenses}
\end{center}
\begin{textebox}[Managing money in Yaoundé]
Arnaud is a young teacher in Yaoundé. At the end of every month, he receives his \ldots and immediately makes a \ldots to plan his \ldots. First, he pays his \ldots: water, electricity and telephone. Then he puts some money into his \ldots account at the credit union. Last year, he had to \ldots money from a friend because his motorbike broke down, and he is still paying back that \ldots. “When I was a student,” he says, “I was always \ldots by the middle of the month. Now I am careful: I never buy what I cannot \ldots, and I always check my \ldots when I pay in the market.”
\end{textebox}
\end{exercice}

\begin{exercice}[Appariement : expressions idiomatiques]
Associe chaque expression (1--8) à sa signification (a--h).
\begin{enumerate}
\item to be broke
\item to cost a fortune
\item to make ends meet
\item to tighten one's belt
\item to be in the red
\item a nest egg
\item to pay off a debt
\item money doesn't grow on trees
\end{enumerate}
\begin{enumerate}[label=\alph*.]
\item to spend less because money is short
\item to have no money at all
\item to manage to live with the money you earn
\item to be very expensive
\item money saved for the future
\item to finish giving back what you owe
\item to owe more money than you have
\item money is not easy to get, so don't waste it
\end{enumerate}
\end{exercice}

\courssub{Je me prépare au Bac}

\begin{exercice}[Compréhension écrite et langue — type Bac]
Lis le texte suivant, puis réponds aux questions.
\begin{textebox}[A young trader's lesson]
Manka, a 22-year-old girl from Buea, sells second-hand clothes at the market. When she started her business two years ago, she spent everything she earned. “I bought expensive shoes and went out with friends every weekend,” she admits. “By the end of each month, I was completely broke and I even had to borrow money to buy new stock.”

One day, an old trader gave her some advice: “Money doesn't grow on trees, my daughter. Make a budget, pay your debts, and build a nest egg.” Manka listened. She opened a savings account, wrote down all her expenses in a notebook, and tightened her belt. Little by little, she paid off her debts. Today, she can afford a bigger stall, and she has even lent money — with interest — to two other young traders. “Managing money is a skill,” she says with a smile. “Nobody taught me at school. I learned it the hard way.”
\end{textebox}
\textbf{A. Comprehension questions} (réponds par des phrases complètes)
\begin{enumerate}
\item What does Manka do for a living?
\item What was her problem when she started her business?
\item What three pieces of advice did the old trader give her?
\item What did Manka do to change her situation? Mention two actions.
\item What does the expression “I learned it the hard way” mean in the last sentence?
\end{enumerate}
\textbf{B. Language work}
\begin{enumerate}
\item Find in the text: (a) an idiomatic expression meaning “to have no money”; (b) a word meaning “money saved for the future”; (c) a verb meaning “to give money that must be returned”.
\item Put into the past simple: “She opens a savings account and writes down her expenses.”
\item Rewrite using \emph{afford}: “The stall was too expensive for her two years ago.”
\end{enumerate}
\end{exercice}

\begin{exercice}[Traduction — type Bac]
Traduis le passage suivant en anglais.
\begin{textebox}[À traduire]
Beaucoup de jeunes Camerounais ne savent pas gérer leur argent. Dès qu'ils reçoivent leur salaire, ils le dépensent en quelques jours. Pourtant, il est facile de faire un budget : il suffit de noter ses dépenses et de mettre un peu d'argent de côté chaque mois. Ceux qui empruntent de l'argent doivent rembourser leurs dettes à temps. Comme on dit, l'argent ne pousse pas sur les arbres !
\end{textebox}
\end{exercice}

\begin{exercice}[Expression écrite — type Bac]
\textbf{Topic:} Your friend spends all his pocket money in one week and then asks you for money. Write him a short letter (80--100 words) giving him advice on how to manage his money better.
\begin{itemize}
\item Use at least \textbf{six words} from the lesson vocabulary (e.g. \eng{budget, savings, afford, borrow, expenses, broke, bills, earn}).
\item Use at least \textbf{two collocations or idiomatic expressions} (e.g. \eng{to make ends meet, to tighten one's belt, money doesn't grow on trees, to pay off a debt}).
\item Respect the layout of an informal letter: greeting, body, closing.
\end{itemize}
\end{exercice}

\newpage

\coursec{Corrigés des exercices}

\begin{corrige}[Exercice 1 — QCM : le bon mot]
\begin{enumerate}
\item \textbf{a) salary} — \eng{My father receives his salary at the end of every month.} Un \eng{salary} est un salaire fixe versé mensuellement (employés de bureau, fonctionnaires). Un \eng{wage} est payé à l'heure ou à la semaine (ouvriers) ; un \eng{tip} est un pourboire.
\item \textbf{b) lend} — \eng{Can you lend me 2,000 francs until Friday?} Le locuteur demande à recevoir l'argent : l'autre personne \emph{prête}. \eng{Lend something to somebody} / \eng{lend somebody something}.
\item \textbf{b) economical} — \eng{This small car is very economical: it uses little fuel.} \eng{Economical} = « économe, qui consomme peu » ; \eng{economic} = « économique » (relatif à l'économie) ; \eng{economy} est un nom, pas un adjectif.
\item \textbf{b) cost} — \eng{The cost of living in Douala is higher than in Bafoussam.} Collocation figée : \eng{the cost of living} (« le coût de la vie »).
\item \textbf{a) bills} — \eng{After paying all his bills, he had nothing left.} \eng{Bills} = factures (eau, électricité). \eng{Coins} = pièces de monnaie, \eng{notes} = billets de banque.
\item \textbf{a) savings} — \eng{She keeps her savings in a bank account.} \eng{Savings} existe toujours au pluriel ; \eng{economies} est un faux ami (« économies » au sens d'argent mis de côté = \eng{savings}).
\item \textbf{a) afford} — \eng{I can't afford a new phone this year.} \eng{Afford} = avoir les moyens de ; \eng{offer} est un faux ami (« offrir », pas « se permettre »).
\item \textbf{a) change} — \eng{The market woman gave me my change after I paid with a 10,000-franc note.} \eng{Change} = la monnaie rendue ; \eng{exchange} = échange, change (devises).
\item \textbf{a) debt} — \eng{He is deeply in debt because of his loan.} Expression figée : \eng{to be in debt} (« être endetté »). Attention à l'orthographe : le \eng{b} ne se prononce pas (« dèt »).
\item \textbf{a) budget} — \eng{A budget helps you plan your expenses.} Il faut un nom : \eng{budget}. \eng{Budgetting} est en outre mal orthographié (le \eng{t} ne double pas : \eng{budgeting}).
\end{enumerate}
\end{corrige}

\begin{corrige}[Exercice 2 — Vrai ou faux ? Justifie]
\begin{enumerate}
\item \textbf{F.} \eng{To lend} means to give money that must be returned, whereas \eng{to borrow} means to receive it. Ce sont deux actions opposées : \eng{lend to} / \eng{borrow from}.
\item \textbf{F.} \eng{Actually} means “in fact” or “really”; “actuellement” is translated by \eng{currently} or \eng{at present}. Faux ami classique.
\item \textbf{V.} \eng{To be broke} is an informal expression meaning to have no money left. \eng{By the end of the month, I am always broke.}
\item \textbf{V.} \eng{Savings} is always plural: \eng{my savings are in a bank account} (jamais \eng{a saving} pour de l'argent mis de côté).
\item \textbf{F.} \eng{Wages} are usually paid weekly or daily to manual workers; office workers receive a monthly \eng{salary}.
\end{enumerate}
\end{corrige}

\begin{corrige}[Exercice 3 — Vocabulaire : définitions]
\begin{enumerate}
\item \textbf{a salary / an income} — argent reçu régulièrement pour son travail.
\item \textbf{savings} — argent mis de côté pour l'avenir (toujours au pluriel).
\item \textbf{a budget} — un plan de dépenses : \eng{to make / to draw up a budget}.
\item \textbf{a debt} — argent que l'on doit rembourser : \eng{to pay back a debt}.
\item \textbf{bills} — factures : \eng{the electricity bill}, \eng{the water bill}.
\item \textbf{to afford} — avoir les moyens d'acheter : \eng{I can't afford it.}
\end{enumerate}
\end{corrige}

\begin{corrige}[Exercice 4 — Conjugaison et formes]
\begin{enumerate}
\item \eng{Yesterday, my uncle \textbf{lent} me his bicycle to go to the market.} Prétérit irrégulier : \eng{lend -- lent -- lent}. Marqueur : \eng{yesterday}.
\item \eng{She has never \textbf{borrowed} money from anyone.} Present perfect (\eng{has} + participe passé) avec \eng{never} ; \eng{borrow} est régulier.
\item \eng{They \textbf{have saved} 50,000 francs since January.} Present perfect obligatoire avec \eng{since} + point de départ. Ne jamais dire \eng{they save since January}.
\item \eng{He \textbf{spent} all his pocket money last week.} Prétérit irrégulier : \eng{spend -- spent -- spent}. Marqueur : \eng{last week}.
\item \eng{My mother usually \textbf{pays} the electricity bill on time.} Présent simple avec \eng{usually} ; 3\up{e} personne : \eng{pay} + \eng{s} (le \eng{y} reste après voyelle : \eng{pays}, pas \eng{paies}).
\item \eng{We \textbf{earned} some money by selling groundnuts during the holidays.} Prétérit régulier (action passée datée). Attention à l'orthographe : \eng{earn} (« èrne »), pas \eng{win} — on \eng{earns} de l'argent par son travail, on \eng{wins} un match ou à la loterie.
\end{enumerate}
\end{corrige}

\begin{corrige}[Exercice 5 — Traduction]
\begin{enumerate}
\item \eng{My brother earns his living by repairing motorbikes in Bamenda.} Expression : \eng{to earn one's living} (« gagner sa vie ») ; \eng{by} + V-ing pour le moyen.
\item \eng{I can't afford to buy these shoes: they cost a fortune.} \eng{Afford} est suivi de l'infinitif avec \eng{to} ; \eng{to cost a fortune} = « coûter une fortune ».
\item \eng{She lent 10,000 francs to her neighbour last week.} (ou \eng{She lent her neighbour 10,000 francs last week.}) Prétérit de \eng{lend} : \eng{lent}. Ne pas confondre avec \eng{borrowed}.
\item \eng{We must save money to pay the school fees.} \eng{School fees} = « frais de scolarité » (toujours au pluriel) ; \eng{to save money} = « économiser de l'argent ».
\item \eng{After paying all his bills, he was completely broke.} \eng{After} + V-ing ; \eng{broke} = « fauché » (familier).
\end{enumerate}
\textbf{Points de vigilance :} ne pas traduire « gagner sa vie » par \eng{win his life} ; « se permettre » n'est pas \eng{permit oneself} mais \eng{afford} ; « économiser » (argent) = \eng{save}, pas \eng{economize} dans ce contexte courant.
\end{corrige}

\begin{corrige}[Exercice 6 — Transformation avec le verbe imposé]
\begin{enumerate}
\item \eng{My brother \textbf{lent} me 5,000 francs.} (ou \eng{lent 5,000 francs to me}) — prétérit de \eng{lend}.
\item \eng{She \textbf{borrowed} 3,000 francs from me.} — \eng{borrow something from somebody}.
\item \eng{I \textbf{can't afford} this bag.} (ou \eng{I cannot afford to buy this bag.}) — \eng{afford} + nom ou + to-infinitif.
\item \eng{He \textbf{saves} 2,000 francs every week.} — présent simple, habitude (\eng{every week}), 3\up{e} personne en \eng{-s}.
\item \eng{They \textbf{paid} the taxi driver 500 francs for the ride.} — prétérit irrégulier : \eng{pay -- paid -- paid}.
\item \eng{I am completely \textbf{broke}.} — adjectif après \eng{be} ; \eng{broke} est ici un adjectif, pas le prétérit de \eng{break}.
\end{enumerate}
\textbf{Points de vigilance :} respecter les temps de la phrase d'origine ; vérifier la construction de chaque verbe (\eng{lend to} / \eng{borrow from} / \eng{pay somebody} / \eng{afford to do}).
\end{corrige}

\begin{corrige}[Exercice 7 — Texte à trous : type Bac]
Voici le texte complété :
\begin{enumerate}
\item \textbf{salary} — il reçoit son salaire à la fin du mois.
\item \textbf{budget} — \eng{to make a budget} : il établit un budget.
\item \textbf{expenses} — le budget sert à planifier ses dépenses.
\item \textbf{bills} — eau, électricité, téléphone = factures.
\item \textbf{savings} — collocation : \eng{a savings account} (« un compte d'épargne »).
\item \textbf{borrow} — après \eng{had to}, base verbale : il a dû emprunter de l'argent.
\item \textbf{debt} — il rembourse cette dette (\eng{to pay back a debt}).
\item \textbf{broke} — « fauché » au milieu du mois.
\item \textbf{afford} — \eng{cannot afford} : il n'achète pas ce qu'il ne peut pas se permettre.
\item \textbf{change} — il vérifie la monnaie qu'on lui rend au marché.
\end{enumerate}
\textbf{Barème indicatif :} 1 point par mot correctement placé (10 points). Aucune faute de copie tolérée (ex. \eng{buget}, \eng{expences} = 0).
\textbf{Points de vigilance :} chaque mot n'est utilisé qu'une fois ; lire tout le texte avant de répondre, car le contexte (articles, possessifs, verbes) guide le choix ; vérifier la cohérence grammaticale après chaque blanc.
\end{corrige}

\begin{corrige}[Exercice 8 — Appariement : expressions idiomatiques]
\begin{center}
\begin{tabular}{|l|l|}\hline
\rowcolor{popblueL} \textbf{Expression} & \textbf{Signification} \\ \hline
1. to be broke & b) to have no money at all \\ \hline
2. to cost a fortune & d) to be very expensive \\ \hline
3. to make ends meet & c) to manage to live with the money you earn \\ \hline
4. to tighten one's belt & a) to spend less because money is short \\ \hline
5. to be in the red & g) to owe more money than you have \\ \hline
6. a nest egg & e) money saved for the future \\ \hline
7. to pay off a debt & f) to finish giving back what you owe \\ \hline
8. money doesn't grow on trees & h) money is not easy to get, so don't waste it \\ \hline
\end{tabular}
\end{center}
\textbf{Traductions :} \eng{to be broke} « être fauché » ; \eng{to cost a fortune} « coûter une fortune » ; \eng{to make ends meet} « joindre les deux bouts » ; \eng{to tighten one's belt} « se serrer la ceinture » ; \eng{to be in the red} « être dans le rouge / à découvert » ; \eng{a nest egg} « un pécule, des économies » ; \eng{to pay off a debt} « solder une dette » ; \eng{money doesn't grow on trees} « l'argent ne pousse pas sur les arbres ».
\end{corrige}

\begin{corrige}[Exercice 9 — Compréhension écrite et langue — type Bac]
\textbf{A. Comprehension questions} (réponses modèles)
\begin{enumerate}
\item \eng{She sells second-hand clothes at the market in Buea.} (ou \eng{She is a trader: she sells second-hand clothes for a living.})
\item \eng{When she started her business, she spent everything she earned, so she was always broke and had to borrow money to buy new stock.}
\item \eng{The old trader told her to make a budget, to pay her debts, and to build a nest egg.}
\item \eng{She opened a savings account and wrote down all her expenses in a notebook.} (autres réponses possibles : \eng{she tightened her belt} ; \eng{she paid off her debts little by little}.)
\item \eng{It means she learned through difficult personal experience, by making mistakes, not from a teacher.}
\end{enumerate}
\textbf{B. Language work}
\begin{enumerate}
\item (a) \eng{to be broke} ; (b) \eng{a nest egg} (ou \eng{savings}) ; (c) \eng{to lend} (ou \eng{to borrow} selon le point de vue — ici le texte dit \eng{lent money to two other young traders}, donc \eng{lend}).
\item \eng{She \textbf{opened} a savings account and \textbf{wrote} down her expenses.} Prétérits irréguliers : \eng{open -- opened} (régulier), \eng{write -- wrote -- written}.
\item \eng{Two years ago, she \textbf{could not afford} the stall.} (ou \eng{she couldn't afford a bigger stall.}) \eng{Afford} s'emploie presque toujours avec \eng{can / could}.
\end{enumerate}
\textbf{Barème indicatif (sur 20) :} A. 5 questions × 2 pts = 10 pts (1 pt pour le contenu, 1 pt pour la langue) ; B. 1 : 3 pts (1 par élément), 2 : 3 pts (1,5 par verbe), 3 : 4 pts (construction correcte de \eng{afford} avec \eng{could}).
\textbf{Points de vigilance :} répondre par des phrases complètes, en reprenant les termes de la question ; ne pas recopier de longs passages sans les reformuler ; au prétérit, attention aux verbes irréguliers (\eng{wrote}, pas \eng{writed} ; \eng{lent}, pas \eng{lended}).
\end{corrige}

\begin{corrige}[Exercice 10 — Traduction — type Bac]
\textbf{Proposition de traduction :}
\eng{Many young Cameroonians do not know how to manage their money. As soon as they receive their salary, they spend it within a few days. Yet it is easy to make a budget: you just have to write down your expenses and put some money aside every month. Those who borrow money must pay back their debts on time. As the saying goes, money doesn't grow on trees!}

\textbf{Barème indicatif (sur 10) :} 5 segments de sens × 2 pts (1 pt pour le sens, 1 pt pour la langue : grammaire, lexique, orthographe).

\textbf{Points de vigilance :}
\begin{itemize}
\item « gérer leur argent » = \eng{manage their money} (pas \eng{control} ni \eng{direct}).
\item « Dès que » = \eng{as soon as} + présent simple (pas de futur après \eng{as soon as}).
\item « faire un budget » = \eng{to make a budget} (collocation avec \eng{make}, pas \eng{do}).
\item « mettre de côté » = \eng{to put aside} / \eng{to set aside} / \eng{to save}.
\item « rembourser leurs dettes à temps » = \eng{pay back their debts on time} ; \eng{on time} (« à temps ») ≠ \eng{in time} (« à temps pour »).
\item Le proverbe final est une expression figée : \eng{money doesn't grow on trees} — ne pas traduire mot à mot.
\end{itemize}
\end{corrige}

\begin{corrige}[Exercice 11 — Expression écrite — type Bac]
\textbf{Proposition de rédaction modèle (environ 95 mots) :}

\eng{Dear Tabi,}

\eng{I heard you spent all your pocket money in one week and now you are completely broke. Money doesn't grow on trees, my friend! Here is my advice: first, make a budget at the beginning of each month and write down all your expenses. Second, put some money into your savings before you spend anything. Third, never buy what you cannot afford, and avoid borrowing money you cannot pay back. If you follow these steps, you will make ends meet easily.}

\eng{Take care,}

\eng{Mireille}

\textbf{Barème indicatif (sur 20) :}
\begin{itemize}
\item Respect de la tâche et du format de la lettre (salutation, corps, formule finale) : 4 pts.
\item Longueur respectée (80--100 mots) et cohérence du message : 4 pts.
\item Lexique de la leçon : au moins 6 mots du thème correctement employés : 6 pts.
\item Au moins 2 collocations ou expressions idiomatiques : 2 pts.
\item Correction grammaticale et orthographique : 4 pts.
\end{itemize}
\textbf{Points de vigilance :}
\begin{itemize}
\item Mots du lexique présents dans le modèle : \eng{pocket money, broke, budget, expenses, savings, afford, borrowing, pay back, make ends meet} — largement plus de six.
\item Expressions idiomatiques : \eng{money doesn't grow on trees}, \eng{to make ends meet}.
\item Ton amical et conseils formulés avec l'impératif (\eng{make}, \eng{put}, \eng{never buy}) ou \eng{should}.
\item Éviter les calques : « gérer ton argent » = \eng{manage your money} ; « dépenser » = \eng{spend} (pas \eng{expend}) ; « emprunter » = \eng{borrow} (pas \eng{borrow to}).
\item Une lettre informelle se termine par \eng{Take care,} / \eng{Your friend,} / \eng{Best wishes,} suivi du prénom seul.
\end{itemize}
\end{corrige}

\newpage

\coursec{Fiche bilan}

\begin{popfigure}
\centering
\begin{tikzpicture}[mindmap, grow cyclic, every node/.style={concept, font=\small, text=white, execute at begin node=\strut},
    level 1/.append style={level distance=4.2cm, sibling angle=360/7, font=\small},
    level 2/.append style={level distance=3cm, sibling angle=45, font=\scriptsize, text width=2.2cm, align=center}]
\node[concept color=popdark, font=\bfseries\normalsize, align=center]{Managing\\ Personal Finances}
    child[concept color=popblue] { node[align=center]{Money In\\(Income)} 
        child[concept color=popblueL] { node {Earn (v) / Salary (n)} }
        child[concept color=popblueL] { node {Wages (n) / Allowance (n)} }
        child[concept color=popblueL] { node {Revenue (n) / Profit (n)} }
    }
    child[concept color=poporange] { node[align=center]{Money Out\\(Expenses)} 
        child[concept color=poporangeL] { node {Spend (v) / Budget (n)} }
        child[concept color=poporangeL] { node {Bills (n) / Rent (n)} }
        child[concept color=poporangeL] { node {Fees (n) / Utilities (n)} }
    }
    child[concept color=popgreen] { node[align=center]{Saving\ \&\\ Investing} 
        child[concept color=popgreenL] { node {Save (v) / Savings (n)} }
        child[concept color=popgreenL] { node {Account (n) / Interest (n)} }
        child[concept color=popgreenL] { node {Invest (v) / Shares (n)} }
    }
    child[concept color=poppurple] { node[align=center]{Borrowing\ \&\\ Debt} 
        child[concept color=poppurpleL] { node {Borrow (v) / Lend (v)} }
        child[concept color=poppurpleL] { node {Loan (n) / Debt (n)} }
        child[concept color=poppurpleL] { node {Credit (n) / Repay (v)} }
    }
    child[concept color=poppink] { node[align=center]{Banking\\ Services} 
        child[concept color=poppinkL] { node {Deposit (v/n) / Withdraw (v)} }
        child[concept color=poppinkL] { node {Transfer (v/n) / Balance (n)} }
        child[concept color=poppinkL] { node {Statement (n) / ATM (n)} }
    }
    child[concept color=popgold] { node[align=center]{Financial\\ Planning} 
        child[concept color=popgoldL] { node {Budget (n) / Forecast (n)} }
        child[concept color=popgoldL] { node {Goal (n) / Emergency fund (n)} }
    }
    child[concept color=popmuted] { node[align=center]{Idioms\ \&\\ Expressions} 
        child[concept color=popmuted] { node {Make ends meet} }
        child[concept color=popmuted] { node {Nest egg / Rainy day} }
        child[concept color=popmuted] { node {Tighten belt / Break the bank} }
    };
\end{tikzpicture}
\legende{Carte mentale : Vocabulaire de la gestion des finances personnelles (Lesson 22)}
\end{popfigure}

\begin{bilan}
\end{bilan}

\courssub{Lexique clé — Anglais / Français}

\begin{center}
\begin{tabularx}{\linewidth}{|>{\bfseries}X|X|X|}\hline
\rowcolor{popblueL} \textbf{Catégorie} & \textbf{English} & \textbf{Français} \\ \hline
\cellcolor{popblue!10} \textbf{Money In} & \eng{earn, income, salary, wages, allowance, revenue, profit, bonus, pension} & gagner, revenu, salaire, appointement, argent de poche, recette, bénéfice, prime, pension \\ \hline
\cellcolor{poporange!10} \textbf{Money Out} & \eng{spend, expense, expenditure, budget, bill, rent, fee, tuition, grocery, utility, subscription} & dépenser, dépense, budget, facture, loyer, frais, scolarité, courses, charges, abonnement \\ \hline
\cellcolor{popgreen!10} \textbf{Saving \& Investing} & \eng{save, savings, deposit, account, interest, rate, invest, investment, share, stock, bond, dividend, portfolio} & épargner, épargne, dépôt, compte, intérêt, taux, investir, investissement, action, obligation, dividende, portefeuille \\ \hline
\cellcolor{poppurple!10} \textbf{Borrowing \& Debt} & \eng{borrow, lend, loan, debt, credit, creditor, debtor, repay, repayment, instalment, interest, overdraft, mortgage} & emprunter, prêter, prêt, dette, crédit, créancier, débiteur, rembourser, remboursement, échéance, intérêt, découvert, hypothèque \\ \hline
\cellcolor{poppink!10} \textbf{Banking} & \eng{bank, branch, ATM, cash machine, withdrawal, transfer, standing order, direct debit, statement, balance, PIN, contactless} & banque, agence, distributeur, retrait, virement, ordre permanent, prélèvement, relevé, solde, code, sans contact \\ \hline
\cellcolor{popgold!10} \textbf{Planning} & \eng{budget, forecast, goal, target, emergency fund, rainy day, financial literacy, inflation, purchasing power} & budget, prévision, objectif, fonds d'urgence, jour de pluie, éducation financière, inflation, pouvoir d'achat \\ \hline
\end{tabularx}
\end{center}

\courssub{Verbes clés — Formes à connaître par cœur (irréguliers et réguliers)}

\begin{center}
\begin{tabular}{|l|l|l|l|}\hline
\rowcolor{popblueL} \textbf{Infinitif} & \textbf{Prétérit} & \textbf{Participe passé} & \textbf{Sens} \\ \hline
\eng{earn} & \eng{earned} & \eng{earned} & gagner (salaire) — régulier \\ \hline
\eng{spend} & \eng{spent} & \eng{spent} & dépenser (argent, temps) — irrégulier \\ \hline
\eng{save} & \eng{saved} & \eng{saved} & épargner, sauver — régulier \\ \hline
\eng{lend} & \eng{lent} & \eng{lent} & prêter (à qqn) — irrégulier \\ \hline
\eng{borrow} & \eng{borrowed} & \eng{borrowed} & emprunter (à qqn) — régulier \\ \hline
\eng{pay} & \eng{paid} & \eng{paid} & payer — irrégulier \\ \hline
\eng{cost} & \eng{cost} & \eng{cost} & coûter — invariant \\ \hline
\eng{owe} & \eng{owed} & \eng{owed} & devoir (argent) — régulier \\ \hline
\end{tabular}
\end{center}

\courssub{Collocations \& Expressions idiomatiques (à utiliser en speaking/writing)}

\begin{itemize}
\item \eng{make ends meet} — joindre les deux bouts (\eng{My parents work hard to \textbf{make ends meet}.} « Mes parents travaillent dur pour joindre les deux bouts. »)
\item \eng{live from hand to mouth} — vivre au jour le jour (\eng{Many families \textbf{live from hand to mouth}.} « De nombreuses familles vivent au jour le jour. »)
\item \eng{save for a rainy day} / \eng{build a nest egg} — mettre de côté pour les jours difficiles / se constituer un pécule (\eng{She \textbf{saves for a rainy day}.} « Elle met de côté pour les jours difficiles. »)
\item \eng{tighten one's belt} — se serrer la ceinture (réduire ses dépenses) (\eng{We must \textbf{tighten our belts} this month.} « Nous devons nous serrer la ceinture ce mois-ci. »)
\item \eng{break the bank} — coûter très cher / ruiner (\eng{This phone won't \textbf{break the bank}.} « Ce téléphone ne coûtera pas les yeux de la tête. »)
\item \eng{be in the red / in the black} — être dans le rouge / dans le noir (solde négatif / positif) (\eng{My account is \textbf{in the red}.} « Mon compte est dans le rouge. »)
\item \eng{pay through the nose} — payer le prix fort (\eng{We \textbf{paid through the nose} for the tickets.} « Nous avons payé les billets très cher. »)
\item \eng{foot the bill} — payer l'addition / régler la note (\eng{Who will \textbf{foot the bill}?} « Qui paiera l'addition ? »)
\item \eng{get by (on)} — s'en sortir (avec) (\eng{They \textbf{get by on} a small income.} « Ils s'en sortent avec un petit revenu. »)
\item \eng{splash out (on)} — se ruiner pour / s'offrir un luxe (\eng{He \textbf{splashed out on} a new car.} « Il s'est ruiné pour une nouvelle voiture. »)
\end{itemize}

\courssub{Faux amis \& Pièges pour francophones}

\begin{center}
\begin{tabularx}{\linewidth}{|>{\bfseries}X|X|X|X|}\hline
\rowcolor{popred!20} \textbf{Mot anglais} & \textbf{Faux ami (FR)} & \textbf{Vrai sens (EN)} & \textbf{Mot correct (FR)} \\ \hline
\eng{money} & monnaie & argent (billets, pièces, compte) & \eng{currency / change} = monnaie \\ \hline
\eng{coin} & coin (pièce) & pièce de monnaie & \eng{corner} = coin (lieu) \\ \hline
\eng{note} & note (billet) & billet de banque (UK) ; \eng{bill} (US) & \eng{mark / grade} = note (scolaire) \\ \hline
\eng{revenue} & revenu (unique) & recettes (souvent pl., entreprise/État) & \eng{income} = revenu (personnel) \\ \hline
\eng{benefit} & bénéfice & avantage, allocation, prestation & \eng{profit} = bénéfice (commercial) \\ \hline
\eng{charge} & charge (lourde) & frais, coût, accusation & \eng{burden / load} = charge (lourde) \\ \hline
\eng{interest} & intérêt (seul sens) & intérêt (financier) \textbf{et} intérêt (curiosité) & \eng{interest rate} = taux d'intérêt \\ \hline
\eng{loan} & prêt (verbe) & prêt (nom, somme prêtée) & \eng{lend} = prêter (verbe) \\ \hline
\eng{borrow} & prêter & emprunter (prendre) & \eng{lend} = prêter (donner) \\ \hline
\eng{expense} & expensive (adj.) & dépense (souvent pl. \eng{expenses}) & \eng{expensive} = cher (adj.) \\ \hline
\eng{afford} & s'offrir & avoir les moyens de (\eng{can afford}) & \eng{offer} = offrir \\ \hline
\end{tabularx}
\end{center}

\courssub{Familles de mots — Dérivation (Noun / Verb / Adj / Adv)}

\begin{center}
\begin{tabular}{|l|l|l|l|l|}\hline
\rowcolor{popblueL} \textbf{Noun} & \textbf{Verb} & \textbf{Adjective} & \textbf{Adverb} & \textbf{Traduction (N)} \\ \hline
\eng{earner / earnings} & \eng{earn} & \eng{earning} & — & gagneur / gains \\ \hline
\eng{spending / spender} & \eng{spend} & \eng{spending} & — & dépenses / dépensier \\ \hline
\eng{savings / saver} & \eng{save} & \eng{saving} & — & épargne / épargnant \\ \hline
\eng{investment / investor} & \eng{invest} & \eng{invested} & — & investissement / investisseur \\ \hline
\eng{borrower / loan} & \eng{borrow} & \eng{borrowed} & — & emprunteur / prêt \\ \hline
\eng{lender / loan} & \eng{lend} & \eng{lent} & — & prêteur / prêt \\ \hline
\eng{budget} & \eng{budget} & \eng{budgetary} & \eng{budgetarily} & budget \\ \hline
\eng{finance / financier} & \eng{finance} & \eng{financial} & \eng{financially} & finance / financier \\ \hline
\eng{economy / economist} & \eng{economise} & \eng{economic / economical} & \eng{economically} & économie / économiste \\ \hline
\eng{poverty / wealth} & — & \eng{poor / wealthy} & \eng{poorly / wealthily} & pauvreté / richesse \\ \hline
\end{tabular}
\end{center}

\courssub{Prononciation \& Orthographe — Points de vigilance}

\begin{itemize}
\item \eng{earn} [« è-ne »] ≠ \eng{yearn} [« ye-ne »] (désirer ardemment).
\item \eng{debt} [« dette »] — \textbf{b} muet ; \eng{subtle} [« se-tle »] — \textbf{b} muet.
\item \eng{mortgage} [« mor-gage »] — \textbf{t} muet.
\item \eng{receipt} [« ri-sit »] — \textbf{p} muet ; \eng{receive} [« ri-siv »].
\item \eng{budget} [« ba-djè-te »] — \textbf{d} = [dj] devant \textbf{u}.
\item \eng{financial} [« fai-nan-chle »] — accent sur \textbf{nan}, \textbf{ci} = [ch].
\item \eng{economic} [« i-co-no-mic »] vs \eng{economical} [« i-co-no-mi-cal »].
\item \eng{afford} [« a-ford »] — pas de \textbf{double f} à l'oral.
\item \eng{poverty} [« po-ver-ti »] vs \eng{poor} [« pore »].
\item \eng{withdrawal} [« wi-dro-al »] — \textbf{w} = [w], \textbf{th} = [th] (pas [t]).
\item \eng{instalment} [« in-stol-ment »] (UK) / \eng{installment} (US) — attention à l'orthographe.
\end{itemize}

\courssub{Méthodes express}

\begin{enumerate}
\item \textbf{Classer par flux :} « Money In » (Income) vs « Money Out » (Expenses) → \eng{Budget = Income - Expenses}.
\item \textbf{Maîtriser les paires verbales :} \eng{Lend} (give) → \eng{to sb} / \eng{Borrow} (take) → \eng{from sb}. \eng{Can you \textbf{lend} me 1000 francs? = Can I \textbf{borrow} 1000 francs from you?}
\item \textbf{Utiliser les collocations :} Ne pas dire « \eng{make a budget} » mais « \eng{draw up a budget} » ou « \eng{stick to a budget} » ; pas « \eng{put money} » mais « \eng{deposit / save / set aside money} ».
\item \textbf{Rédiger un paragraphe « Budget » :} 1. Income sources. 2. Fixed vs variable expenses. 3. Savings goal. 4. Emergency fund. Connecteurs : \eng{Firstly, Moreover, However, Consequently, Finally}.
\item \textbf{Oral (Speaking) :} Simuler un rendez-vous bancaire (ouvrir un compte, demander un prêt) ou une discussion en famille sur le budget mensuel. Vocabulaire fonctionnel : \eng{I'd like to... / Could you tell me...? / What are the interest rates?}
\end{enumerate}


\begin{aretenir}[Checklist avant le Bac]
$\square$ Je distingue \cle{income} (revenus) et \cle{expenses} (dépenses) et je sais lister 5 exemples de chaque.
$\square$ Je conjugue sans erreur les verbes \eng{earn, spend, save, lend, borrow, pay, cost, owe} (irréguliers : spend/spent/spent, lend/lent/lent, pay/paid/paid, cost/cost/cost).
$\square$ Je ne confonds plus \eng{lend} (prêter) et \eng{borrow} (emprunter) : \eng{lend TO / borrow FROM}.
$\square$ Je connais 5 collocations idiomatiques (\eng{make ends meet, tighten belt, nest egg, rainy day, in the red}).
$\square$ J'évite les faux amis : \eng{money ≠ monnaie, note ≠ note (scolaire), benefit ≠ bénéfice, charge ≠ charge (lourde)}.
$\square$ Je sais former les noms en \eng{-er/-or} (\eng{borrower, lender, investor, saver, spender}) et les adjectifs en \eng{-al/-ial} (\eng{financial, budgetary, economical}).
$\square$ Je prononce correctement \eng{debt, mortgage, receipt, budget, financial, afford, withdrawal, instalment}.
$\square$ Je sais rédiger un budget personnel simple en anglais (structure : Income / Fixed Expenses / Variable Expenses / Savings / Balance).
$\square$ Je peux expliquer la différence entre \eng{saving} (épargne de précaution, faible risque) et \eng{investing} (placement à risque/rendement).
$\square$ Je maîtrise le vocabulaire bancaire de base : \eng{account, deposit, withdrawal, transfer, statement, balance, PIN, overdraft, standing order, direct debit}.
$\square$ Je sais utiliser \eng{afford} correctement : \eng{can/can't afford TO + V} (jamais \eng{afford buying}).
$\square$ Je suis capable de participer à un dialogue « À la banque » ou « Budget familial » en employant le lexique de la leçon.
\end{aretenir}

\findecours

\end{document}
